\begin{abstract}
We develop a general $\infty$-categorical framework for internal Eilenberg-MacLane objects, internal homotopy groups, and an internal notion of covering spaces, and study their behavior under suitable localizations. Applying this to the ordinary motivic localization, we identify internal $n$-Eilenberg-MacLane objects with strongly $\mathbb{A}^1$-invariant sheaves of (abelian when $n\geq 2$) groups, yielding a formal obstruction to the motivic homotopy category being an $\infty$-topos. We prove that taking $\mathbb{A}^1$-localizations induces an equivalence between classical $\mathbb{A}^1$-coverings of a Nisnevich local space and the internal motivic coverings of its motivic localization, providing a streamlined proof of a generalized motivic Van Kampen theorem. Along the way, we establish a Nisnevich-local-to-global $\mathbb{A}^1$-connectivity criterion: a $k$-scheme is $\mathbb{A}^1$-connected if and only if it admits a Nisnevich cover by $\mathbb{A}^1$-connected schemes such that each pairwise intersection has a ($k$-)point. Finally, we show that over a general Qcqs base, passing to the birational motivic homotopy category recovers certain essential topos-theoretic properties absent in the ordinary $\mathbb{A}^1$-setting. In fact, for a scheme with finitely many generic points, we show that the birational motivic homotopy category is a Postnikov-complete $\infty$-topos of cohomological dimension $0$.
\end{abstract}

\tableofcontents

\section{Introduction}

\subsection{Notations and terminologies}
\textit{Categorical.} 
We shall freely use the language of $\infty$-categories without referring to any particular choice of a model for the theory. The reader may freely use their preferred model. A standard reference is [\cite{lurie2009higher}], where the author models these on Kan complexes as models for $\infty$-groupoids.

We choose and fix Grothendieck universes $\mathcal{G}_0\in \mathcal{G}_1\in \mathcal{G}_2$ of small, large, and very large sets, respectively. Unless otherwise specified, by a space we shall always mean a $\mathcal{G}_0$-space. $Spc$ will denote the $\mathcal{G}_1$-$\infty$-category of $\mathcal{G}_0$-spaces (or Kan complexes, if the reader prefers). Unless otherwise specified, by an $\infty$-category we mean a $\mathcal{G}_1$-$\infty$-category. We say that a $\mathcal{G}_1$-$\infty$-category is locally small if it is $\mathcal{G}_0\mbox{-}Spc$-enriched. We shall use the term "very large" $\infty$-categories for locally-$\mathcal{G}_1$-small $\mathcal{G}_2$-$\infty$-categories. 

 By $Pr^L$ we shall mean the $\infty$-category of large presentable $\infty$-categories, while $Cat_{\infty}$ will denote the $\infty$-category of all large $\infty$-categories. $\mathcal{T}op_\infty$ will stand for the $\infty$-category of $\infty$-toposes and geometric morphisms. Thus $Cat_\infty$, $Pr^L$, etc., are examples of very large $\infty$-categories.
 


For an $\infty$-category $\mathcal{C}$, we use $\mathrm{Map}_\mathcal{C}(-,-)$ to denote the $Spc$-valued hom in $\mathcal{C}$. By $\mathcal{P}(\mathcal{C})$ we mean the category of presheaves of spaces on $\mathcal{C}$. If $\mathcal{C}$ has coproducts, by $\mathcal{P}_\Sigma(\mathcal{C})$ we mean the full subcategory of $\mathcal{P}(\mathcal{C})$ consisting of presheaves that send finite coproducts to products.  

\textit{Algebraic geometry.}
Throughout this paper, $S$ will denote a Qcqs scheme, and $k$ will denote a field. We do not assume $k$ is perfect unless otherwise specified. By a morphism of schemes, we shall always mean a quasicompact, separated morphism. In particular, smooth morphisms will be assumed to be of finite presentation.


For a scheme (Qcqs) $S$, we denote by $Sm_S$ the category of smooth schemes over $S$. By $\mathcal{P}(S)$ we shall mean the $\infty$-category of presheaves of spaces on $Sm_S$ and call an object of $\mathcal{P}(S)$ an $S$-space. Similarly, $\mathcal{P}_\Sigma(S):=\mathcal{P}_\Sigma(Sm_S)$. If $\sigma$ is a Grothendieck topology on $Sm_S$, then $\mathcal{P}_\sigma(S)$ stands for the $\infty$-topos of $S$-spaces satisfying descent with respect to $\sigma$-sieves.

 \textit{Motivic homotopy.}
We will use the notation $\mathcal{H}^{mot}(S)$ for the full subcategory of $\mathcal{P}(S)$ consisting of $\mathbb{A}^1$-local, nisnevich-local objects. The corresponding localization functors will be denoted $L_{mot}$. Similarly, $\mathcal{H}^{bir}(S)$ and $L_{bir}$ will handle birational motivic localization, i.e., the localization of $\mathcal{H}^{mot}(S)$ at all dense open immersions. When working with the $\mathbb{A}^1$-motivic homotopy category, most of our important results will be sensitive to the dimension of the base scheme and will work best over fields. Thus, we shall mostly be working with $\mathcal{H}^{mot}(k)$, where $k$ is a field. For birational local results, however, we will work over a general Qcqs base scheme. The point is that, due to the main results of [\cite{0bat}], the birational motivic homotopy category admits a simpler construction.

By $\pi_0$ (similarly $\pi_n$), we shall always mean the presheaf of homotopy groups. This is also the $0$th truncation functor for $\mathcal{P}(S)$. $\pi_n^{nis}$ will denote the homotopy groups internal to the nisnevich topos, i.e., the nisnevich sheafification of $\pi_n$. $\pi_i^{\mathbb{A}^1}$ will, as usual, stand for $\pi_i^{nis}L_{mot}$. Similarly, $\pi_i^b$ is defined as $\pi_i^{nis}L_{bir}$, which is known to be equivalent to $\pi_iL_{bir}$ [\cite{0bat}, Corollary 2.1.8].
We will use italic $X,Y,Z, R,S,T$ for schemes, calligraphic letters $\mathcal{C},\mathcal{D}$, etc., for categories, and $\esc{X}, \esc{Y}$, etc., for $S$-spaces.

\subsection{The philosophy of the current article.}

Let us return to the starting point of $\mathbb{A}^1$ homotopy theory from the universal homotopical viewpoint. The $\mathbb{A}^1$ homotopy category of [\cite{morel19991}] is obtained from the category of smooth schemes as the universal presentable $\infty$-category that `satisfies' Nisnevich descent and `$\mathbb{A}^1$-invariance' [\cite{robalo2012noncommutative}, Theorem 5.2].  

The classical analogue of this concept is the following: start with the category of smooth \textbf{\textit{manifolds}}, then impose open cover descent and $\mathbb{R}^1$-invariance. The beauty of this construction is that it yields an $\infty$-topos. Indeed, since smooth manifolds have good open covers, the resulting category is the $\infty$-category $Spc$ of spaces itself [\cite{dugger1999sheaves}, Remark 3.4.10], which is an $\infty$-topos.  

It is a well-established property that the $\mathbb{A}^1$-motivic homotopy category does not constitute an $\infty$-topos, even over fields, which introduces significant complexities to motivic calculations. This structural constraint was initially noted in [\cite{spitzweck2012motivic}, Remark 3.5]. While the core logic of their remark is sound, a complete formalization necessitates explicitly identifying the appropriate Eilenberg-MacLane objects within the motivic framework, a technical detail that naturally fell outside the immediate scope of their discussion. 

In this work, we will formalize this observation using a categorical framework and then \textit{explicitly identify} the `$n$-Eilenberg-MacLane objects' in the $\mathbb{A}^1$ homotopy category with strongly $\mathbb{A}^1$-invariant sheaves of (abelian when $n\geq 2$) groups. The same arguments as in \textit{loc. cit.} will then show that the motivic homotopy category is not an $\infty$-topos.

The real reason for such a failure seems to be that, without a controlled etale localization, smooth schemes have many more `non-linear' homotopy dimensions. Without imposing such an `extreme' etale localization, the only way to reduce homotopy dimensions is to identify schemes with their generic points. This is exactly what the birational localization promises to achieve (see \S4.0). Throughout this paper, we will see that, over arbitrary base schemes, the birational homotopy category is very close to being an $\infty$-topos, in that it has many properties similar to those of an $\infty$-topos. We will finally show that it is indeed an $\infty$-topos when the base scheme is geometrically unibranch with finitely many generic points (for example, a field) (see \Cref{Hbir a topos}).

\subsection{A Quick summary of the article}
In the second section, we define and study the categorical theory of truncations, homotopy groups, and Eilenberg-MacLane spaces, and discuss their behavior under an accessible (Cartesian) localization. We identify suitable hypotheses on the localization functor to ensure that these notions behave as expected. We recall the definitions here for the convenience of the reader. 

Given an $\infty$-category $\mathcal{C}$, we say that an object $X\in \mathcal{C}$ is $n$-truncated if the presheaf $\mathrm{Map}_\mathcal{C}(-,X):\mathcal{C}^{op}\to Spc$ is (sectionwise) $n$-truncated. The full subcategory of $n$-truncated objects is denoted $\mathcal{C}_{\leq n}$ and is closed under small limits [\cite{lurie2009higher}]. Thus, when $\mathcal{C}$ is presentable, this subcategory is given by a localization $\tau_{\leq n}^{\mathcal{C}}:\mathcal{C}\to \mathcal{C}_{\leq n}$. At the other extreme lies the iterated loop space construction for pointed objects. Given a pointed object $X$, one defines the loop space of $X$, denoted $\Omega_\mathcal{C}X$, as the kernel pair of $*\to X$, which is canonically pointed, so this operation can be iterated. One defines the internal homotopy groups $\pi_n^\mathcal{C}X$ as the $0$th internal truncation of $\Omega^n_\mathcal{C}X$. The orthogonal theory of truncated objects is the theory of internal $n$-connective objects, i.e., pointed objects whose internal homotopy groups vanish below $n$. Finally, one defines an $n$-Eilenberg-MacLane object as an $n$-truncated $n$-connective object. 

To this end, we will discuss the cases of $0$- and $1$-covering spaces under the names geometric and spatial coverings, respectively. When the presentable category is given by a localization of an $\infty$-topos, there are two primary ways to define spatial coverings. One is internal to the localized subcategory, while the other is an $L$-local theory developed inside the base $\infty$-topos. We will find appropriate conditions so that these agree under the localization.

We study the behavior of each of these constructions under a suitable localization. Let us skip the discussion of the general techniques and see their applications to various motivic categories. 

The third section is thus dedicated to the internal homotopy theory of the $\mathbb{A}^1$ motivic homotopy category, $\mathcal{H}^{mot}:=L_{\mathbb{A}^1}\mathcal{P}_{nis}$ [\cite{morel19991}]. Since this is a presentable $\infty$-category, by the general method mentioned above, there is an internal homotopy path component functor providing a let adjoint to the inclusion of $\mathbb{A}^1$-invariant nisnevich sheaves of sets $Sh^{\mathbb{A}^1}$ into $\mathcal{H}^{mot}$ which we denote by $\pi_0^{mot}$. There is a canonical natural transformation $\pi_0^{\mathbb{A}^1}\to \pi_0^{mot}$. Morel's conjecture [\cite{morel2012a1}, Conjecture 1.12], which has been proven to be false by Ayoub [\cite{ayoubcounterexamples}], is then read as saying that this transformation is an isomorphism. In the first \S3.1, we study some properties of this morphism. To this end, we say that a motivic space $\esc{X}$ is motivically connected if $\pi_0^{mot}\esc{X}=*$. The first key result we obtain is:
\begin{recall*}{motivic 0 truncation is injective}
  Over a field $k$, a motivic space is $\mathbb{A}^1$-connected iff it is motivically connected. 
\end{recall*}

 Using this, we will show that:
 \begin{recall*}{motivic em}
Let $k$ be a perfect field. The restriction of the internal motivic homotopy groups functor to motivic Eilenberg-MacLane spaces induces equivalences of categories 
\begin{flalign*}
    Sh_k^{\mathbb{A}^1}&\simeq \mathcal{EM}_0(\mathcal{H}^{mot}):\pi_0^{mot}\\
    Grp_k^{\mathbb{A}^1}&\simeq \mathcal{EM}_1(\mathcal{H}^{mot}): \pi_1^{mot}\\
    Ab_k^{\mathbb{A}^1}&\simeq \mathcal{EM}_n(\mathcal{H}^{mot}):\pi_n^{mot} \text{  for all  }n\geq 2
\end{flalign*}
 \end{recall*}
This helps us present a complete argument for [\cite{spitzweck2012motivic}, Remark 3.5] as a concrete theorem:
\begin{recall*}{mot not topos}
    Over any perfect field $k$, the motivic homotopy category $\mathcal{H}^{mot}(k)$ is not an $\infty$-topos. 
    \end{recall*} 
And that:
     \begin{recall*}{mot bar omega}
Over a perfect field, there is an equivalence of categories $$ \mathrm{B}^{nis}:  \mathcal{G}rp_{mot}(\mathcal{H}^{mot}(k))\simeq\mathcal{H}_\bullet^{mot}(k)_{\geq 1}:\Omega$$  where $\mathcal{G}rp_{mot}(\mathcal{H}^{mot}(S))$ denotes the full subcategory of $ Grp(\mathcal{H}^{mot}(k))$ consisting of motivic monoids with strongly $L_{mot}$-local $\pi_0^{mot}.$ This provides further evidence that $\mathcal{H}^{mot}$ is not an $\infty$-topos.
    \end{recall*}  

In the next subsection, we study the geometric and spatial theories of coverings for the motivic homotopy category. The first significant result implied by \Cref{motivic 0 truncation is injective} is the following counterpart of the classical result stating that ``if $X$ is a topological space with a cover $U_\alpha$ such that each $U_\alpha$ is path connected and $U_{\alpha\beta}\neq \emptyset$, then $X$ is path connected.''
\begin{recall*}{a1 geom conn}
 A scheme $X/k$ is $\mathbb{A}^1$-connected iff there is a nisnevich covering $U_i\to X$ such that each $U_i$ is $\mathbb{A}^1$-connected and for every pair $i,j$ the fiber product $U_i\times_XU_j$ has a ($k$-)point.
\end{recall*}


Next, we apply the $\infty$-categorical notion of covering spaces developed in \S2.4 to the $\mathbb{A}^1$-motivic setup. Using the methods from \S2.4, we reprove various results from [\cite{morel2012a1}, \S7]. As the very first step, we have
\begin{recall*}{Lmot covering is covering of Lmot}
    Suppose $\esc{X}$ is a Nisnevich sheaf of spaces. We then have an equivalence of categories $$\eta^*_{mot}:Cov_{{\mathbb{A}^1}}({L_{mot}}\esc{X)}=Cov_{\mathcal{H}^{mot}}({L_{mot}}\esc{X)}\leftrightarrows  Cov_{{\mathbb{A}^1}}(\esc{X)}: L_{mot}$$ where $\eta_{mot}$ is the localization unit $1\to L_{mot}$. The category on the left-hand side is the category of covering spaces of $L_{mot}\esc{X}$ internal to the motivic homotopy category, while the right-hand side is the $\infty$-categorical analogue of Morel's category of $\mathbb{A}^1$-coverings of $\esc{X}$, as developed in \textup{[\cite{morel2012a1}, \S7]}.
\end{recall*}
Using this, we re-prove Theorem 7.8 of [\cite{morel2012a1}] (see \Cref{a1 simply connected covering space}).

The final section aims to show that much of the extra care required in the $\mathbb{A}^1$-local case is unnecessary in the birational-local case, $\mathcal{H}^{bir}:=L_{bir}\mathcal{H}^{mot}$ [\cite{0bat}]. and in fact, they work over arbitrary Qcqs base schemes:
\begin{recall*}{path injectivity of Lbir}
 Over a general Qcqs scheme $S$, the birational path component functor $\pi_0^{b}:=\pi_0L_{bir}$ is the internal homotopy path component of $\mathcal{H}^{bir}$.
\end{recall*}
Let us quickly list some of the topos-theoretic properties enjoyed by the birational motivic homotopy category: 
\begin{recall*}{bir em}
Similar to that of an $\infty$-topos \textup{[\cite{lurie2009higher}, Lemma 7.2.2.11 (1)]}, then there are equivalences of categories: 
\begin{flalign*}
    \mathrm{Disc}(\mathcal{H}^{bir})_\bullet&\simeq \mathcal{EM}_0(\mathcal{H}^{bir}):\pi_0^{bir}\\
    Grp(\mathrm{Disc}(\mathcal{H}^{bir}))&\simeq \mathcal{EM}_1(\mathcal{H}^{bir}): \pi_1^{bir}\\
    Ab(\mathrm{Disc}(\mathcal{H}^{bir}))&\simeq \mathcal{EM}_n(\mathcal{H}^{bir}):\pi_n^{bir}  \text{ for all }n\geq 2
\end{flalign*}
over an arbitrary Qcqs bse scheme $S$.
\end{recall*}
\begin{recall*}{bir omega bar}
   Over a general Qcqs base scheme $S$, the functor $$\Omega: \mathcal{H}^{bir}_{\geq 1} \to  Grp(\mathcal{H}^{bir})_\bullet\simeq Grp(\mathcal{H}^{bir})$$ is an equivalence of infinity categories, with inverse given by the internal bar construction $\mathbf{B}_{bir}$.
\end{recall*}
However, since $L_{bir}$ is neither locally cartesian (see [\cite{0bat}, Counterexample 3.2.13]) nor effective (see [\cite{0bat}, Counterexample 3.5.7]), the theory of birational covering spaces is not as strong as its motivic counterpart. Specifically, we can construct universal birational covers only for $\mathbb{A}^1$-motivically connected spaces, rather than for all birationally connected spaces (\Cref{universal birational covering}), which follows from:
 \begin{recall*}{Lbir covering is covering of Lbir}
For a motivically connected space $\esc{X}$, we have an equivalence of categories $$\eta^*_{bir}:Cov_{b}({L_{bir}}\esc{X)}=Cov_{\mathcal{H}^{bir}}({L_{bir}}\esc{X)}\leftrightarrows  Cov_{{b}}(\esc{X)}: L_{bir}$$ where $\eta_{bir}$ is the localization unit $1\to L_{bir}$, $Cov_{\mathcal{H}^{bir}}(-)$ is the category of internal coverings in $\mathcal{H}^{bir} $, while $Cov_{{b}}(\esc{X)}$ is the category of birational coverings of $\esc{X}$.
\end{recall*}
In \S4.5, we finally reach the main result of this paper.: 
\begin{recall*}{bir topos on fin gen pt}
    Let $S$ be a Qcqs scheme with finitely many generic points. Then $\mathcal{H}^{bir}(S)$ is an $\infty$-topos.
\end{recall*}
We summarize the remaining properties enjoyed by the birational motivic homotopy category construction as follows:
\begin{recall*}{bir yoneda}
Let $S$ be a qcqs scheme with finitely many generic points. Then the birational motivic space functor$$h^b_S(-):=L_{bir}h_S(-): Sm_S \to \mathcal{H}^{\mathrm{bir}}(S)$$ exhibits the following structural properties:
\begin{enumerate}
    \item It contracts the affine line $\mathbb{A}^1_S$.
    \item It maps Nisnevich and cdh squares to pushout squares.
    \item The target $\mathcal{H}^{\mathrm{bir}}(S)$ is a Postnikov complete $\infty$-topos of cohomological dimension $0$.
\end{enumerate}

\end{recall*}

\section{Internal homotopy theory in an \texorpdfstring{$\infty$}{}-category}
In this section, we introduce the internal notions of truncated objects, homotopy groups, and Eilenberg-MacLane objects. We will study their behavior under suitable localizations. For various aspects of this analysis to work as expected, we will require different properties of the localization functor. (One such important condition is the path injectivity condition [\Cref{path injectivity}], which is key to the most important results of the first three subsections of this section.) 


In the fourth and final subsection, we will formalize the theory of covering objects internal to an $\infty$-category and study this under a localization. (The key condition of this section on localization functors is the locality of the fundamental of local objects and effectivity.)

\subsection{Truncations}
A $(-1)$-connected space $X$ is called $n$-truncated if, for all choices of base points $x\in \mathrm{Map}_{Spc}(*,X)$, the set $\pi_0X$ is a singleton and, for all $i>n$, the groups $\pi_i(X,x)$ are trivial. The following notion from [\cite{lurie2009higher}, Definition 5.5.6.1] offers an immediate $\infty$-categorical generalization of this concept:
\begin{defn}
 Let $\mathcal{C}$ be an $\infty$-category. An object $\esc{F}$ of $\mathcal{C}$ is said to be $n$-truncated if, for all $\esc{E}\in \mathcal{C}$, the $\infty$-groupoid $\mathrm{Map}_\mathcal{C}(\esc{E},\esc{F})$ is $n$-truncated. In other words, the presheaf $\mathrm{Map}_{\mathcal{C}}(-,\esc{F}):\mathcal{C}^{op}\to Spc$ factors through the $\infty$-category $Spc_{\leq n}$ of $n$-truncated spaces. 
\end{defn}


We denote the category of $n$-truncated objects in $\mathcal{C}$ by $\mathcal{C}_{\leq n}$. In particular, $\mathcal{C}_{\leq 0}$ is a discrete category and we denote $\text{Disc}\mathcal{C}:=\text{N}(\mathcal{C}_{\leq 0})$.

If $\mathcal{C}$ is a presentable $\infty$-category, the full subcategory of $n$-truncated objects is closed under limits and is therefore given by a localization $\tau_{\leq n}^\mathcal{C}$ [\cite{lurie2009higher}, Proposition 5.5.6.5]. 

We now analyze the concept of $n$-truncated objects under a localization. 
\begin{lem}\label[lem]{truncation of localization}
Let $L\mathcal{C}\subset \mathcal{C}$ be a localizing subcategory. Then the inclusion $\mathcal{C}_{\leq n}\bigcap L\mathcal{C\subset}(L\mathcal{C})_{\leq n} $ is an equivalence, i.e., we have an equality $(L\mathcal{C})_{\leq n}=\mathcal{C}_{\leq n}\bigcap L\mathcal{C}$.
\end{lem}
\begin{proof}
 The containment is obvious. Indeed, if $\esc{A}$ belongs to $ L\mathcal{C}$ and $ \mathcal{C}_{\leq n}$, then for any $\esc{B}\in  L\mathcal{C}$, one has $ \text{Map}_{ L\mathcal{C}}(\esc{B},\esc{A})= \text{Map}_{ \mathcal{C}}(\esc{B},\esc{A})$. But this last term is $n$-truncated since $\esc{A}\in \mathcal{C}_{\leq n}$. 
 
 To see the reverse inclusion, note that since $(L\mathcal{C})_{\leq n}\subset L\mathcal{C}$, it suffices to show that $(L\mathcal{C})_{\leq n}\subset \mathcal{C}_{\leq n}$. So let $\esc{X}\in L\mathcal{C}_{\leq n}$, and let $\esc{Y}\in \mathcal{C}$ be an arbitrary object. Then we have $\text{Map}_{ \mathcal{C}}(\esc{Y},\esc{X})\simeq \text{Map}_{ L\mathcal{C}}(L\esc{Y},\esc{X})$. But this last term is an $n$-truncated $\infty$-groupoid because $\esc{X}\in (L\mathcal{C})_{\leq n}$. This shows that $\esc{X}\in \mathcal{C}_{\leq n}$ as well. Thus, $(L\mathcal{C})_{\leq n}\subset \mathcal{C}_{\leq n}$. 
\end{proof}
In other words, the following diagram is a pullback square in $Cat_\infty$:
\[
\xymatrix{
(L\mathcal{C})_{\leq n}\ar@{^(->}[r]\ar@{^(->}[d]&
L\mathcal{C}\ar@{^(->}[d]\\
\mathcal{C}_{\leq n}\ar@{^(->}[r]&\mathcal{C}
}
\]
   
\begin{prop}\label[prop]{L tunc is a localization}
    Let $L\mathcal{C}\subset \mathcal{C}$ be an accessible localization of a presentable $\infty$-category $\mathcal{C}$. Then $(L\mathcal{C})_{\leq n}$ is a localization of $\mathcal{C}_{\leq n}$ and of $\mathcal{C}$.
\end{prop}
\begin{proof}
    Since both $\mathcal{C}_{\leq n}$ and $L\mathcal{C}$ are closed under limits in $\mathcal{C}$, we conclude immediately from the lemma above that so is $(L\mathcal{C})_{\leq n}\subset \mathcal{C}_{\leq n}$. 
\end{proof}
Let us denote the corresponding localization functor by $L_{\leq n}:\mathcal{C}_{{\leq n} }\to(L\mathcal{C})_{\leq n}$.

On the other hand, when $L$ is an accessible localization of a presentable $\infty$-category $\mathcal{C}$, so that $L\mathcal{C}$ is also presentable, we know, by [\cite{lurie2009higher}, Proposition 5.5.6.5], that $(L\mathcal{C})_{\leq n}\subset L\mathcal{C}$ is also a localization. We denote the internal truncation functor for the $\infty$-category $L\mathcal{C}$ by $$\tau^{L\mathcal{C}}_{\leq n}:L\mathcal{C}\to (L\mathcal{C})_{\leq n} .$$ It follows that,
\begin{lem}\label[lem]{formulae for tau under localization}
    $\tau^{L\mathcal{C}}_{\leq n} \simeq {L_{\leq n}\tau_{\leq n}}_{|L\mathcal{C}}$.
\end{lem}

\begin{proof}
Suppose $\esc{X}\in L\mathcal{C}$ and $\esc{Y}\in (L\mathcal{C})_{\leq n}=L\mathcal{C}\bigcap \mathcal{C}_{\leq n}$. We get \begin{flalign}
    \text{Map}_{(L\mathcal{C})_{\leq n}}(\tau^{L\mathcal{C}}_{\leq n}\esc{X}, \esc{Y})&=\text{Map}_{L\mathcal{C}}(\esc{X}, \esc{Y})\text{ [by definition of } \tau^{L\mathcal{C}}_{\leq n}\text{]}\\&=\text{Map}_{\mathcal{C}}(\esc{X}, \esc{Y})\\&=\text{Map}_{\mathcal{C}_{\leq n}}(\tau_{\leq n}^{\mathcal{C}}\esc{X}, \esc{Y})\text{ [since } \esc{Y}\in \mathcal{C}_{\leq n}\text{]}
    \\&
    =\text{Map}_{(L\mathcal{C})_{\leq n}}(L_{\leq n}\tau_{\leq n}^{\mathcal{C}}\esc{X}, \esc{Y})     
     \end{flalign}
     The result now follows from the Yoneda lemma. 
 \end{proof}

\textbf{Internal homotopy groups}

Let $\mathcal{C}$ be a category with finite limits. The projection map $\Lambda^2_2\to \Delta^1$ sending $\{0,1\}$ to $0$ and $2$ to $1$ induces a functor $\mathcal{C}^{\Delta^1}\to \mathcal{C}^{\Lambda ^{2}_2}$. Composing this with the canonical functor $\mathcal{C}_\bullet:=\mathcal{C_{*/}}\to \mathcal{C}^{\Delta^1}$, induces another functor $\mathcal{C}_*\to \mathcal{C}^{\Lambda^2_2}$. Finally, taking limits over $\Lambda^2_2$ we arrive at the loop space functor $\Omega :\mathcal{C}_\bullet\to \mathcal{C}_\bullet$ as the obvious factorization of the composition:
$$\mathcal{C}_\bullet\to \mathcal{C}^{\Lambda^2_2}\xrightarrow{lim}\mathcal{C}.$$ By iteration, we arrive at the $n$-fold loop space functor $\Omega ^n: \mathcal{C}_\bullet\to \mathcal{C}_\bullet$. 

Actually, there are multiple ways to think about this. For example, given a pointing $x:*\to X$, we may define $\Omega_{\mathcal{C}}X $ by either of the following two equivalent pullback diagrams:
\begin{figure}
\[
\xymatrix{
\Omega_{\mathcal{C}}X \ar[r] \ar[d] & pt\ar[d]^{x} & & \Omega_{\mathcal{C}}X \ar[r] \ar[d] & X \ar[d]^{\Delta }\\pt
\ar[r]^{x} & X & & pt \ar[r]^{(x,x)} & X \times X
}
\]\caption{loop object}\label{loop object}
\end{figure}
\begin{defn}\label[defn]{LC homotopy group}
 The $0$-th internal homotopy path components object functor of a presentable $\infty$-category $\mathcal{C}$ is defined as $\pi_0^\mathcal{C}:=\tau_{\leq 0}^{\mathcal{C}}$. For $n>0$, the $n$-th internal homotopy group $\pi_n^\mathcal{C}: \mathcal{C}_*\to \text{Disc}(\mathcal{C})_*$ is defined as $\tau_{\leq 0}^{\mathcal{C}}\Omega ^n_{\mathcal{C}}=\pi_0^\mathcal{C}\Omega^n_\mathcal{C}$.
\end{defn}
\begin{defn}\label[defn]{path cart}
    We say that the localization is path cartesian if $L_{\leq 0}$ is cartesian.
\end{defn}
In the rest of this section, we assume that $L$ is path Cartesian, so that $\tau_{\leq 0}^\mathcal{C}$ is also cartesian. By formal arguments, it is clear that $\pi_n^\mathcal{C}X$ is an $n$-group object in $\text{Disc}(\mathcal{C})$. In short, if we let $n\mbox{-}Grp(\mathcal{D})$ stand for $n$-fold group objects of a $1$-category $\mathcal{D}$, i.e., $n\mbox{-}Grp(\mathcal{D})=Grp(Grp(...(Grp(\mathcal{D}))))$, it is clear by the Eckmann-Hilton argument that for every $n\geq 1$ \[n\mbox{-}Grp(\mathcal{D})=
\begin{cases}
 Grp(\mathcal{D}) \text{ for }n=1  \\
  Ab(\mathcal{D}) \text{ for }n\geq 2
 
\end{cases}\]
we have induced functors $\pi_n^\mathcal{C}: \mathcal{C_\bullet}\to n\mbox{-}Grp(\text{Disc}(\mathcal{C}))$.
\begin{lem}\label[lem]{pi0 to Lpi0}
    Let $L\mathcal{C}$ be an accessible localization of a presentable $\infty$-category $\mathcal{C}$. Then $\pi_0^{L\mathcal{C}}\simeq L_{\leq 0}\pi_0^{\mathcal{C}}$. \end{lem}
    
    \begin{proof}
        This follows at once from the identification $\tau_{\leq 0}^{L\mathcal{C}}=\pi_0^{L\mathcal{C}}$ and \Cref{formulae for tau under localization}.         \end{proof}
    In the case of a localization, the previously obtained truncation formula tells us that $$\pi_0^{L\mathcal{C}}:=\tau^{L\mathcal{C}}_{\leq 0} \simeq {L_{\leq 0}\tau^\mathcal{C}_{\leq 0}}_{|L\mathcal{C}}={L_{\leq 0}\pi_{ 0}^\mathcal{C}}_{|L\mathcal{C}}.$$ In fact, 

\begin{cor}\label[cor]{homotopy groups of localization}
    Let $L\mathcal{C}\subset \mathcal{C}$ be a Bousfield localization of a presentable $\infty$-category $\mathcal{C}$. Then $\pi_i^{L\mathcal{C}}: L\mathcal{C}\to \text{Disc}(LC)$ is equivalent to ${L_{\leq 0}\pi_i^{\mathcal{C}}}_{|L\mathcal{C}}$. In fact, this is also equivalent to $ {L_{\leq 0}^{L\mathcal{C}}\pi_i^{\mathcal{C}}}_{|L\mathcal{C}}: L\mathcal{C}\to \text{Disc}(LC)$.
\end{cor}
\begin{proof}
  First of all, since $L\mathcal{C}\subset\mathcal{C}$ is closed under limits and $\Omega$ is constructed from limits, we find that $\Omega_{L\mathcal{C}}^n=(\Omega^n_\mathcal{C})_{|L\mathcal{C}}$.
  
  
  By \Cref{truncation of localization}, for any $\esc{X}\in L\mathcal{C}$, we have $$\pi_i^{L\mathcal{C}}\esc{X}:=\tau^{L\mathcal{C}}_{\leq 0}\Omega^n_{L\mathcal{C}}\esc{X}\simeq L_{\leq 0}\tau^{\mathcal{C}}_{\leq 0}\Omega_\mathcal{C}^n\esc{X}\simeq L_{\leq 0} \pi_i^{\mathcal{C}}\esc{X}.$$
\end{proof}
\begin{rem}
Note that the above corollary does not say that $L$ takes homotopy groups to homotopy groups. Rather, it says that one must choose a localization at the level of discrete (aka $0$-truncated) objects to obtain the correct notion of internal homotopy `groups'.
\end{rem}

\begin{lem}\label[lem]{omega reduce trancation level}
    Let $\mathcal{C}$ be an $\infty$-category with finite limits, and let $X\in \mathcal{C}_\bullet$ be an $n$-truncated object for some $n$. Then $\Omega_{\mathcal{C}}X$ is $(n-1)$-truncated. 
\end{lem}
\begin{proof}
    Recall that an object $X$ in a presentable $\infty$-category $\mathcal{C}$ is $n$-truncated iff the diagonal $X\xrightarrow{\Delta}X\times X$ is $(n-1)$-truncated [\cite{lurie2009higher}, Lemma 5.5.6.15]. From the pullback diagram on the right of \Cref{loop object}, we deduce that $\Omega_{\mathcal{C}}X$ is an $(n-1)$-truncated object. To see this, note that the base change functor $x^*: \mathcal{C}_{/X}\to \mathcal{C}$ is left exact (in fcat, it is a right adjoint). So by [\cite{lurie2009higher}, Proposition 5.5.6.16], $x^*$ preserves $(n-1)$-truncated objects. In particular, because the right vertical map of the right square of \Cref{loop object} is $(n-1)$-truncated, the pullback condition yields that $\Omega_{\mathcal{C}}X\to *$ is $(n-1)$-truncated. But this is equivalent to saying that $\Omega_\mathcal{C}X$ is $(n-1)$-truncated. 
\end{proof}
\begin{rem}
    The slice condition (i.e., $\Omega \tau_{\leq n}\simeq \tau_{\leq n-1}\Omega$) need not hold unless $\mathcal{C}$ is an $\infty$-topos. In other words, the canonical map $\tau_{\leq n-1}\Omega \to \Omega\tau_{\leq n}$ need not be an equivalence in general. 
\end{rem}
\begin{prop}\label[prop]{htpy groups of n truncated objects}
    Let $\mathcal{C}$ be a presentable $\infty$-category and $X\in \mathcal{C}$ be an $n$-truncated object for some $n\geq 0$. Then for every choice of base point $x:*\to X$ and every $i\geq n$, $\pi_i^\mathcal{C}(X,x)=*\in \mathrm{Disc}(\mathcal{C})$.
\end{prop}
\begin{proof}
 Using either of the pullback squares of \Cref{loop object}, observe that when $X$ is pointed, so is $\Omega_{\mathcal{C}}X$. Fixing a base point of $X$, we will thus drop the base point from the picture. Using \Cref{omega reduce trancation level}, it follows by induction that $\Omega^{i}X$ is $(n-i)$-truncated. So when $i>n+1$, we are done (since $(-2)$-truncated objects are contractible). It remains to show that $\pi_{n+1}^\mathcal{C}(X)=*$. Now, $\Omega^{n+1}X$ is $(-1)$-truncated. So it suffices to know that a pointed $(-1)$-truncated object is contractible. This follows from the similar statement in $Spc$, the Yoneda lemma, and the definition of truncated objects.
 \end{proof} 
 \begin{rem}\label[rem]{truncated condition}
Once again, the converse of the proposition need not be true in general. In fact, this need not even be true in arbitrary $\infty$-toposes. A minimal hypothesis under which the converse holds is hypercompleteness. Indeed, in any $\infty$-topos, such a morphism induces an $\infty$-connective morphism $X\to \tau_{\leq n}^\mathcal{C}X$. So the converse is true, for example, in $\mathcal{P}(S)$ [\cite{lurie2009higher}, Example 7.2.1.9], or in $\mathcal{P}_{nis}(S)$ when $S$ is Qcqs of finite Krull dimension (use [\cite{MR4296353}, Theorem 3.18] and [\cite{lurie2009higher}, Corollary 7.2.1.12] along with the fact that the collection of functors $\{p^*:\mathcal{P}_{nis}(Sm_S)\to \mathcal{P}_{nis}(\acute{e}t_X)\}_{p:X\to S\in Sm_S}$ is conservative).
 \end{rem}
We end this subsection with a few more definitions that we will use in the future:
\begin{defn}\label[defn]{pinL}
When $L\mathcal{C}\subset \mathcal{C}$ is an accessible localization of presentable $\infty$-categories, we define $\pi_0^L:\mathcal{C}\to \mathrm{Disc}(\mathcal{C})$ as the composition $\pi_0^{\mathcal{C}}\circ L$.
\end{defn}
\begin{defn}\label[defn]{L connected and LC connected}
Let $X\in \mathcal{C}$ be an object and let $L$ be a localization of $\mathcal{C}$. We say that $X$ is $L$-connected if $LX$ is a connected object of $\mathcal{C}$, i.e., $\pi_0^LX=\pi_0^\mathcal{C}LX\simeq *$. We say that $X$ is $L\mathcal{C}$-connected if $\pi_0^{L\mathcal{C}}LX\simeq *$.
\end{defn}
\begin{lem}
    An $L$-connected object $X\in \mathcal{C}$ is $L\mathcal{C}$-connected.
\end{lem}
\begin{proof}
    This follows at once from \Cref{pi0 to Lpi0}.
\end{proof}
\subsection{Eilenberg Maclane objects}
Throughout this subsection, $\mathcal{C}$ is a presentable $\infty$-category. In a general presentable $\infty$-category, there could possibly be many ways to define the notion of $n$-connective objects. We choose to define it dually to \Cref{htpy groups of n truncated objects}:
\begin{defn}\label[defn]{L n connective}
    We say that a pointed object $x:*\to X$ in a presentable $\infty$-category $\mathcal{C}$ is 
    \begin{enumerate}
    \item ($0$-)connected, or equivalently, $1$-connective, if $\pi_0^\mathcal{C}X=\tau^\mathcal{C}_{\leq 0}X=*$.
        \item globally $(n+1)$-connective (for some $n\geq 0$) (or $n$-connected) if it is connected and the object $\pi_i^\mathcal{C}(X,x)$ is the terminal object for all $0<i\leq n$.
    \end{enumerate}
\end{defn}

\begin{rem}
 By convention, we say that every object is $(-2)$-connected, and that an object is $(-1)$-connected if $\tau_{\leq -1}^\mathcal{C}X=*$. Similarly, one can define the $0$-connectivity condition even when $\esc{X}$ is not given to be pointed. When we say that an object $X$ is $0$-connected, it is inherently implied that $X$ is $(-1)$-truncated. This is because $\tau_{\leq -1}\tau_{\leq 0}\simeq \tau_{\leq -1}$. However, the above definition is not a generalization of [\cite{lurie2009higher}, Definition 6.5.1.10], where the author defines it for an $\infty$-topos. The problem is that the local homotopy group object $\pi_iX$ cannot always be assembled from the global homotopy groups $\pi_i(X,x)$.
\end{rem}
\begin{rem}
    A version of [\cite{lurie2009higher}, Proposition 6.5.1.12], i.e., `an object $X$ is $n$-connective iff $\tau_{\leq n}^\mathcal{C}X\simeq *$', might be false unless $\mathcal{C}$ is an $\infty$-topos.
\end{rem}


We denote the full subcategory of $\mathcal{C}$ consisting of $n$-connective objects by $\mathcal{C}_{\geq n}$.  
\begin{defn}
An object in $\mathcal{C}$ is called $n$-Eilenberg-MacLane if it is $n$-connective and $n$-truncated in $\mathcal{C}$. We denote by $\mathcal{EM}_n(\mathcal{C})$ the full subcategory of $n$-EM objects in $\mathcal{C}$. 
\end{defn}
Let $\tau_{\leq 0}^\mathcal{C}(L\mathcal{C}) $ denote the image of the composite $L\mathcal{C}\subset \mathcal{C}\xrightarrow[]{\tau_{\leq 0}^\mathcal{C}}\text{Disc}\mathcal{C}$.

\begin{lem}\label[lem]{connective em objects of ordinary localization}
 For any localization $L\mathcal{C}\subset \mathcal{C}$, there are containments $$L\mathcal{C}\bigcap \mathcal{C}{_{\geq n} }\subset (L\mathcal{C}){_{\geq n} }\text{ and }L\mathcal{C}\bigcap \mathcal{EM}_n(\mathcal{C})\subset \mathcal{EM}_n(L\mathcal{C}).$$
\end{lem}
\begin{proof}
    Let $\esc{X}\in L\mathcal{C}\bigcap \mathcal{C}{_{\geq n} }$.
    We shall show that $\esc{X}\in(L\mathcal{C}){_{\geq n} }$. That is $\pi_i^{L\mathcal{C}}\esc{X}=0$ for all $0\leq i\leq n-1$. But we know from \Cref{homotopy groups of localization} that $\pi_i^{L\mathcal{C}}\esc{X}=L_{\leq 0}\pi^\mathcal{C}_{i}\esc{X}=0$ (since $\esc{X}\in\mathcal{C}{_{\geq n} }$).

    Using \Cref{truncation of localization} 
    \begin{flalign*}
   L\mathcal{C}\bigcap \mathcal{EM}_n(\mathcal{C})&=L\mathcal{C}\bigcap \mathcal{C}_{\leq n}\bigcap\mathcal{C}{_{\geq n} }\\&=(L\mathcal{C})_{\leq n}\bigcap L\mathcal{C}\bigcap \mathcal{C}{_{\geq n} }\\&
   \subset (L\mathcal{C})_{\leq n}\bigcap( L\mathcal{C})_{\geq n}=\mathcal{EM}_n(L\mathcal{C})      
    \end{flalign*}
\end{proof}
\begin{defn}\label[defn]{path injectivity}
    We say that the localization $L$ of the $\infty$-topos $\mathcal{C}$ is path-injective if the discrete localization functor $\tau_{\leq 0}^{L\mathcal{C}}: \tau_{\leq 0}^{\mathcal{C}}L\mathcal{C}\to \text{Disc}\mathcal{C}$ detects the terminal object.
\end{defn}
In other words: A localization $L$ is path injective if one has the following: an object $X\in \mathcal{C}$ is $L$ connected iff it is $L\mathcal{C}$ connected.
\begin{prop}\label[prop]{em of localization}
    Let $L$ be a path-injective localization of an $\infty$-topos $\mathcal{C}$. Then there are equalities $$ L\mathcal{C}\bigcap \mathcal{C}{_{\geq n} }=(L\mathcal{C}){_{\geq n} }\text{ and }\mathcal{EM}_n(L\mathcal{C})=L\mathcal{C}\bigcap \mathcal{EM}_n(\mathcal{C}).   $$
\end{prop}
\begin{proof}
Using \Cref{truncation of localization}, $$L\mathcal{C}\bigcap \mathcal{EM}_n(\mathcal{C})=L\mathcal{C}\bigcap \mathcal{C}_{\leq n}\bigcap\mathcal{C}{_{\geq n} }=(L\mathcal{C})_{\leq n}\bigcap \mathcal{C}{_{\geq n} }.$$ It suffices to show that $\mathcal{C}{_{\geq n} }\bigcap L\mathcal{C}=(L\mathcal{C}){_{\geq n} }.$

By \Cref{connective em objects of ordinary localization}, it suffices to show that if $\esc{X}\in L\mathcal{C}{_{\geq n} }$, then $\esc{X}\in \mathcal{C}{_{\geq n} }\bigcap L\mathcal{C}$. This is equivalent to showing that $\pi_i^\mathcal{C}\esc{X}=0$ for all $i\leq n-1$. Since $\tau_{\leq 0}^{L\mathcal{C}}$ is injective on the image of the restriction $\tau_{\leq 0}^\mathcal{C}: L\mathcal{C}\to \text{Disc}\mathcal{C}$, it suffices to show that $\tau_{\leq 0}^{L\mathcal{C}}\pi_i^\mathcal{C}\esc{X}=0$. But this is nothing but $$\tau_{\leq 0}^{L\mathcal{C}}\pi_i^{\mathcal{C}}\esc{X}=L_{\leq 0}\tau_{\leq 0}^{\mathcal{C}}\pi_i^{\mathcal{C}}\esc{X}=L_{\leq 0}\pi_i^{\mathcal{C}}\esc{X}=\pi_i^{L\mathcal{C}}\esc{X}$$ [thanks to \Cref{homotopy groups of localization}], which is given to be $0$ $\forall i\leq n-1$ (by the condition $\esc{X}\in (L\mathcal{C}){_{\geq n} }$).
\end{proof}

\begin{rem}\label[rem]{em in infinity topos}
A key feature of an $\infty$-topos (for example, $\text{Spc}$ of spaces, $\mathcal{P}(Sm_S)$, $\mathcal{P}_\Sigma(Sm_S)$, $\mathcal{P}_{nis}(Sm_S)$, etc.) is that in an $\infty$-topos an object is $n$-connective if and only if its $n-1$-truncation is trivial. In other words, $\mathcal{C}_{\geq n}$ coincides with the fiber of $\tau_{\leq n-1}:\mathcal{C}\to \mathcal{C}$ at the final object [\cite{lurie2009higher}, Proposition 6.5.1.12], provided $\mathcal{C}$ is an $\infty$-topos. 
Following this, it is not hard to show that in an $\infty$-topos $\mathcal{C}$, the restriction of the homotopy group functors $${\pi_n^\mathcal{C}}_{|\mathcal{EM}_n}: \mathcal{EM}_n(\mathcal{C}_*)\to n\mbox{-}Grp(\text{Disc}(\mathcal{C}))$$ induces equivalences of categories [HTT, \cite{lurie2009higher}, Proposition 7.2.2.12]. In words, in an $\infty$-topos, Eilenberg-MacLane objects are nothing more than discrete "group" objects.
\end{rem}

\begin{lem}
 Let $\mathcal{C}$ be an $\infty$-topos and $L$ an accessible localization. Assume that $L_{\leq 0}$ is cartesian. Then for every $\esc{X}\in L\mathcal{C}$, $\pi_n^{L\mathcal{C}}$ is an $n$-group object of $\mathrm{Disc}(L\mathcal{C})$.
\end{lem} 
\begin{proof}
Since $\pi_i^{L\mathcal{C}}=L_{\leq 0}\pi_i^\mathcal{C}$ (\Cref{homotopy groups of localization}), this follows from the fact that $\pi_i^\mathcal{C}$ has the same properties when $\mathcal{C}$ is an $\infty$-topos [\cite{lurie2009higher}, the paragraph after Lemma 6.5.1.2].
\end{proof} 
Given a cartesian $\infty$-category $\mathcal{C}$, an $n$-group object is a grouplike $E_n$-monoid object of $\mathcal{C}$. We denote the full subcategory of $\mathcal{C}$ consisting of $n$-group objects by $n\mbox{-}Grp(\mathcal{C})$. When $\mathcal{C}$ is an ordinary $1$-category, $0\mbox{-}Grp(\mathcal{C})\simeq \mathcal{C}_*$, $1\mbox{-}Grp(\mathcal{C})\simeq Grp(\mathcal{C})$, and finally, $n\mbox{-}Grp(\mathcal{C})\simeq Ab(\mathcal{C})$ for every $n\geq 2$.
\begin{defn}\label[defn]{strictly n local}
    Let $G\in n\mbox{-}Grp(\mathrm{Disc}(\mathcal{C}))$. We say that $G$ is $n$-strictly $L$-local if the canonical map $\mathrm{B}^n_\mathcal{C}G\to \mathrm{B}^n_{L\mathcal{C}}G$ is an equivalence. Denote the full subcategory of $n\mbox{-}Grp(\mathrm{Disc}(\mathcal{C}))$ consisting of $n$-strictly $L$-local objects as $n\mbox{-}Grp^L$.
 \end{defn}  

\begin{prop}\label[prop]{grp em}
Let $L$ be a path-injective localization of an $\infty$-topos. Then for every $n\geq 0$, there is an equivalence of categories $\mathrm{B}^n_{L\mathcal{C}}: n\mbox{-}Grp^L\simeq \mathcal{EM}_n(L\mathcal{C}):\pi_n^{L\mathcal{C}}$. In fact, when restricted to $\mathcal{EM}_n(L\mathcal{C})$, there is an equivalence $\pi_n^{L\mathcal{C}}\simeq \pi_n^\mathcal{C}$.
\end{prop}
\begin{proof}
 Since, by definition, $${\mathrm{B}^n_{L\mathcal{C}}}_{|n\mbox{-}Grp_L(L\mathcal{C})}\simeq\mathrm{B}^n_\mathcal{C}, $$ it follows that the functor $$\mathrm{B}^n_{L\mathcal{C}}: n\mbox{-}Grp_L(L\mathcal{C})\to\mathcal{EM}_n(L\mathcal{C})$$ is well-defined. Now the equivalence  $$\mathrm{B}^n_{\mathcal{C}}: n\mbox{-}Grp(\mathcal{C})\simeq \mathcal{EM}_n(\mathcal{C})$$ [\cite{lurie2009higher}, Proposition 7.2.2.12.] guarantees that the above functor is fully faithful. The equivalence, in fact, guarantees that this is an equivalence, by the definition of $n\mbox{-}Grp^L$. (Observe that the inverse is given by ${\pi_n^\mathcal{C}}_{|\mathcal{EM}_n(L\mathcal{C}}$, i.e., $\pi_n^\mathcal{C}$ has $L$-local $n$-groups as images. Therefore, the inverse is identical to $\pi_n^{L\mathcal{C}}$.) 
\end{proof}

\begin{cor}\label[cor]{not a topos localization}
    Suppose $L$ is a path-injective (and path-cartesian) localization of an $\infty$-topos such that the inclusion $Grp_L(L\mathcal{C})\subset Grp(\mathrm{Disc}(L\mathcal{C}))$ is strict. Then $L\mathcal{C}$ is not an $\infty$-topos.
\end{cor}\begin{proof}
    If $L\mathcal{C}$ is an $\infty$-topos, then by [\cite{lurie2009higher}, Proposition 7.2.2.12], $\pi_1^{L\mathcal{C}}:\mathcal{EM}_n(L\mathcal{C})\to Grp(\mathrm{Disc}(L\mathcal{C}))$ should be an equivalence. But by the above Proposition, this should imply $Grp_L(L\mathcal{C})=Grp(\mathrm{Disc}(L\mathcal{C}))$.
\end{proof}

\subsection{Truncation and group objects}

 Let us briefly recall the bar construction more generally for monoids, for the convenience of the reader. Roughly, a monoid object in a presentable $\infty$-category $\mathcal{C}$ is a Segal object in the monoidal category $\mathcal{C}^\times$, i.e., a simplicial object $M_\bullet$ equipped with a compatible family of equivalences $M_n\simeq M_1\times_{M_0}\cdots \times _{M_0}M_1$. The bar construction $\mathrm{B}M_\bullet$ is defined as the colimit of the simplicial object $M_\bullet$ regarded as a simplicial diagram in $\mathcal{C}$. Because monoids are canonically pointed by the unit morphism, the resulting colimit is also pointed. Thus, the bar construction is given by a functor $\mathrm{B}_\mathcal{C}:\mathrm{Mon}(\mathcal{C})\to \mathcal{C}_\bullet$. 
 
 
 Now, let us return to the internal loop space construction $\Omega_\mathcal{C}:\mathcal{C}_\bullet \to \mathcal{C}$. Given a pointed object $*\xrightarrow{x} X$ in $\mathcal{C}$, we define the loop object of $x$ as the following pullback diagram:
 \[
 \xymatrix{
 \Omega_{\mathcal{C}}X\ar[r]^{p_1}\ar[d]_{p_2}&pt\ar[d]\\
 pt \ar[r]&X
 }
 \]
 In the category of spaces, it is well known that the loop space of any pointed space has a monoid structure. In general, the loop construction in an arbitrary presentable $\infty$-category can be promoted to a functor $\Omega : \mathcal{C}_\bullet \to \mathrm{Mon}(\mathcal{C})$. Indeed, given a pointed space $x: *\to \esc{X}$, one can define the monoid structure on $\Omega \esc{X}$ by identifying it with the $1$-simplex of the Čech nerve $\check{C}x$. But $\check{C}x$ is actually an augmented object, i.e., it belongs to $\mathcal{P}(\Delta_+,\mathcal{C})$ (we shall denote this category by $Aug(\mathcal{C})$). To identify the monoid structure of $\Omega_{\mathcal{C}}X$, we take the unaugmented part of $\check{C}x$, i.e., $(\Omega_{\mathcal{C}}X)_\bullet =(\check{C}x)_{\geq 0}$.
\begin{prop}\label[prop]{bar loop adjunction}
     Let $\mathcal{C}$ be a presentable $\infty$-category. Then there exists an adjunction:
     $$\mathrm{B}_{\mathcal{C}}: \mathrm{Mon}(\mathcal{C}) \leftrightarrows \mathcal{C}_\bullet:\Omega$$
\end{prop}

\begin{proof}
For a monoid object $M_\bullet $ and a pointed object $x:*\to X$, we have,
   \begin{flalign*}
       \mathrm{Map}_{\mathrm{Mon}(\mathcal{C})}(M_\bullet, \Omega_{\mathcal{C}}X)&\simeq\mathrm{Map}_{\mathrm{Mon}(\mathcal{C})}(M_\bullet, \check{C}x_{\geq 0 })&\\
       &\simeq \mathrm{Map}_{Aug(\mathcal{C})}(M_\bullet\to |M_\bullet|, \check{C}x)&\\& \simeq \mathrm{Map}_{\mathcal{P}(\Delta^{\leq 0}_{+}, \mathcal{C})}(*\xrightarrow{1}|M_\bullet| , *\xrightarrow{x} X)\text{ [\cite{lurie2009higher}, Proposition 6.1.2.11]}&\\&
 = \mathrm{Map}_{\mathcal{C}_{\bullet}}(|M_\bullet| , X)&\end{flalign*}
\end{proof}
It is worth recalling that in the category of spaces, the loop space construction yields a grouplike space. One way to see this is to note that the path component of the loop space of a space is a group (namely the fundamental group at the given base point), with the inverse of a based loop $\gamma$ given by $t\mapsto \gamma(-t)$.

In general, the loop construction in an arbitrary presentable $\infty$-category is also a grouplike monoid in $\mathcal{C}$. To see this, it suffices to note that the switching of $p_1$ and $p_2$ (defining $\Omega_\mathcal{C}X$) induces an inverse operation $\iota : \Omega_{\mathcal{C}}X\to \Omega_{\mathcal{C}}X$ for the aforementioned monoid structure of $\Omega_{\mathcal{C}}X$. Therefore, the adjunction of \Cref{bar loop adjunction} restricts to an adjunction:
$$\mathrm{B}_{\mathcal{C}}: \mathrm{Grp}(\mathcal{C}) \leftrightarrows \mathcal{C}_\bullet:\Omega$$
\begin{rem}\label[rem]{em equiv of a topos}
If $\mathscr{C}$ is an $\infty$-topos, then the adjunction $\mathrm{B}_\mathcal{C }\dashv\Omega$ is an equivalence of $\infty$-categories [\cite{lurie2009higher}, Lemma 7.2.2.11 (1)]. 
\end{rem}
When we are in the situation of an accessible localization $L\mathcal{C}\subset \mathcal{C}$, the equivalence from the remark above does not immediately restrict to a similar equivalence for the presentable $\infty$-category $L\mathcal{C}$. Of course, the first issue is that $(L\mathcal{C})_{\geq 1}$ is not always identical to $\mathcal{C}_{\geq 1}\bigcap L\mathcal{C}$. Under the hypothesis of \Cref{em of localization}, we can get rid of this issue. However, there is another issue, which we have already introduced in the previous section, namely, that $\mathrm{B}_{\mathcal{C}}$ need not take local objects to local objects, and therefore $\mathrm{B}_{L\mathcal{C}}\neq {\mathrm{B}_\mathcal{C}}_{|L\mathcal{C}}$.  
\begin{lem}\label[lem]{bar of cart localization}
    Suppose $L$ is a cartesian localization of an $\infty$-topos. Then $\mathrm{B}_{L\mathcal{C}}L\simeq L {\mathrm{B}_\mathcal{C}}$.
\end{lem}
\begin{proof}
    This is obvious because the bar construction is given by finite products and geometric realizations, and $L:\mathcal{C}\to L\mathcal{C}$ commutes with both. 
\end{proof}
\begin{rem}
Clearly, $\mathrm{B}_{L\mathcal{C}}\simeq {L\mathrm{B}_\mathcal{C}}_{|L\mathcal{C}}$. So the adjunction of \Cref{bar loop adjunction} can be deduced for $L\mathcal{C}$ using the adjunction for $\mathcal{C}$. In particular, since every presentable $\infty$-category arises as a localization of a presheaf $\infty$-topos, the adjunction of \Cref{bar loop adjunction} can be deduced even more easily, since the adjunction is well known for presheaf $\infty$-toposes. 
\end{rem}
Let $ Grp_{L}L\mathcal{C}$ denote the full subcategory of $ GrpL\mathcal{C}$ consisting of objects $\esc{X}$ such that $\mathrm{B}_\mathcal{C}\esc{X}\in L\mathcal{C}$, i.e., $\mathrm{B}_{L\mathcal{C}}X\simeq \mathrm{B}_{\mathcal{C}}X$. Following [\cite{elmanto2021motivic}, \S3.4], we will call such objects strongly $L$-local monoids.
\begin{lem}
    For a localization $L\mathcal{C}\subset \mathcal{C}$, there is a canonical equivalence $\Omega_{L\mathcal{C}}\simeq {\Omega_\mathcal{C}}_{|L\mathcal{C}}$.
\end{lem} 
\begin{proof}
    This is clear because $L\mathcal{C}\subset \mathcal{C}$ is closed under limits and, in particular, under fiber products.
\end{proof}
\begin{thm}\label{loop bar for localization}
   Suppose $L$ is a path-injective localization of an $\infty$-topos $\mathcal{C}$. Then the adjunction $$ \mathrm{B}_{L\mathcal{C}} : Grp(L\mathscr{C})\leftrightarrows L\mathcal{C}_\bullet: \Omega_{L\mathcal{C}}\simeq \Omega_\mathcal{C}$$ restricts to an equivalence $$\mathrm{B}_{L\mathcal{C}} : Grp_L(L\mathscr{C})\simeq (L\mathcal{C})_{\geq 1}: \Omega_{L\mathcal{C}}\simeq \Omega_\mathcal{C}.$$ 
 \end{thm}
\begin{proof}
     Under the hypothesis, we know that $(L\mathscr{C})_{\geq 1}= L\mathcal{C}\bigcap \mathscr{C}_{\geq 1}$ (\Cref{em of localization}). So it suffices to show that if $\esc{X}\in L\mathcal{C}\bigcap \mathscr{C}_{\geq 1}$, then $\Omega^{\mathcal{C}}\esc{X}  \in  Grp_L(L\mathscr{C})$. But we know that $\mathrm{B}_\mathcal{C}\Omega^\mathcal{C}\esc{X}\simeq \esc{X}\in L\mathcal{C}$ (the first equivalence holds because $\esc{X}\in \mathcal{C}_{\geq 1}$, [\cite{lurie2017higher}, Theorem 5.2.6.10]), and we are done.
\end{proof}

 
\subsection{Covering theory and localizations}
\textbf{Geometric coverings: (a version of $(-1)$-truncated morphisms)}

We first standardize the notion of an effective epimorphism from [\cite{lurie2009higher}, the paragraph following Corollary 6.2.3.5]. Recall that this notion is defined in \textit{loc.cit.} only for semi toposes. In this text, we adopt the following definition for arbitrary $\infty$-categories:
\begin{defn}[[\cite{nlab:effective_epimorphism_in_ainfinity1category}, effective epimorphism\text{]}]
We say that a morphism in an $\infty$-category is an effective epimorphism if its \v Cech nerve is a colimit diagram.    
\end{defn}
\begin{rem}
    This, however, should be contrasted with the notion of an $(-1)$-truncated morphism in $\mathcal{C}$. Unless $\mathcal{C}$ is a semi-topos, there is no good reason for these two notions to coincide, as in 
    [\cite{lurie2009higher}, Corollary 6.2.3.5].
\end{rem}


\begin{defn}[Geometric covering] 
Let $\mathcal{C}$ be an $\infty$-category and $L\mathcal{C}$ a localization. We say that a collection of morphisms $\{p_i:U_i\to X\}_{i\in I}$ in $\mathcal{C}$ is:
    \begin{enumerate}
        \item A geometric $\mathcal{C}$ cover of $X$ if they are jointly effective epimorphic, i.e., $p:\bigsqcup U_i\to X$ is an effective epimorphism, or equivalently, if the \v Cech nerve of $p$ is a colimit diagram in $\mathcal{C}$.
        \item A geometric $L$-cover of $X$ if $L\check{C}_\bullet p$ is a colimit diagram in $L\mathcal{C}$.
    \end{enumerate}
\end{defn} 
\begin{lem}
    A geometric $\mathcal{C}$-cover is a geometric $L$-cover.
\end{lem}
\begin{proof}
This is obvious since $L$ preserves colimits. 
\end{proof}
\begin{lem}
    An $L$ geometric covering of $L$-local objects is an $L\mathcal{C}$-geometric covering. 
\end{lem}
\begin{proof}
    This is clear since $L\mathcal{C}$ is closed under fiber product in $\mathcal{C}$.
\end{proof}
\begin{lem}[local to global $L$-connectivity]\label{L geom L connectivity}
Suppose $L$ is a path-injective localization of an $\infty$-topos $\mathcal{C}$. Let $\{p_i:U_i\to X\}_{i\in I}$ be a geometric $L$-covering such that each $U_i$ is $L$-connected and each $U_i\times_XU_j$ is a globally nontrivial $L$-local path component (i.e., $\mathrm{Disc}(L\mathcal{C})(*,\pi_0^L(U_i\times_XU_j))\neq \emptyset$; (for example, when these fiber products are pointed). Then $X$ is $L$-connected.
\end{lem}
\begin{proof}
 First, note that the functor $\tau_{\leq 0}^{L\mathcal{C}}: L\mathcal{C}\to \mathrm{Disc}(L\mathcal{C})$, being a localization, preserves colimits. So by applying $\tau_{\leq 0}^{L\mathcal{C}}$ to the colimit diagram $L\check{C}_\bullet p$ in $L\mathcal{C}$ (by definition of geometric $L$-covering), we see that $\tau_{\leq 0}^{L\mathcal{C}}LX$ is the colimit of the path components of $L\check{C}_\bullet p$ in $L\mathcal{C}$, computed in $\mathrm{Disc}(L\mathcal{C})$. As $\mathrm{Disc}(L\mathcal{C})$ is an ordinary category, this colimit can actually be computed as the colimit of the 1-skeleton of the simplicial diagram. In other words, $\tau_{\leq 0}^{L\mathcal{C}}LX$ is the coequalizer of the following diagram in $\mathrm{Disc}(L\mathcal{C})$:
$$\underset{i,j}{\bigsqcup} \pi_0^{L\mathcal{C}}({U_i\times_XU_j})\rightrightarrows \underset{i}{\bigsqcup}\pi_0^{L\mathcal{C}}({U_i})$$
But by formal arguments, this colimit is equal to the colimit of the pairwise pushouts  $$*\simeq\pi_0^{L\mathcal{C}}U_i\leftarrow \pi_0^{L\mathcal{C}}U_i\times_XU_j\to \pi_0^{L\mathcal{C}}U_j\simeq *$$ in $L\mathcal{C}$, say denoted $E_{ij}$. Since every leg $\pi_0^{L\mathcal{C}}U_i\times_XU_j\to \pi_0^{L\mathcal{C}U_i}\simeq *$ is identical and is an epimorphism (by the given condition of the existence of a global section), it follows that the pushout maps $\pi_0^{L\mathcal{C}}U_j\to E_{ij}$ are all isomorphisms. So $E_{ij}\simeq *$. Thus, we have shown that the colimit of each of the pairwise diagrams in $\mathrm{Disc}(L\mathcal{C})$:
$$*\simeq\pi_0^{L\mathcal{C}}U_i\leftarrow \pi_0^{L\mathcal{C}}U_i\times_XU_j\to \pi_0^{L\mathcal{C}}U_j\simeq *$$
yields the terminal object. Hence, the final colimit in $\mathrm{Disc}(L\mathcal{C})$, which is $\pi_0^{L\mathcal{C}}LX$, is equivalent to the colimit of the diagram $*\simeq \pi_0^{L\mathcal{C}}LU_i\to *$ in $\mathrm{Disc}(L\mathcal{C})$ and is thus equivalent to $*$ itself. Since $L$ is path injective, it follows that $\pi_0^LX\simeq*$.
\end{proof}
\begin{cor}\label[cor]{geom L connectivity}
 Let $\{p_i:U_i\to X\}_{i\in I}$ be a geometric $\mathcal{C}$ covering such that each $U_i$ is $L$-locally $0$-connected and each $U_i\times_XU_j$ is globally $L$-connected (e.g., pointed). Then $X$ is $L$-locally $0$-connected.   
\end{cor}

\begin{thm}[$L$-local Van Kampen theorem]\label{L van campen}
Let $L$ be a path-injective localization of an $\infty$-topos $\mathcal{C}$. Let $\{p_i:U_i\to X\}_{i\in I}$ be a well-pointed $\mathcal{C}$ (or, more generally $L$) geometric covering in $\mathcal{C}$. Assume that $U_i\times_XU_j\times_XU_k$ is $L$-connected for every $i,j,k \in \{I\}\cup I$ (where we denote $U_I:=X$) (i.e., 1,2-, 2-, and 3-fold intersections are connected). Then the $L$-local Van Kampen theorem states that there is a coequalizer diagram in $Grp^L$:
$$\underset{i,j}{*}\pi_1^{L\mathcal{C}}({U_i\times_XU_j})\rightrightarrows \underset{i}{*}\pi_1^{L\mathcal{C}}({U_i})\to \pi_1^{L\mathcal{C}}(X)\to *$$
\end{thm}
\begin{proof}
Let us construct a diagram $I^{\Delta^{inj}}$ as follows. Its objects are pairs $(n,i_1\cdots i_n)$ with $n\geq 0$ and $i_1\cdots i_n\in I^n$, and its morphisms are projection maps. Consider the functor $C:(I^{\Delta^{inj}})^{{op}}\to \mathcal{C}$ given by $$(n,i_1\cdots i_n)\mapsto U_{i_1}\times_X\cdots\times _XU_{i_n}.$$ By the effectivity condition, we have a colimit diagram in $\mathcal{C}$: $$\colim_{x\in I^{\Delta^{inj},op}}\Pi x\simeq X.$$ Applying the left adjoint, we obtain a colimit diagram in $L\mathcal{C}$: $$\colim_{x\in I^{\Delta^{inj},op}}L\Pi x\simeq LX.$$ Applying the truncation $\tau_{\leq  1}^{L\mathcal{C}}$ (which is a left adjoint), we find the colimit, $$\colim_{x\in I^{\Delta^{inj},op}}\tau_{\leq 1}^{L\mathcal{C}}L\Pi x\simeq \tau_{\leq 1}^{L\mathcal{C}}LX$$ in $(L\mathcal{C})_{\leq 1}$. Moreover, because $I^{\Delta^{inj},op}$ is a semisimplicial diagram and $(L\mathcal{C})_{\leq 1}$ is a $(2,1)$-category, the colimit can be computed using its $2$-skeleton. 
Since each object in this diagram is connected in $L\mathcal{C}$ and so is the colimit $LX$ (by \Cref{L geom L connectivity}), it follows that  $$\colim_{x\in ({I^{\Delta^{inj}}_{\leq 2}})^{op}}\tau_{\leq 1}^{L\mathcal{C}}L\Pi x\simeq \tau_{\leq 1}^{L\mathcal{C}}LX$$ in $\mathcal{EM}_1(L\mathcal{C})$. For $x\in I^{\Delta^{inj},op}_{\leq 2}$, $L\Pi x$ is given to be $0$-connected in $L\mathcal{C}$, hence so is $\tau_{\leq 1}^{L\mathcal{C}}L\Pi x$. Thus, $$\tau_{\leq 1}^{L\mathcal{C}}L\Pi x\in \mathcal{EM}_1(L\mathcal{C}).$$ Since the colimit in $(L\mathcal{C})_{\leq 1}$ is connected, it is computed in $\mathcal{EM}_1(L\mathcal{C})$. By \Cref{grp em}, the functor $$\pi_1^{L\mathcal{C}}:\mathcal{EM}_1(L\mathcal{C})\to Grp^L$$ preserves colimits. Therefore, we have a diagram in $Grp^L$ (for $(i,j,k)\in I^3$:
$$\pi_1^{L\mathcal{C}}({U_i\times_XU_j}\times U_k)\to \pi_1^{L\mathcal{C}}({U_i\times_XU_j})\to\pi_1^{L\mathcal{C}}({U_i})$$ whose colimit is $\pi_1^{L\mathcal{C}}X$. Since $Grp^L$ is a discrete category, this colimit can also be computed as the colimit of the subdiagram (for $(i,j)\in I^2$)$$\pi_1^{L\mathcal{C}}({U_j})\leftarrow\pi_1^{L\mathcal{C}}({U_i\times_XU_j})\to\pi_1^{L\mathcal{C}}({U_i}).$$ The colimit of this last diagram can be obtained from the diagram in question, namely the coequalizer in $Grp^L$ of the parallel arrows: 
$$\underset{i,j}{*}\pi_1^{L\mathcal{C}}({U_i\times_XU_j})\rightrightarrows \underset{i}{*}\pi_1^{L\mathcal{C}}({U_i})$$
\end{proof}

\textbf{Spatial Covering ($0$-truncated morphisms):}

\begin{defn}\label[defn]{defn for covering}
    Let $\mathcal{C}$ be a presentable $\infty$-category. A morphism $f:Y\to X$ in $\mathcal{C}$ is said to be a $\mathcal{C}$-covering if $f\in ({\mathcal{C}_{/X}})_{\leq 0}$, i.e., $f\to \tau^{\mathcal{C}_{/X}}_{\leq 0}f$ is an equivalence. The category of coverings of $X$ is the full subcategory of $\mathcal{C}_{/X}$ consisting of coverings of $X$. We denote this category by $\mathcal{C}ov_{\mathcal{C}}(X):=  ({\mathcal{C}_{/X}})_{\leq 0}$. 
\end{defn}
\begin{prop}\label[prop]{Covering of Localization}
   When $L\mathcal{C}\subset \mathcal{C}$ is a localization, a morphism $X\to Y$ in $L\mathcal{C}$ is a covering in $L\mathcal{C}$ if and only if it is a covering in $\mathcal{C}$.
\end{prop}
\begin{proof}
    When $X\in L\mathcal{C}$, the inclusion $L\mathcal{C}_{/X}\subset \mathcal{C}_{/X}$ is closed under limits. Indeed, if $D:\mathscr{I}\to L\mathcal{C}_{/X}$ is a diagram, then the limit of the composite $\mathscr{I}\to L\mathcal{C}_{/X}\subset \mathcal{C}_{/X}$ is the limit of the extended cone $D^\triangleright:\mathscr{I}^{\triangleright}\to \mathcal{C}$ (mapping the terminal object to $X$), with structure map to $X$ given by the leg to the terminal object $X$. But since $L\mathcal{C}\subset \mathcal{C}$ is closed under limits, this limit belongs to $L\mathcal{C}_{/X}$. Since both $\mathcal{C}_{/X}$ and $L\mathcal{C}_{/X}$ are presentable, it follows that $L\mathcal{C}_{/X}\subset \mathcal{C}_{/X}$ is a localization. Therefore, \Cref{truncation of localization} applies, and it follows that $$ {L\mathcal{C}_{/X}}_{\leq 0}= ({\mathcal{C}_{/X}})_{\leq 0}\bigcap L\mathcal{C}_{/X}.$$ 
\end{proof}

\begin{rem}\label[rem]{Effective localization remark}
For convenience, we say that a localization is effective if, for every object $X$, the localization unit $X\to LX$ is an effective epimorphism. An example of an effective localization is the localization of a sheaf topos at an interval object, in the sense of [\cite{morel19991}]. Indeed, when a sheaf topos $\mathcal{T}$ is localized at an interval object $I$, there is an epimorphism $\tau_{\leq 0}^{\mathcal{T}}\esc{X}\to \tau_{\leq 0}^{\mathcal{T}}L_{I}\esc{X}$ [\cite{morel19991}, Corollary 2.3]. By [\cite{lurie2009higher}, Proposition 7.2.1.14], it follows that $\esc{X}\to L_{I}\esc{X}$ is an effective epimorphism. More generally, one can prove that a cartesian localization at any split object (objects admitting global sections) is effective (see [\cite{p1algtop}]). 
\end{rem} 
\begin{lem}
    An effective localization $L:\mathcal{C}\to \mathcal{C}$ preserves effective epimorphisms. 
\end{lem}
\begin{proof}
   Since $\mathcal{C}$ is an $\infty$-topos, by [\cite{lurie2009higher}, Proposition 7.2.1.14], this follows from the fact that the property of being an epimorphism in an ordinary category has right cancellation.\end{proof}
\begin{lem}\label[lem]{covering lemma}
Let $\mathcal{C}$ be an $\infty$-topos and $L$ a localization of $\mathcal{C}$. Then
    \begin{enumerate}
        \item A morphism of $L\mathcal{C}$-coverings is itself an $L\mathcal{C}$-covering.
        \item Composition of $L\mathcal{C}$ coverings is an $L \mathcal{C}$ covering.
        \item Let $f:Y\to X$ be a morphism in $\mathcal{C}$. Then the pull back functor $\mathcal{C}/LX\to \mathcal{C}/LY$ takes $L\mathcal{C}$ coverings to $L\mathcal{C}$ coverings. 
        \item If $L$ is effective, then the canonical map $\mathcal{C}_{/LX}\to \mathcal{C}_{/X}$ (induced by pullback along the localization unit $X\to LX$) preserves and detects coverings. 
        \item  If $L$ is effective, then a morphism $Z\to LX$ is an $L\mathcal{C}$-covering if and only if $Z$ is $L$-local and $Z\times _XLX\to X$ is a $\mathcal{C}$-covering. 
    \end{enumerate}
\end{lem}
\begin{proof}
(1,2). Using the proposition above, for (1) and (2), it suffices to assume $L=id_\mathcal{C}$. Now, suppose we have a commutative triangle:
\[
\xymatrix{
X\ar[rr]^f\ar[dr]_h&&Y\ar[dl]^g\\
&Z&
}
\]
thought of as a morphism $f: h\to g$ in the slice category $\mathcal{C}_{Z}$. By [\cite{lurie2009higher}, Lemma 5.5.6.14], it follows that when $g$ is $0$-truncated, $h$ is $0$-truncated if and only if $f$ is. 

(3, 4).
When $L$ is effective, i.e., $X\to LX$ is an effective epimorphism for every $X\in \mathcal{C}$, the result [\cite{lurie2009higher}, Proposition 6.2.3.17] applies.

(5). Follows immediately from (4) and the definition of $L\mathcal{C}$ coverings.
\end{proof}
\begin{rem}\label[rem]{covering pullback description}
    Under the conditions of statement 5 of the above proposition, that is, when the localization $L$ is effective, it follows that there is a pullback square:
\[
\xymatrix{
\mathcal{C}ov_{L\mathcal{C}}(LX)\ar@{^(->}[r]\ar[d]&L\mathcal{C}_{/LX}\ar[d]\\
\mathcal{C}ov_{\mathcal{C}}(X)\ar@{^(->}[r] &\mathcal{C}_{/X}
}
\]
\end{rem}
\begin{lem}\label[lem]{covering iff fiber is discrete}
 Suppose $X$ is a pointed, connected object. Then a morphism $f: Z\to X$ in $\mathcal{C}$ is a covering if and only if the fiber of $f$ is a discrete object of $\mathcal{C}$. 
\end{lem}
\begin{proof}
    When $X$ is connected, the pointing map $*\to X$ is an effective epimorphism. As before [\cite{lurie2009higher}, Proposition 6.2.3.17], the result applies again. The result follows from the fact that a morphism $Z\to *$ is $0$-truncated iff $Z$ is $0$-truncated.
\end{proof}
\begin{cor}\label[cor]{L covering iff fiber is L discrete}
Suppose $L$ is closed under extensions over $LX$. Then a morphism $f:Z\to LX$ with $LX$ pointed and connected is an $L\mathcal{C}$ covering if and only if $fib(f)\in Disc(L\mathcal{C})$.
\end{cor} 
\begin{proof}
    It follows that $Z$ is $L$-local. So it suffices to know that $f$ is a $\mathcal{C}$-covering (see \Cref{Covering of Localization}). This is immediate from the Lemma above. 
\end{proof}
\begin{defn}\label[defn]{definition for L covering}
Let $\mathcal{C}$ be a presentable $\infty$-category and $L$ an accessible localization. Let $X\in \mathcal{C}$ be an object. The category of $L$-fibrations over $X$ is the full subcategory of $\mathcal{C}_{/X}$ consisting of objects over $X$ that are local with respect to $L$-local equivalences in $\mathcal{C}_{/X}$. We denote it by $L\mathcal{F}ib_X$. An object $Y\in L\mathcal{F}ib_X$ is said to be an $L$-fibration over $X$. The category of $L$-covering fibrations of $X$, denoted $Cov_L(X)$, is the truncation $\tau_{\leq 0} L\mathcal{F}ib_X$.
\end{defn}


\begin{rem}
Since the forgetful functor $\mathcal{C}_{/X}\to \mathcal{C}$ preserves and creates colimits, the class of $L$-local equivalences in $\mathcal{C}_{/X}$ is again saturated. Hence the inclusion $L\mathcal{F}ib_X\subset \mathcal{C}_{/X}$ is an accessible localization, say given by $L^{fib}$. Therefore, using \Cref{truncation of localization}, we find: $$(L\mathcal{F}ib_X)_{\leq 0} =(\mathcal{C}_{/X})_{\leq 0} \bigcap  L\mathcal{F}ib_X.$$
\end{rem}
\begin{lem}\label[lem]{L fib over L local object}
    A morphism between $L$-local objects is always an $L$-fibration.
    \end{lem}
    \begin{proof}
Given a morphism $f:Y\to X$, let us consider the $L^{fib}$ localization defined by the following triangle
\[
\xymatrix{
X\ar[r]^-{e}\ar[dr]_f& dom(L^{fib}f)\ar[d]^{L^{fib}f}\\
& Y
}
\]
Because $Y$ is $L$-local and $L^{fib}f$ is an $L$-fibration, it follows that $dom(L^{fib}f)$ is $L$-local. Since $e$ is an $L$-equivalence by definition, and both $X$ and $dom(L^{fib}f)$ are $L$-local, it follows that $e$ is an equivalence, and therefore $f$ is equivalent to the $L$-fibration $L^{fib}f$.   \end{proof}
\begin{lem}\label[lem]{lemma for L covering} 
\begin{enumerate}
   \item  The pullback of an $L$-fibration (resp. covering) is a fibration (resp. covering). 
   \item  The composition of $L$ fibrations (resp. coverings) is a fibration (resp. covering). 
\end{enumerate}
\end{lem}
\begin{proof}
    (1). Since the left adjoint $f: \mathcal{C}_{/Y}\to \mathcal{C}_{/X}$ given by precomposition with $f$ clearly preserves $L$-equivalences, the right adjoint $f^*$ preserves the local objects, i.e., $L$-fibrations. On the other hand, the pullback functor $f^*: \mathcal{C}_{/X}\to \mathcal{C}_{/Y}$ is left exact. By [\cite{lurie2009higher}, Proposition 5.5.6.16], it preserves $0$-truncated objects.  
    
    (2). The fact that the composition of $L$-fibrations is an $L$-fibration follows easily. That the composition of $0$-truncated morphisms is $0$-truncated is well known. In fact, this is true in general for $n$-truncated morphisms. For example, if $f: X\to Y$ and $g:Y\to Z$ are $n$-truncated, the claim follows from applying [\cite{lurie2009higher}, Lemma 5.5.6.14] to the morphism $f$ thought of as a morphism from $gf$ to $g$ in the slice category $\mathcal{C}_{/Z}$.
    \end{proof}
\begin{prop}\label[prop]{L covering of L local object}
Let $X\in L\mathcal{C}$ be an $L$-local object. Then $Cov_{L}(X)= Cov_{L\mathcal{C}}(X) $.
\end{prop}
\begin{proof}
    The point is that a morphism $f:Y\to X$ over an $L$-local object $X$ is an $L$-fibration iff $Y\in L\mathcal{C}$ (see \Cref{L fib over L local object}). The claim now follows easily from \Cref{lemma for L covering}.
\end{proof} 
For the rest of this section, assume that $\mathcal{C}$ is an $\infty$-topos and that $L$ is an accessible localization (not necessarily left exact).
\begin{cor}\label[cor]{torsors are covering}
Let $G\in ( Grp^L(\mathcal{C}))$ and let $Y\to X$ be a $G$-torsor [\cite{MR3423073} $\S$3]. Then $Y\to X$ is an $L$-covering. 
\end{cor}
\begin{proof}
    Note that for any $G\in Grp(\mathrm{Disc}{(\mathcal{C}}))$, $*\to \mathrm{B}_\mathcal{C}G$ is a $\mathcal{C}$ covering by [\Cref{L covering iff fiber is L discrete}]. Thus, when $G\in Grp^L\mathcal{C}$ so that $\mathrm{B}_\mathcal{C}G$ is $L$ local, the morphism is an $L\mathcal{C}$ covering. But then, by the above proposition \Cref{L covering of L local object}, it is also an $L$ covering. Because $Y\to X$ is pulled back from $*\to \mathrm{B}_\mathcal{C}G$ along some map $X\to \mathrm{B}_\mathcal{C}G$, \Cref{lemma for L covering} (1) yields the desired result.
\end{proof}


For an $X\in \mathcal{C}$, let $\eta:X\to LX$ be the corresponding localization unit. Pulling back along this unit morphism yields a canonical map $$\eta^*: Cov_{L\mathcal{C}}{LX}=Cov_L(LX)\to Cov_{L}(X).$$ This map fits into the following pullback square (even when $L$ is not effective):
\[
\xymatrix{
\mathcal{C}ov_{L\mathcal{C}}(LX)\ar@{^(->}[r]\ar[d]^{\eta^*}&L\mathcal{C}_{/LX}\ar[d]^{\eta^*}\\
\mathcal{C}ov_{L}(X)\ar@{^(->}[r] &\mathcal{C}_{/X}
}
\]

\begin{prop}\label[prop]{L covering is covering of L}
If $L$ is locally cartesian and effective on $X$, then the functor $$\eta^*_X: Cov_{L\mathcal{C}}{LX}=Cov_L(LX)\to Cov_{L}(X)$$ given by pull back along the localization $\eta_X:X\to LX$ is an equivalence of categories, with inverse given by $L$.
\end{prop}
\begin{proof}
    Given an $L$-covering $f:Y\to X$, consider the $L$-fibration $Lf: LY\to LX$ (see \Cref{L fib over L local object}) and the pullback 
    \[
    \xymatrix{
    Y\ar[r]\ar[dr]& LY\times_{LX}{X}\ar[r]\ar[d]&X\ar[d]\\
   &LY\ar[r]^{Lf}&LX
    }
    \]
Because $L$ is locally cartesian on $X$, the morphism $ Y\to LY\times_{LX}{X}$ is an $L$-equivalence. Since $Lf$ (being a morphism of $L$-local objects) is an $L$-fibration (by \Cref{L fib over L local object}), so is the base change $LY\times_{LX}X\to X$ (use \Cref{lemma for L covering}). Since $f$ (being an $L$-covering) is also an $L$-fibration, the $L$-equivalence $ Y\to LY\times_{LX}{X}$ is, in fact, an equivalence (being a morphism between $L$-fibrations). It follows that the following is a pullback square:
\[
\xymatrix{ Y\ar[r]\ar[d]&X\ar[d]\\
LY\ar[r] &LX}
\]
Since $Y\to X$ is a covering and $X\to LX$ is an effective epimorphism, it follows from [\cite{lurie2009higher}, Proposition 6.2.3.17] that $LY\to LX$ is also a covering. 

From the proof, it is clear that the functor $L:Cov_L(X)\to L\mathcal{C}_{/LX}$ factors through $Cov_{L\mathcal{C}}(LX)$ (which we again denote by $L:Cov_L(X)\to Cov_{L\mathcal{C}}(LX)$) and that $\eta^*L$ is equivalent to the identity. On the other hand, local cartesianity guarantees that the composition $L\eta^*$ is also equivalent to the identity. 
\end{proof}

\begin{rem}\label[rem]{initial object of cov in a topos}
Suppose $\mathcal{C}$ is an $\infty$-topos with an initial object $\emptyset$. If $X\in \mathcal{C}$, the unique map $\emptyset\to X$ is always $0$-truncated and hence forms an (uninteresting) initial object of the category $Cov_\mathcal{C}(X)$. When coverings are pointed (so the empty covering is discarded), $Cov_\mathcal{C}(X)_*$ has an interesting initial object. To see this, let $X\in \mathcal{C}_*$. Consider the zero truncation $\tau_{\leq 0}^{/X}$ (which is a left adjoint) of the pointing map $*\to X$ and denote it by $\tau_{>1}^\mathcal{C}X\to X$. Since $*\to X$ is an initial object of $\mathcal{C}_{/X,*}$, and the truncation functor $$\tau_{\leq 0}^{/X}:\mathcal{C}_{/X,*}\to Cov_\mathcal{C}(X)_*$$ preserves colimits, it follows that $\tau_{>1}^\mathcal{C}X\to X$ is an initial object of $Cov_{\mathcal{C}}(X)_*$. In fact, by computing homotopy groups, one can show that this is part of a fiber sequence $$\tau^\mathcal{C}_{>1}X\to X\to \tau_{\leq 1}^\mathcal{C}X.$$
\end{rem}
\begin{prop}\label[prop]{Universal cover for fundamental localization}
Suppose $\mathcal{C}$ is an $\infty$-topos with an accessible localization $L$. Suppose moreover that $X\in \mathcal{C}_*$ is such that $\tau_{\leq 1}LX\simeq \tau_{\leq 1}^{L\mathcal{C}}LX$. Then the category $\mathcal{C}ov_{L\mathcal{C}}(LX)_*$ has an initial object $\tilde{X}:=\tau_{>1}LX$. 
\end{prop}
\begin{proof} 
Clearly, $\mathcal{C}ov_{L\mathcal{C}}(LX)\subset \mathcal{C}ov_{\mathcal{C}}(LX)$. Therefore, it suffices to show that the initial object of $\mathcal{C}ov_{\mathcal{C}}(LX)$ is $L$-local. This initial object is the 1-connected covering $\tau_{>1}^\mathcal{C}LX\to LX$. In the fiber sequence $$\tau_{>1}^\mathcal{C}LX\to LX\to \tau^\mathcal{C}_{\leq 1}LX,$$ the term $\tau^\mathcal{C}_{\leq 1}LX$ is $L$-local. Since $L\mathcal{C}\subset \mathcal{C}$ is closed under limits, we find that $\tau_{>1}^\mathcal{C}LX$ is $L$-local and hence $$(\tau_{>1}^\mathcal{C}LX\to LX)\in \mathcal{C}ov_{L\mathcal{C}}(L
X).$$
\end{proof}
We will call $\tau_{\leq 1}^\mathcal{C}Y$ the fundamental groupoid of $Y$ and denote it by  
$\Pi_1^\mathcal{C}(Y)$. We will use the following notation:
\begin{flalign*}
\Pi^LY:=\tau_{\leq 1}^\mathcal{C}LY\text{ and }\Pi^{L\mathcal{C}}Y:=\tau_{\leq 1}^{L\mathcal{C}}LY\simeq L_{\leq 1}\tau^\mathcal{C}_{\leq 1}LY\simeq L_{\leq 1}\Pi^{L}Y.
\end{flalign*}

\begin{ex}
For a space $\esc{X}\in\mathcal{P}(Sm_S)$ and a Grothendieck topology $\sigma$,  
 the ordinary $\sigma$-local fundamental groupoid will be denoted by $\Pi_1^\sigma\esc{X}$. Similarly, when $L=L_{mot} $ and $L=L_{bir} $, the corresponding fundamental groupoids will be denoted by $\Pi_1^{\mathbb{A}^1}\esc{X}$ and $\Pi_1^{b}\esc{X}$, respectively. 
\end{ex}
The condition in the above proposition is then read as $\Pi^LX\simeq \Pi^{L\mathcal{C}}X$. While it might be possible to deduce more general statements without assuming this equivalence, at this point we will work with this condition.  
\begin{defn}
    We say that a localization $L$ is fundamental on $X$ if $\Pi^LX\simeq \Pi^{L\mathcal{C}}X$. We say that $L$ is fundamental if it is fundamental on every object.
\end{defn} 
\begin{lem}
    Suppose $L$ is fundamental on a pointed object $y:*\to Y$. Then 
    $$\pi_1^L(Y,y)=\pi_1^\mathcal{C}(LY,Ly)\simeq \pi_1^{L\mathcal{C}}(LY,Ly).$$
    \end{lem}
    \begin{proof}
        This follows from the fact that $$\pi_1^\mathcal{C}(LY,Ly)\simeq \Omega \tau_{\leq 1}^{\mathcal{C}}LY\simeq \Omega \tau_{\leq 1}^{L\mathcal{C}}LY,$$
        and that the loop space of a local object is local.
    \end{proof}
 Suppose $Y\in \mathcal{C}$ admits a global section $y: *\to Y$. Then, by \Cref{covering lemma} (3), pullback along $y$ induces a functor $$y^*:Cov_{L\mathcal{C}}(LY)\to Cov_L(*)=Disc(LC).$$

 \begin{lem}\label[lem]{action of the fundamerntal groupoid}
     Let $Z\to Y$ be a covering in $\mathcal{C}$. Then the fiber $y^*Z$ has an action by $\pi_1(Y,y)$. When $\Pi^LX\simeq \Pi^{L\mathcal{C}}X$, $y^*$ of an $L\mathcal{C}$ covering is a $\pi_1^L(Y,y)=\pi_1^\mathcal{C}(LY,Ly)\simeq \pi_1^{L\mathcal{C}}(LY,Ly)$ module. 
 \end{lem}
\begin{proof}
Again, we may reduce to $L=id$. By applying $Map(W,-)$ and using the equivalence of the Yoneda functor $\mathcal{C}\to \infty \mbox{-}Top(\mathcal{C}^{op},Spc)$, we reduce this to the case of spaces, where it is well known that $\Omega Y$ acts coherently on $y^*Z$, inducing the required action on $\pi_0$.
\end{proof}


Thus, there is a canonical functor $$\Gamma_y^L:Cov_{L\mathcal{C}}(LY)\longrightarrow \pi_1^L(Y,y)\mbox{-}Disc(L\mathcal{C})$$ provided $L$ is fundamental on $Y$.
\begin{prop}\label[prop]{L galois theory}
Suppose $L\mathcal{C}$ is closed under extensions over $LY$. When $LY$ is connected, the functor $$\Gamma_y^L:Cov_{L\mathcal{C}}(LY)\longrightarrow \pi_1^L(Y,y)\mbox{-}Disc(L\mathcal{C})$$ is an equivalence. 
\end{prop}
\begin{proof}
When $L=id$, the inverse is given by $$F\mapsto \big (\mathrm{B}_\mathcal{C}(F,\Omega Y,*)\to \mathrm{B}_\mathcal{C}(*,\Omega Y,*)\simeq Y \big ).$$ Since the morphism $\mathrm{B}_\mathcal{C}(F,\Omega Y,*)\to \mathrm{B}_\mathcal{C}(*,\Omega Y,*)\simeq Y$ has a canonical fiber equal to the discrete space $F$, connectivity of $Y$ implies that it is a covering (see \Cref{covering iff fiber is discrete}).

  When $L$ is nontrivial, it suffices to know the $L$-local counterpart, namely \Cref{L covering iff fiber is L discrete}.
\end{proof}
Before we state the next result, it is worth noting that when $X$ is an $n$-truncated object in a presentable $\infty$-category $\mathcal{D}$, then for any object $Y$, the object $\underline{Map}_\mathcal{C}(Y,X)$ is also $n$-truncated. Hence, the following makes sense, where $\underline{End}_\mathcal{D}(X):=\underline{Map}_\mathcal{D}(X)$.
\begin{lem}\label[lem]{aut=pi_1}
Let $L$ be a localization of an $\infty$-topos $\mathcal{C}$. Suppose $L$ is locally cartesian and fundamental on a pointed object $X$. Then the canonical map $  \pi_1^LX\to \underline{End}_{L\mathcal{C}_{/LX}}(\tilde{X})$ in $\mathcal{C}$ is an isomorphism. It follows that this induces isomorphisms in $\mathcal{C}$ $$ \pi_1^{L\mathcal{C}}LX\simeq  \pi_1^LX\cong\underline{Aut}_{L\mathcal{C}_{/LX}}(\tilde{X})\cong\underline{End}_{L\mathcal{C}_{/LX}}(\tilde{X}).$$
\end{lem}
\begin{proof}
Again, it suffices to prove this for $L=id_\mathcal{C}$, since $$\underline{End}^{L\mathcal{C}}_{LX}(\tilde{X})\simeq \underline{End}^{\mathcal{C}}_{LX}(\tilde{X})\text{ and }\pi_1^L= \pi_1^\mathcal{C}L.$$ Indeed, for any $Z\in \mathcal{C}_{/X}$,
\begin{flalign*}
    \mathrm{Map}_{\mathcal{C}_{/LX}}(Z,\underline{End}_{L\mathcal{C}_{/LX}}(\tilde{X}))&\simeq \mathrm{Map}_{L\mathcal{C}_{/LX}}(LZ,\underline{End}_{L\mathcal{C}_{/LX}}(\tilde{X}))\\
    &\simeq \mathrm{Map}_{L\mathcal{C}_{/LX}}(LZ\times _{LX}\tilde{X},\tilde{X})\\
    &\simeq \mathrm{Map}_{L\mathcal{C}_{/LX}}(L(Z\times _{LX}\tilde{X}),\tilde{X})\text{ [since } L \text{ is locally cartesian]}\\
    &\simeq \mathrm{Map}_{\mathcal{C}_{/LX}}(Z\times _{LX}\tilde{X},\tilde{X})\\
    &\simeq \mathrm{Map}_{\mathcal{C}_{/LX}}(Z,\underline{End}_{\mathcal{C}_{/LX}}(\tilde{X})
\end{flalign*}


So we are reduced to proving that $$Aut_{\mathcal{C}_{/Y}}(\tau_{>1}Y)\simeq \underline{End}_{\mathcal{C}_{/Y}} 
(\tau_{>1}Y)\cong \pi_1^\mathcal{C}Y$$ for an object $Y\in \mathcal{C}$.
Let $u: Y \to \tau_{\le 1}Y$ be the canonical truncation with basepoint $f: * \to \tau_{\le 1}Y$. We wish to compute the absolute mapping object $End_{\mathcal{C}_{/Y}}(\tau_{>1}Y) \in \mathcal{C}$.

Let $i: \tau_{>1}Y \to Y$ denote the canonical inclusion. By definition of morphisms in the slice category $\mathcal{C}_{/Y}$, we obtain a fiber sequence in $\mathcal{C}$:
$$ \underline{End}_{\mathcal{C}_{/Y}}(\tau_{>1}Y) \to \underline{Map}_{\mathcal{C}}(\tau_{>1}Y, \tau_{>1}Y) \xrightarrow{i \circ -} \underline{Map}_{\mathcal{C}}(\tau_{>1}Y, Y) $$
taken over the structure map $i$.

Next, we map the object $\tau_{>1}Y$ into the defining fiber sequence of the connected cover:
$$ \tau_{>1}Y \xrightarrow{i} Y \xrightarrow{u} \tau_{\le 1}Y. $$
Because mapping objects in an $\infty$-topos preserve limits, this induces another fiber sequence in $\mathcal{C}$:
$$ \underline{Map}_{\mathcal{C}}(\tau_{>1}Y, \tau_{>1}Y) \xrightarrow{i \circ -} \underline{Map}_{\mathcal{C}}(\tau_{>1}Y, Y) \xrightarrow{u \circ -} \underline{Map}_{\mathcal{C}}(\tau_{>1}Y, \tau_{\le 1}Y). $$
By the adjunction properties of truncations, $\tau_{>1}Y$ is $>1$-connected and $\tau_{\le 1}Y$ is $1$-truncated. Thus, any map between them factors uniquely through the terminal object. $*$:
$$ \underline{Map}_{\mathcal{C}}(\tau_{>1}Y, \tau_{\le 1}Y) \simeq \underline{Map}_{\mathcal{C}}(*, \tau_{\le 1}Y) \simeq \tau_{\le 1}Y. $$
The composition $u \circ i$ is canonically identified with the basepoint $f: * \to \tau_{\le 1}Y$. Therefore, the fiber of the map $u \circ -$ is precisely the loop space of $\tau_{\le 1}Y$ at the basepoint $f$.

By comparing the fiber sequences, we see that the object $\underline{End}_{\mathcal{C}_{/Y}}(\tau_{>1}Y)$ is exactly the fiber of $u \circ -$ over $f$. Hence:
$$ \underline{End}_{\mathcal{C}_{/Y}}(\tau_{>1}Y) \simeq \Omega_f \tau_{\le 1}Y \simeq \pi_1^\mathcal{C}(Y). $$

Finally, because $\tau_{>1}Y$ is an object of $\mathcal{C}_{/Y}$, its endomorphism object $\underline{End}_{\mathcal{C}_{/Y}}(\tau_{>1}Y)$ naturally inherits a monoid structure in $\mathcal{C}$. Since $\pi_1^\mathcal{C}(Y)$ is inherently group-like, the inclusion of the maximal subgroupoid is an equivalence. This yields the desired isomorphism:
$$ \underline{Aut}_{\mathcal{C}_{/Y}}(\tau_{>1}Y) \simeq \underline{End}_{\mathcal{C}_{/Y}}(\tau_{>1}Y) \simeq \pi_1^\mathcal{C}(Y). $$
\end{proof}

\section{The internal homotopy theory of \texorpdfstring{$\mathcal{H}^{mot}$}{h}}

In this section, we apply the definitions and results from the previous section to the motivic localization functor $L_{mot}:\mathcal{P}_{nis}(S)\to\mathcal{H}^{mot}(S) $. Before we do that, let us have a quick recap of how this localization is constructed.

Let $S$ be a Qcqs scheme. By an $S$-space, we mean a presheaf of spaces on $Sm_S$. We denote by $\mathcal{P}(S)$ the $\infty$-category of $S$-spaces. We say that an $S$-space is Nisnevich-local if it takes Nisnevich cd squares to pullback squares of spaces (and the empty scheme to the contractible space). Recall that a Nisnevich cd square is a pullback square:

\[
\xymatrix{
\ar@{}[d]_{Q:}&V\ar@{^(->}[r]\ar[d]&Y\ar[d]^p\\
&U\ar@{^(->}[r]_j&X
}
\]
of schemes smooth over $S$, where $p$ is etale and $j$ is an open immersion such that the complements induce an isomorphism with the reduced scheme structure. We denote the full subcategory of $\mathcal{P}(S)$ consisting of 
Nisnevich local spaces by $\mathcal{P}_{nis}(S)$. It is easy to see that the Nisnevich locality condition is equivalent to the condition of being local with respect to the following morphism of $k$-spaces:
$$\{h_SU\bigsqcup_{h_SV}h_SY\to h_SX\mid Q \text{ is a Nisnevich square}\}\bigcup \{\emptyset\to h_S\emptyset\}$$ 
Since the collection above is small, it follows that $\mathcal{P}_{nis}(S)\subset \mathcal{P}(S)$ is an accessible localization [\cite{lurie2009higher}, Proposition 5.5.1.14]. We denote the corresponding localization functor by $L_{nis}$. It turns out that $L_{nis}$ is left exact, and thus $\mathcal{P}_{nis}(S)$ is an $\infty$-topos [\cite{hoyois2017six}, Corollary 3.9].


We say that an $S$-space $\esc{X}$ is $\mathbb{A}^1$-local if, for every $X\in Sm_S$, the restriction map $\esc{X}(X)\xrightarrow{p^*}\esc{X}(\mathbb{A}^1\times X)$, induced by the projection $\mathbb{A}^1\times X\xrightarrow{p} X$, is an equivalence of spaces. Again, since the collection of $\mathbb{A}^1$-projections is a small set, it follows that the full subcategory of $\mathcal{P}(S)$ consisting of $\mathbb{A}^1$-local objects is an accessible localization, which we denote by $L_{\mathbb{A}^1}$.

Finally, an $S$-space is called motivic if it is both Nisnevich local and $\mathbb{A}^1$-local. We denote the full subcategory of $\mathcal{P}(S)$ consisting of motivic $S$-spaces by $\mathcal{H}^{mot}(S)$. Because both conditions are given by small sets, it follows that the subcategory $\mathcal{H}^{mot}(S)\subset \mathcal{P}(S)$ is an accessible localization. This localization is denoted by $L_{mot}$. It is a cartesian, locally cartesian, and accessible localization. 

Morel defines the $\mathbb{A}^1$-homotopy group of a space as the Nisnevich homotopy group of its motivic localization, i.e., $$\pi_i^{\mathbb{A}^1}:=\pi_i^{nis}L_{mot}:=L_{nis}\pi_iL_{mot}.$$ It is evident that, in terms of \Cref{pinL}, this is precisely $\pi_i^{L_{mot}}$. 
Nonetheless, it is evident that Morel's $\mathbb{A}^1$-homotopy groups of a motivic space are simply its Nisnevich homotopy group sheaves. 
 
Unfortunately, the $\mathbb{A}^1$-homotopy path component construction $\pi_0^{\mathbb{A}^1}$ is not as well behaved categorically. In particular, it fails to be the internal $0$-truncation functor for the motivic homotopy category, as it is in an $\infty$-topos. Consequently, it does not preserve (motivic) colimits. The agenda of this section is to elaborate on this issue and correct it by introducing the internal notions of motivic truncations and motivic homotopy groups.



\subsection{Motivic truncations and Motivic homotopy groups} 

For the convenience of the motivic readers (those who are reading this section directly without reading \S2), we shall first restate the definitions of \S2.1 for the motivic localization $L_{mot}:\mathcal{P}(S)\to \mathcal{H}^{mot}(S)$. To simplify notation, we shall drop the scheme $S$, assuming it to be fixed. This assumption shall continue until \Cref{strong A1 invariance of pin}, after which we shall assume that the base is a field (though not necessarily perfect, unless otherwise specified).

\textit{Motivic truncations:}

\begin{defn}\label[defn]{motivically n truncated space}
    We say that a motivic space $\esc{X}$ is motivically $n$-truncated if, for every motivic space $\esc{Y}$, the space $\mathrm{Map}_{\mathcal{H}^{mot}}(\esc{Y},\esc{X})$ is $n$-truncated. We denote the full subcategory of $n$-truncated motivic spaces by $\mathcal{H}^{mot}_{\leq n}$ or, equivalently, by $\tau_{\leq n}^{mot}\mathcal{H}^{mot}$.
\end{defn}
 From \Cref{truncation of localization}, it follows that:
 \begin{prop}\label[prop]{motivic trunc is nis trunc}
     A motivic space is motivically $n$-truncated if and only if it is Nisnevich locally $n$-truncated, i.e., $$\mathcal{H}^{mot}_{\leq n}=\mathcal{H}^{mot}\bigcap\mathcal{P}^{nis}_{\leq n}.$$
\end{prop}
\begin{rem}
     Recall that, when $n\geq 0$ and $S$ has finite Krull dimension, by \Cref{truncated condition} this last condition is equivalent to the requirement that for all $i>n$ the Nisnevich-local space has trivial $\pi_i^{nis}$.
 \end{rem}
 Since $\mathcal{P}$ is an $\infty$-topos, \Cref{L tunc is a localization} implies that $\mathcal{H}^{mot}_{\leq n}\subset \mathcal{P}$ is also a localization. We let $$L_{\leq n}^{mot}: \mathcal{P}\to \mathcal{H}^{mot}_{\leq n}$$
stand for the induced motivic truncation functors. In particular, these restrict to the internal motivic truncation functors $$\tau_{\leq n}^{mot}: \mathcal{H}^{mot}\to \mathcal{H}^{mot}_{\leq n},$$ i.e., $\tau_{\leq n}^{mot}$ is the localization functor for $\mathcal{H}^{mot}_{\leq n}\subset \mathcal{H}^{mot}$.

\begin{con}\label{motivic truncation formulae}
      To set this in concrete terms, let us discuss how one constructs a formula for $L_{\leq n}^{mot}$.
      
    We keep in mind that, from \Cref{truncation of localization}, $$\mathcal{H}^{mot}_{\leq n}=\mathcal{H}^{mot}\bigcap\mathcal{P}_{\leq n} =L_{\mathbb{A}^1}\mathcal{P}\bigcap\mathcal{P}_{nis}\bigcap \mathcal{P}_{\leq n}.$$ The idea is then standard: define 
    $$L_{\leq n}^{mot}:=\colim_{m}\Big (L_{nis}\tau_{\leq n}L_{\mathbb{A}^1}\Big)^m$$ where $L_{nis}:\mathcal{P}\to\mathcal{P}^{{nis}}$ is the nisnevich $\infty$-sheafification, and $\tau_{\leq n}: \mathcal{P}\to\mathcal{P}_{\leq n}$ is the ordinary Postnikov truncation of the presheaf topos (which can, for example, be modeled sectionwise by any simplicial truncation model, such as the Moore-Postnikov model $P_n$) and $L_{\mathbb{A}^1}: \mathcal{P}\to L_{\mathbb{A}^1}\mathcal{P}$ is the (presheaf) $\mathbb{A}^1$-localization functor (which can be modeled by taking the singular Suslin-type construction). In fact, since $\tau_{\leq n}^{nis}=L_{nis}\tau_{\leq n}$, this formula reads $L_{\leq n}^{mot}=\colim_{m} (\tau^{nis}_{\leq n}L_{\mathbb{A}^1})^m$.
\end{con}
\begin{thm}\label{formula for motivic truncation}
    The construction above provides the desired localization functor $$L_{\leq n}^{mot}: \mathcal{P}\to \mathcal{H}^{mot}_{\leq n}.$$
\end{thm}
\begin{proof}
    To see why this works, first observe that all three subcategories in question are closed under filtered colimits. Indeed, $\mathcal{P}_{nis}\subset\mathcal{P}$ is closed under filtered colimits because the nisnevich topology is a cd topology; $L_{\mathbb{A}^1}\mathcal{P}\subset\mathcal{P}$ is closed under all small colimits, as can be checked on sections; and $ \mathcal{P}_{\leq n}\subset\mathcal{P}$ is closed under filtered colimits, which can be proven by induction: the base case $n=-2$ follows from [\cite{lurie2009higher}, Definition 5.5.6.1], while the inductive step follows from [Proposition 5.5.6.5, \cite{lurie2009higher}] and the commutation of pullbacks and filtered colimits in toposes. A cofinality argument then shows that $$\dcolim_{m}(L_{nis}\tau_{\leq n}L_{\mathbb{A}^1}\esc{X})^m\in \mathcal{P}^{nis}\bigcap L_{\mathbb{A}^1}\mathcal{P}\bigcap\mathcal{P}_{\leq n}=\mathcal{H}^{mot}\bigcap\mathcal{P}_{\leq n}=\mathcal{H}^{mot}_{\leq n}.$$ That the natural transformation $$1\to \colim_{m}(L_{nis}\tau_{\leq n}L_{\mathbb{A}^1})^m$$ is an $L_{\leq n}^{mot }$ equivalence follows from the definition of associated saturated classes.
\end{proof}


\begin{rem}
    It follows that $\tau_{\leq n}^{mot}\simeq {L_{\leq n}^{mot}}_{|\mathcal{H}^{mot}}$.
\end{rem}
\begin{rem}\label[rem]{remark on truncation of localization}
    From the arguments of the above proof, it is not hard to see that the same truncation formula can equivalently be given by:
    $$L_{\leq n}^{mot}=\colim_{m} (\tau^{nis}_{\leq n}L_{mot})^m=\colim_{m} (\tau_{\leq n}L_{mot})^m.$$ 
\end{rem}

The following is not generally true for arbitrary localizing categories. In fact, it is a property of an $\infty$-topos [\cite{lurie2009higher}, Lemma 6.5.1.2]:  
 \begin{prop}
     The localization $L_{\leq n}^{mot}$ is effective, i.e., for every $S$-space $\esc{X}$, the map $\esc{X}\to L_{\leq n}^{mot}\esc{X}$ is an effective epimorphism.
 \end{prop} 
 \begin{proof}
  By [\cite{lurie2009higher}, Proposition 7.2.1.14], it suffices to prove that the morphism is a surjection of Nisnevich sheaves under $\pi_0^{nis}$. But this morphism of Nisnevich sheaves is identified with $\pi_0^{nis}\esc{X}\to \colim_{m}\pi_0^{\mathbb{A}^1}\esc{X}$. The claim thus follows from [\cite{morel19991}, Corollary 1.2.3].
 \end{proof}
 \begin{rem}
    It is clear from the above proof that:
    $$\tau_{\leq 0}^{mot}=\colim_{m}(\pi_0^{\mathbb{A}^1})^{\circ m}.$$
\end{rem} 

\begin{prop}
    $L^{mot}_{\leq n}$ preserves products. 
\end{prop}
\begin{proof}
    We use the formula constructed above. Since filtered colimits commute with finite products, it suffices to observe that each component functor in the formula has the desired property. The case of $L_{nis}$ is standard (it is, in fact, left exact); the case of $L_{\mathbb{A}^1}$ is well known (using the fact that $\mathbb{A}^1$-projections are stable under arbitrary base change); for $\tau_{\leq n}$, use [\cite{lurie2009higher}, Lemma 6.5.1.2].
\end{proof}
\begin{cor}\label[cor]{mot truncation cart}
    $\tau^{mot}_{\leq n}$ preserves products. 
\end{cor}

\begin{cor}
    $L_{mot}$ is a path cartesian localization (see \Cref{path cart}).
\end{cor}
In the particular case of discrete motivic spaces, one finds $$\mathrm{Disc}(\mathcal{H}^{mot})=\mathcal{H}^{mot}\bigcap disc(\mathcal{P}_{nis})=\mathcal{H}^{mot}\bigcap Sh_{nis}(Sm_S, Sets)=Sh_{nis}^{\mathbb{A}^1}$$ (here $Sh_{nis}^{\mathbb{A}^1}$ is the category of $\mathbb{A}^1$-invariant nisnevich sheaves of sets on $Sm_S$). Following the above machinery, a localization for the inclusion $Sh_{nis}^{\mathbb{A}^1}\subset \mathcal{P}$ can be given by the formula:
$$L_{\leq 0}^{mot}:=\colim_{m} (L_{nis}\tau_{\leq 0}L_{\mathbb{A}^1})^m.$$

A well-developed version of this was obtained for sheaves of sets in [\cite{amit}], where it was denoted by $\mathcal{L}$. 
\begin{cor}
    For motivic spaces $\esc{X}$, one has the sheaf isomorphism $\mathcal{L}\pi_0^{nis}\esc{X}\cong \pi_0^{mot}\esc{X}$.
\end{cor}
\begin{proof}
    By taking the $\mathbb{A}^1$-singular formulae $L_{\mathbb{A}^1}: = |(-)^{\mathbb{A}^\bullet}|$, one finds that $$\tau_{\leq 0}|(-)^{\mathbb{A}^\bullet}|\simeq |(\tau_{\leq 0}(-))^{\mathbb{A}^\bullet}|\cong coeq\big ((-)^{\mathbb{A}^1}\rightrightarrows  (-)\big ).$$ But this is precisely the formulae obtained in [\cite{amit}] for presheaves of sets, which they called $\mathcal{S}^{pre}$. Consequently, $$L_{nis}\tau_{\leq 0}L_{\mathbb{A}^1}=L_{nis}\pi_0^{ch}=\mathcal{S}.$$ Thus their formulae $\mathcal{L}:=\colim_{m}\mathcal{S}^m\equiv \colim_{m} (L_{nis}\tau_{\leq 0}L_{\mathbb{A}^1})^m $ is indeed the motivic truncation functor $$L_{\leq 0}^{mot}: Sh_{nis}\to disc(\mathcal{H}^{mot})=Sh_{nis}^{\mathbb{A}^1}$$ as constructed in \Cref{motivic truncation formulae}. 
\end{proof}
The actual motivic truncation is thus given by $\tau_{\leq 0}^{mot}:={L_{\leq 0}^{mot}}_{|\mathcal{H}^{mot}}\simeq \mathcal{L}({\pi_0^{nis}}_{|{\mathcal{H}^{mot}}})$.

\textit{Motivic homotopy groups: } 

Let $\pi_i^{mot}$ be the internal homotopy groups of $\mathcal{H}^{mot}(S)$, following \Cref{LC homotopy group}. By definition, it is the internal motivic $0$-truncation of the $i$-fold loop space: $$\pi_i^{mot}:=\tau_{\leq 0}^{mot}\Omega_{\mathcal{H}^{mot}(k)}^i\simeq \tau_{\leq 0}^{mot}\Omega^i.$$

\begin{cor}\label[cor]{piimot is Lpiinis}
For all $i\geq 0$ and for $\esc{X}$ a motivic space, we have $$\pi_i^{mot}\esc{X}=\tau_{\leq 0}^{mot}\pi_i^{nis}\esc{X}\simeq \mathcal{L}\pi_i^{nis}\esc{X}.$$ where $\mathcal{L
}$ is a localization for the inclusion $Sh^{\mathbb{A}^1}_{nis}\subset Sh_{nis}$.
\end{cor}
\begin{proof}
    This is \Cref{homotopy groups of localization} applied to $L_{mot}$, together with the observation ${L_{\leq 0}^{mot}}\simeq \mathcal{L}({\pi_0^{nis}})$.
\end{proof}
(For the rest of this section, we shall work over fields. Thus, we let $S=Speck$ for a field $k$.)

\begin{rem}\label[rem]{strong A1 invariance of pin}
     In [\cite{morel2012a1}], Morel proved that, at least over perfect fields, all higher Nisnevich homotopy sheaves of a motivic local space are $\mathbb{A}^1$-invariant (in fact, strongly/strictly!), rendering a further discrete motivic localization unnecessary. Following this, he conjectured that the same would hold even for $\pi_0$. Since $\mathcal{L}$ is a localization for the inclusion $Sh^{\mathbb{A}^1}_{nis}\subset Sh_{nis}$, this conjecture is equivalent to the statement that $$\pi_0^{mot}\simeq {\mathcal{L}\pi_0^{nis}}_{|\mathcal{H}^{mot}}\simeq {\pi_0^{nis}}_{|\mathcal{H}^{mot}}$$ [\cite{morel2012a1}, Conjecture 1.12]. Thus, in the early stages of motivic homotopy theory, no importance was given to the internal truncation functor $L_{\leq 0}^{mot}$ and internal homotopy groups.
 \end{rem}
 
 Though it was known that the above conjecture is very close to being true:
        \begin{lem}
   If $\esc{X}$ is a $k$-motivic space over a perfect field $k$, then $\pi_0^{nis}\esc{X}$ is $\mathbb{A}^1$-invariant if and only if it is weakly unramified if and only if, for any irreducible $X\in Sm_k$, the map $\pi_0^{nis}\esc{X}(X)\to \pi_0^{nis}\esc{X}(k(X))$ is injective. 
\end{lem}
\begin{proof}
  If $\pi_0^{nis}\esc{X}$ is $\mathbb{A}^1$-invariant, then it is weakly unramified by [\cite{bachmann2024strongly}, Corollary 2.4].
  
  Conversely, if $\pi_0^{nis}\esc{X}$ is weakly unramified, then its $\mathbb{A}^1$-invariance follows from [\cite{choudhury2014connectivity}, Corollary 3.2 and Lemma 4.14].
\end{proof}
It was known that $\pi_0^{nis}\esc{X}(X)\to \pi_0^{nis}\esc{X}(k(X))$ has a trivial fiber for any pointed motivic space $\esc{X}$. Building on this observation and using the lemma above, Choudhury proved in [\cite{choudhury2014connectivity}] that, for motivic spaces with a minimal amount of group structure, the Nisnevich homotopy path component $\pi_0^{nis}$ is indeed $\mathbb{A}^1$-invariant:
\begin{prop}\label[prop]{pi0 of h space}
       Let $\esc{X}$ be a $k$-motivic $H$-group space (i.e., a group object in the homotopy $1$-category $h\mathcal{H}^{\mathbb{A}^1}(k)$) over a perfect field $k$. Then $\pi_0^{mot}\esc{X}\simeq \pi_0^{nis}\esc{X}$. 
   \end{prop} 
   \begin{proof}
       This is [\cite{choudhury2014connectivity}, Theorem 4.18] using \Cref{piimot is Lpiinis}.
   \end{proof}
Unfortunately, despite all that, it is now known, thanks to Ayoub, that there are motivic spaces $\esc{X}$ such that $\pi_0^{nis}\esc{X}$ is not $\mathbb{A}^1$-invariant [\cite{ayoubcounterexamples}]. Thus motivic truncation functors $L_{\leq n}^{mot}$ are unavoidable. At best, what is known is:
\begin{prop}\label[prop]{pi0 A1 inv in one vsariable}
    Let the base be a field $k$.  Then for a $k$-motivic space $\esc{X}$, the canonical map $$\pi_0^{\mathbb{A}^1}\esc{X}\simeq \pi_0^{nis}\esc{X}\to \pi_0^{mot}\esc{X}$$ 
    is an isomorphism on $\mathcal{F}_k$ and on $\mathbb{A}^1_{\mathcal{F}_k}$.
\end{prop}

\begin{proof}
    This is a restatement of [\cite{choudhury2022characterisation}, Corollary 2.15].
\end{proof} 


\begin{rem}
    Since $\tau_{\leq0}^{mot}$ preserves products by \Cref{mot truncation cart}, it follows that $\pi_1^{mot}$ is a group object in $Sh^{\mathbb{A}^1}_k$ and that $\pi_i^{mot}$ is an abelian group object in $Sh_k^{\mathbb{A}^1}$. Although, these observations are immediate when $k$ is perfect:
\end{rem}

  \begin{prop}\label[prop]{pi_i mot=nis}
Let $k$ be a perfect field. Then for all $i\geq 1$ and for every $k$-motivic space $\esc{X}$, there are equivalences $$\pi_i^{mot}\esc{X}\simeq \pi_i^{nis}\esc{X}\simeq \pi_i^{\mathbb{A}^i}\esc{X}.$$      
  \end{prop}
\begin{proof}
By the corollary above, this is an immediate consequence of the fact that for every $i\geq 1$ the sheaf $\pi_i^{\mathbb{A}^1}\esc{X}$ is $\mathbb{A}^1$-invariant [\cite{morel2012a1}, Corollary 6.3].  \end{proof}
 \begin{prop}\label[prop]{layers of connected motivic space}
Let $\esc{X}\in  \mathcal{P}_{nis}(k)_{\geq 1}$ be pointed. Then the following conditions are equivalent:
\begin{enumerate}
    \item $\esc{X}\in \mathcal{H}^{mot}(k)$.
    \item $\mathrm{K}_{nis}(\pi_i^{nis}\esc{X},i)\in \mathcal{H}^{mot}(k)$ for all $i$.
    \item $\tau_{\leq i}^{nis}\esc{X}\in \mathcal{H}^{mot}(k)$ for all $i$.
\end{enumerate}
 In this situation, $\tau^{nis}_{\leq i}\esc{X}\simeq \tau^{mot}_{\leq i}\esc{X}$ for all $i\geq 0$.
 \end{prop}
\begin{proof}
The equivalence of the three conditions is well known (see [Corollary 6.3 in \cite{morel2012a1}]); we merely repeat it to set the stage. 

$(1)\implies (2)$ follows from [Corollary 6.2, \cite{morel2012a1}].   
 
Suppose (2) holds, i.e., for all $i$ we have $\mathrm{K}_{nis}(\pi_i^{nis}\esc{X},i)\in \mathcal{H}^{mot}(k)$. Since $\esc{X}$ is nisnevich locally connected, we have $\tau_{\leq 0}^{nis}\esc{X}\simeq *\in \mathcal{H}^{mot}(k)$. This proves the base case of $(3)$. Suppose inductively that $\tau_{\leq i}^{nis}\esc{X}\in \mathcal{H}^{mot}(k)$ for some $i$. Consider the fiber sequence $$\mathrm{K}^{nis}(\pi^{nis}_{i+1}(\esc{X}),i+1)\to \tau^{nis}_{\leq i+1}\esc{X}\to\tau^{nis}_{\leq i}\esc{X}$$
Because of (2), we know that the fiber $\mathrm{K}^{nis}(\pi^{nis}_{i+1}(\esc{X}),i+1)$ lies in $\mathcal{H}^{mot}(k)$. By [\cite{morel2012a1}, Lemma 6.51], it then follows that $\tau^{nis}_{\leq i+1}\esc{X}$ lies in $\mathcal{H}^{mot}(k)$. 

Finally, if (3) holds, i.e., $\tau^{nis}_{\leq n}\esc{X}\in \mathcal{H}^{mot}(k)$. By the Postnikov convergence of $\mathcal{P}_{nis}(Sm_k)$, we have $$\esc{X}\simeq  \plim_{n}\tau^{nis}_{\leq n}\esc{X}.$$ Since the inclusion $\esc{H}^{mot}(k)\subset \mathcal{P}_{nis}(Sm_k)$ is closed under limits, it follows that $\plim_{n}\tau^{nis}_{\leq n}\esc{X}\in \mathcal{H}^{mot}(k)$, and consequently, $\esc{X}\in \mathcal{H}^{mot}(k)$, i.e., (1) holds.

Now we come to the final statement. From \cref{formula for motivic truncation} and \Cref{motivic truncation formulae}  we know that $\tau^{mot}_{\leq i}\simeq\colim_{m} (\tau^{nis}_{\leq n}L_{\mathbb{A}^1})^m$. Since $\esc{X}$ is motivic local, it follows that $\tau^{nis}_{\leq i}L_{\mathbb{A}^1}\esc{X}\simeq \tau_{\leq i}^{nis}\esc{X} $, which from the second condition is motivic local as well. Hence for all $n$ we have  $(\tau^{nis}_{\leq n}L_{\mathbb{A}^1})^m\esc{X}\simeq \tau_{\leq n}^{nis}\esc{X}$. Hence, $$\tau^{mot}_{\leq i}\esc{X}\simeq\dcolim_{m} (\tau^{nis}_{\leq n}L_{\mathbb{A}^1})^m\esc{X}\simeq \dcolim_{m}\tau^{nis}_{\leq n}\esc{X}\simeq \tau_{\leq i}^{nis}\esc{X}$$ as claimed.
\end{proof}

\begin{rem}[On the postnicov convergence of the motivic hhomotopy category]
 To the best of my knowledge, the Postnikov convergence of the motivic homotopy category is not yet known. However, it is known that the presentable subcategory of Nisnevich-connected motivic spaces $\tau^{nis}_{\geq 1}\mathcal{H}^{mot}$ is Postnikov convergent, at least over base fields. Indeed, the higher motivic truncations $\tau^{mot}_{\leq n}$ for objects in $\tau^{mot}_{\geq 1}\mathcal{H}^{nis}$ agree with their Nisnevich truncation $\tau_{\leq n}^{nis}$ in this case (see \Cref{layers of connected motivic space}). Since, by closedness under limits and Remark 2.9 of \textit{loc.cit.} (again), $\tau^{nis}_{\geq 1}\mathcal{H}^{mot}$ is Postnikov convergent, we conclude that so is $\tau^{mot}_{\geq 1}\mathcal{H}^{mot}$.
\end{rem}
We now come to one of the most important observations of this paper, namely that the motivic localization functor is path-injective (see \Cref{path injectivity}). As we have seen in the categorical section, \S2, this will have far-reaching consequences.

\begin{lem}\label[lem]{motivic 0 truncation is injective}
  Let $\tau_{\leq 0}^{{nis}}\mathcal{H}^{mot}$ denote the image of $\mathcal{H}^{mot}$ under $\tau_{\leq 0}^{nis}$ (we may call it the category of $\mathbb{A}^1$-homotopy path components). Then the restriction functor $$L_{\leq 0}^{mot}\simeq \mathcal{L}:\tau_{\leq 0}^{{nis}}\mathcal{H}^{mot}\to {Sh_{nis}^{\mathbb{A}^1}}$$ is injective. In other words, the motivic localization is path injective (\Cref{path injectivity}).
\end{lem}
\begin{proof}
    Using the isomorphism on $\mathcal{F}_k$ as in \Cref{pi0 A1 inv in one vsariable}, this follows immediately from [Corollary 2.10, \cite{bachmann2024strongly}].
\end{proof}
\begin{rem}
Here is a complete argument for the benefit of the reader: First, since $\pi_0\esc{X}(K)\simeq
    \pi_0^{nis}\esc{X(K)}$ for every $K\in \mathcal{F}_k$, we want to show that for any $K\in \mathcal{F}_k$ any two $t,t'\in \esc{X}_0(K)$ are connected by a simplicial homotopy in $\esc{X}(K)$. Since $\mathcal{L}\pi_0^{nis}\esc{X}(K)=*$, there is an $n$-ghost homotopy $$H=(\pi_0^{nis}\esc{X}\xleftarrow{h}V\xrightarrow[]{nis} \mathbb{A}^1_K,h^W:W\mathrel{\substack{\xrightarrow{ha} \\ \xrightarrow[hb]{}}} \pi_0^{nis}\esc{X}, \tilde{\sigma_i}),$$ where the last two maps $a$ and $b$ are the following composites (in which the first map is a Nisnevich cover): $$a,b: W\xrightarrow{nis}V\times_{\mathbb{A}^1_K}V\mathrel{\substack{\xrightarrow{pr_1} \\ \xrightarrow[pr_2]{}}} V.$$ Note that since $V$ is a $1$-dimensional scheme, it has homotopy dimension 1, and $\esc{X}$ has Nisnevich descent. By [\cite{bachmann2024strongly}, Lemma 2.17], we have a surjection $[V,\esc{X}]\to [V, \pi_0^{nis}\esc{X}]$. Therefore, up to homotopy, there is a map $u: V\to \esc{X}$ descending to $h$.
    

Now, we claim that for every $K\in \mathcal{F}_k$ and any two sections $ t,t'$, as above, are simplicially homotopic. We prove this by induction on $n$. As in [\cite{10.2140/akt.2022.7.385}, Theorem 2.2], we may assume that $V$ is an elementary Nisnevich cover $V=V_1\sqcup V_2$ and let $h_i=h_{|V_i}$ and $u_i=u_{|V_i}$, and denote $W'=V_1\times _{\mathbb{A}^1_K}V_2$. Let $W'=\cup W'_j$ be the irreducible components of $W'$, and let $\eta _j$ be the corresponding generic points. We deduce that $H'_{\eta_j}$ is an $(n-1) $ ghost homotopy of ${(u_1)}_{\eta_j}$ and $(u_2)_{\eta_j}$. By the induction hypothesis, this implies that $(u_1)_{\eta_j}$ and $(u_2)_{\eta_j} $ are simplicially homotopic in $\esc{X}$. By definition, these simplicial homotopies are defined over an open subset $W ''\subset  W '$. Consequently, $(u_1)_{|W''}$ and $(u_2)_{|W''}$ are simplicially homotopic in $\esc{X}$. Moreover, by further shrinking $V_2$, we may assume that actually $ W''= V_1 \times _{\mathbb{A}^1_K}V_2$. 

Therefore, we may assume that $u_1:V_1\to \esc{X}$ and $u_2: V_2\to \esc{X}$ agree on $V_1\times_{\mathbb{A}^1_k}V_2$ up to a simplicial homotopy. Since $V=V_1\sqcup V_2\to \mathbb{A}^1_K$ is a Nisnevich cover and $\esc{X}$ is a Nisnevich $\infty$-sheaf, we may therefore find a lift of $u: V\to \esc{X}$ to $u': \mathbb{A}^1_K\to \esc{X}$. But since $\esc{X}$ is $\mathbb{A}^1$-local, this yields a simplicial homotopy between the endpoints, namely $t$ and $t'$. 

The base case is easy. Indeed, a $0$-ghost homotopy from $t$ to $t'$ in $\pi_0^{nis}\esc{X}$ is given by an explicit map $\mathbb{A}^1_K\to \pi_0^{nis}\esc{X}$, which, by appeal to the same homotopy dimension argument [\cite{bachmann2024strongly}, Lemma 2.17], comes from $[\mathbb{A}^1_K, \esc{X}]$ and, by the $\mathbb{A}^1$-locality of $\esc{X}$, is nothing but a simplicial homotopy.

\end{rem}
\subsection{Motivic Eilenberg Maclane spaces}
Let $L_{\mathbb{A}^1}Grp_k$ and $L_{\mathbb{A}^1}Ab_k$ denote the full subcategories of $Sh^{nis}_k$ consisting of $\mathbb{A}^1$-invariant sheaves of groups and abelian groups, respectively. Both are full subcategories of $Sh^{\mathbb{A}^1}_k$. As introduced by Morel in his textbook [\cite{morel2012a1}], $Grp_k^{\mathbb{A}^1}$ denotes the full subcategory of $L_{\mathbb{A}^1}Grp_k$ consisting of strongly $\mathbb{A}^1$-invariant sheaves of groups, while $Ab_k^{\mathbb{A}^1}$ is the full subcategory consisting of strictly $\mathbb{A}^1$-invariant sheaves of abelian groups. In the rest of section \S3, we will require the base to be a field $k$ that we don't assume to be perfect unless otherwise specified.

As we did in the previous subsection, we shall restate the definitions of \S3.2 for the motivic localization. 

\begin{defn}\label[defn]{motivically n connective}
A motivic space $\esc{X}$ is motivically $1$-connective (equivalently, $0$-connected) if $\pi_0^{mot}\esc{X}$ is trivial. We say that a pointed motivic space $x: *\to \esc{X}$ is motivically $(n+1)$-connective (or $n$-connected) for $n> 0$ if, in addition, $\pi_i^{mot}(\esc{X},x )\simeq *$ for all $i\leq n$. We denote the full subcategory of $n$-connective objects of $\mathcal{H}^{mot}(k)$ by $\tau_{\geq n}^{mot}\mathcal{H}^{mot}(k)$.
\end{defn} 
\begin{thm}\label{nis connective vs mot connective of motivic spaces}
    For every $n\geq 1$ and $k$ a field (assumed perfect when $n\geq 2$), $$\tau_{\geq n}^{mot}\mathcal{H}^{mot}(k)\simeq \tau_{\geq n}^{nis}\mathcal{H}^{mot}(k).$$ 
\end{thm} 
\begin{proof}
    When $n=1$, this is \Cref{motivic 0 truncation is injective}. With this, the $n \geq 2$, the cases follow immediately from \Cref{pi_i mot=nis}.
\end{proof}

\begin{defn}
 Let $n\geq 1$. A Nisnevich sheaf of $n$-groups (i.e., of groups when $n=1$ and of abelian groups when $n\geq 2$) $G$ is said to be $n$-strictly $\mathbb{A}^1$-invariant if $\mathrm{B}^n_{nis}G$ is $\mathbb{A}^1$-local. We denote the full subcategory of $n-Grp^{nis}_k$ consisting of $n$-strictly $\mathbb{A}^1$-invariant sheaves by $n\mbox{-}Grp^{L_{mot}}$.
\end{defn}  
 \begin{rem}\label[rem]{strict local is local}
  Since $\Omega^n$ preserves local objects, it follows from the equivalence $G\simeq \Omega ^n\mathrm{B}_{nis}^nG$ that $n$-strictly $\mathbb{A}^1$-invariant sheaves are $\mathbb{A}^1$-invariant.
 \end{rem}
\begin{rem}
This is equivalent to saying that for all $i\leq n$, the presheaves $H^i_{nis}(-,G)$ are $\mathbb{A}^1$-local.
\end{rem}
\begin{rem}
    Clearly, $1$-strictly $\mathbb{A}^1$-invariant sheaves are precisely the strongly $\mathbb{A}^1$-invariant sheaves in the sense of [\cite{morel2012a1}, Definition 7 (2)]. In other words, $1\mbox{-}Grp^{L_{mot}}=Grp^{\mathbb{A}^1}_k$. Whereas $G$ is strictly $\mathbb{A}^1$-invariant in the sense of [\cite{morel2012a1}, Definition 7 (3)] if and only if it is $n$-strictly $\mathbb{A}^1$-invariant for all $n\geq 1$. 
\end{rem}
\begin{lem}\label[lem]{n-strict groups}
 Let $k$ be a perfect field. Then, for all $ n\geq 2$ $$n\mbox{-}Grp_{L_{mot}}=Ab^{\mathbb{A}^1}_k.$$
\end{lem} 
\begin{proof}
    First, observe that for every $n\geq 2$, an object of $n\mbox{-}Grp^{L_{mot}}(\mathcal{H}^{mot}(k))$ is, in particular, a strongly $\mathbb{A}^1$-invariant sheaf of abelian groups. By [\cite{morel2012a1}, Theorem 4.46], these are strictly $\mathbb{A}^1$-invariant when $k$ is perfect.
\end{proof}  

\begin{defn}
 For $n\geq 0$, a motivic space $\esc{X}$ is an $n$-Eilenberg-MacLane space if it is motivically $n$-truncated (see \Cref{motivically n truncated space}) and motivically $n$-connective. 
\end{defn} 
 \begin{lem}\label[lem]{motivic em vs nis em}
  For every $n\geq 0$ and every field $k$, there is an equality  \begin{flalign}
\mathcal{EM}_n(\mathcal{H}^{mot}(S))=L_{\mathbb{A}^1}\mathcal{P}(S)\bigcap \mathcal{EM}_{n}(\mathcal{P}^{nis}).\end{flalign}
 \end{lem}


\begin{lem}
    Let $k$ be a perfect field. Then there is an adjunction:
    $$\mathrm{B}^{}_{\mathcal{H}^{mot}(k)}:Grp (\mathcal{H}^{mot}(k))\leftrightarrows \mathcal{EM}_{1}(\mathcal{H}^{mot}): \pi_1^{mot}$$ and for every $n\geq 2$ an adjunction:
    $$\mathrm{B}^{n}_{\mathcal{H}^{mot}(k)}:Ab(\mathcal{H}^{mot}(k))\leftrightarrows \mathcal{EM}_{n}(\mathcal{H}^{mot}): \pi_n^{mot}$$ 
\end{lem}
\begin{proof}
    This follows from the facts that $\mathrm{B}_{\mathcal{H}^{mot}(k)}\simeq L_{mot}\mathrm{B}$ and that, for every $n\geq 1$, \Cref{pi_i mot=nis} gives $\pi_n^{mot}\simeq {\pi_i^{nis}}_{|\mathcal{H}^{mot}(k)}$.
\end{proof}
\begin{thm}\label{motivic em}
Let $k$ be a perfect field. Then the adjunction in the above lemma 
induces equivalences of categories: $$Sh_{*,k}^{\mathbb{A}^1}\simeq \mathcal{EM}_0(\mathcal{H}^{mot}(k)) \text{ and  }Grp_k^{\mathbb{A}^1}\simeq \mathcal{EM}_1(\mathcal{H}^{mot}(k)): \pi_1^{mot}$$ and for all $n\geq 2$ $$Ab_k^{\mathbb{A}^1}\simeq \mathcal{EM}_n(\mathcal{H}^{mot}(k)):\pi_n^{mot}.$$
\end{thm}
\begin{proof}
Since $L_{mot}$ is path-injective, this follows from \Cref{grp em} and \Cref{n-strict groups}.
\end{proof}
 Such characterizations are largely unknown over arbitrary base schemes. Nevertheless, over an arbitrary base, we have the equivalence $$\mathcal{EM}_0(\mathcal{H}^{mot})\simeq Sh_k^{\mathbb{A}^1}=\text{Disc}(\mathcal{H}^{mot})$$, which is expected of an $\infty$-topos. Unfortunately, 
 \begin{cor}[\cite{spitzweck2012motivic}, Remark 3.5]\label{mot not topos}
    Over any perfect field $k$, the motivic homotopy category $\mathcal{H}^{mot}(k)$ is not an $\infty$-topos. 
    \end{cor} 
    \begin{proof}
     First note that the inclusion $Grp^{\mathbb{A}^1}_k\subset L_{\mathbb{A}^1}Grp_k$ is not essentially surjective; indeed the reduced subgroup $\mathbb{Z}(\mathbb{G}m,1)$ is $\mathbb{A}^1$-invariant [\cite{choudhury2014connectivity}, see the paragraph after Lemma 5.4] but not strongly $\mathbb{A}^1$-invariant, see [Lemma 5.6]. At this point one may directly apply \Cref{not a topos localization}, but we shall sketch the argument again in the present (motivic) context. 

     If $\mathcal{H}^{mot}(k)$ were to be an $\infty$-topos, one would have equivalences of categories given by the functor $$\pi_1^{mot}\simeq \pi_1^{nis}:\mathcal{EM}_1(\mathcal{H}^{mot})\to  \text{Grp}(\text{Disc}(\mathcal{H}^{mot})) $$ [\cite{lurie2009higher}, Proposition 7.2.2.12]. However we see that by the above theorem, the essential image of the same is the full subcategory $Grp_k^{\mathbb{A}^1}$ of $\text{Grp}(\text{Disc}(\mathcal{H}^{mot})) = L_{\mathbb{A}^1}Grp_k^{nis}$.  But
     \end{proof}

\begin{rem}
Since $\mathbb{Z}(\mathbb{G}_m,1)$ is a sheaf of abelian groups, the proof above applies to all $n\geq 1$. In particular, there are $\mathbb{A}^1$-invariant sheaves of groups that are not strongly $\mathbb{A}^1$-invariant and hence not strictly $\mathbb{A}^1$-invariant.  Hence, for any $n\geq 1$ $$
\pi_n^{mot}\simeq \pi_n^{nis}:\mathcal{EM}_n(\mathcal{H}^{mot})\to  n{\mbox{-}}\text{Grp}(\text{Disc}(\mathcal{H}^{mot})) $$ is never an equivalence. This failure of the motivic homotopy category to be a topos was first brought to notice by [\cite{spitzweck2012motivic}, Remark 3.5]. In their argument, however, they did not explain what all the internal Eilenberg-MacLane objects of the motivic homotopy category are, making their argument somewhat incomplete.
\end{rem} 

    \begin{rem}[On the postnicov convergence of the motivic connected motivic homotopy category]\label{post comvg of conn}
 The presentable subcategory of categorically connected motivic spaces $\tau^{mot}_{\geq 1}\mathcal{H}^{mot}$ is Postnikov convergent, at least over base fields. Indeed, $$\tau^{mot}_{\geq 1}\mathcal{H}^{mot}=\tau^{nis}_{\geq 1}\mathcal{H}^{mot},$$ (\Cref{nis connective vs mot connective of motivic spaces}), and $\tau^{nis}_{\geq 1}\mathcal{H}^{mot}$ is Postnikov convergent by \Cref{post comvg of conn}. Thus, $\tau^{mot}_{\geq 1}\mathcal{H}^{mot}$ is Postnikov convergent.
\end{rem}
 \begin{cor}\label[cor]{internal layers of connected motivic space}
Let $\esc{X}\in  \mathcal{P}_{nis}(k)_{\geq 1}$ be a pointed object. Then the following conditions are equivalent:
\begin{enumerate}
    \item $\esc{X}\in \mathcal{H}^{mot}(k)$.
    \item $\mathrm{K}_{nis}(\pi_i^{nis}\esc{X},i)\in \mathcal{EM}_i(\mathcal{H}^{mot}(k))$ for all $i$.
    \item $\tau_{\leq i}^{nis}\esc{X}\in \mathcal{H}^{mot}(k)_{\leq i}$ for all $i$.
\end{enumerate}
 \end{cor} 
 \begin{proof}
    By \Cref{layers of connected motivic space}, this follows from identifying the conditions.

    For (2), use the identification in \Cref{motivic em vs nis em}. To identify the condition (3) here with the condition (3) there, use \Cref{motivic trunc is nis trunc}.
 \end{proof}
\subsection{Motivic connective objects and grouplike objects}
In the previous section, we have seen that the motivic homotopy category fails to internally deloop arbitrary discrete motivic groups. In fact, we have seen that those that can be delooped are precisely those that arise as the internal homotopy groups of (internal) motivic Eilenberg-MacLane objects. This section is devoted to understanding the condition required for the internal delooping of arbitrary motivic grouplike monoids and thus to establishing a motivic version of the loop-deloop equivalence between the $\infty$-categories of internal motivic connected objects and of those special grouplike objects.

The failure of the equivalence of the functor $\pi_1^{mot}:\mathcal{EM}_1(\mathcal{H}^{mot}(k))\to L_{\mathbb{A}^1}Grp_k^{nis}$, as explained in \Cref{mot not topos}, extends to the deeper fact that the functor $$\Omega: \mathcal{H}^{mot}(S)_{\geq 1}\to  Grp(\mathcal{H}^{mot}(S))$$ is not an equivalence of infinity categories. The primary issue is that, as shown by the same example of Choudhury ([\cite{choudhury2014connectivity}], which we have cited multiple times), this functor is not essentially surjective. It is easy to see that, over perfect fields, the spaces in the image must have strongly $\mathbb{A}^1$-invariant motivic $\pi_0$, since these are just motivic $\pi_1^{\mathbb{A}^1}$'s of the starting space [\cite{morel2012a1}, Corollary 6.3]. In fact, this condition suffices to ensure a motivic delooping. Before we explain this, let us first restate the main definition and result of \S2.3 applied to the motivic case:
\begin{defn}
We say that a motivic monoid $\esc{M}$ is strongly $L_{mot}$-local if $\mathrm{B}_{mot}\esc{M}\simeq \mathrm{B}_{nis}\esc{M}$, i.e., if $\mathrm{B}_{nis}\esc{M}$ is motivic local. We denote the full subcategory of strongly $L_{mot}$-local spaces by $\mathrm{Mon}_{L_{mot}}(\mathcal{H}^{mot}(k))$. We denote the further full subcategory of grouplike objects by $\mathrm{Grp}_{L_{mot}}(\mathcal{H}^{mot}(k))$.
\end{defn}
\begin{cor}\label[cor]{strong motivic groups and motivic connected spaces}
        There is an equivalence of $\infty$-categories:
        $$\mathrm{B}_{\mathcal{H}^{mot}(k)}: \mathrm{Grp}_{L_{mot}}(\mathcal{H}^{mot}(k))\simeq \mathcal{H}^{mot}(k)_{\geq 1}:\Omega.$$
\end{cor}
\begin{proof}
         Due to \Cref{motivic 0 truncation is injective}, the categorical theorem \Cref{loop bar for localization} applies to the motivic localization $L_{mot}$ of the Nisnevich topos.
\end{proof}
    
    We first make a preliminary observation about the motivic counterparts of notons from \S2.3:
\begin{prop}\label[prop]{strong mot is strong a1}
    A monoid motivic space is strongly $L_{mot}$-local iff it is strongly $\mathbb{A}^1$-invariant in the sense of \textup{[\cite{elmanto2021motivic}, Definition 3.1.6(1)]}. If $k$ is a perfect field, then this is iff $\pi_0^{nis}\esc{M}$ is strongly $\mathbb{A}^1$-invariant.
\end{prop} \begin{proof}
    The first line follows from comparing the definitions of the two terms. When $k$ is a perfect field, the second line is a restatement of [\cite{elmanto2021motivic}, Proposition3.1.12(1)] under the first line. 
\end{proof}
\begin{lem}\label[lem]{pi0 of monoid}
    Let $\esc{G}$ be a grouplike motivic monoid space. Then $\pi_0^{mot}\esc{G}\simeq \pi_0^{nis}\esc{G}$.
\end{lem} 
\begin{proof}
    Since grouplike motivic monoid spaces are, in particular, motivic $H$-groups, this is a special case of \Cref{pi0 of h space}.
\end{proof}
  \begin{cor}\label[cor]{strongly Lmot iff pi0 strongly Lmot}
Over a perfect field, a grouplike monoid motivic space $\esc{G}$ is strongly $L_{mot}$-local if and only if $\pi_0^{mot}\esc{G}$ is strongly $L_{mot}$-local.      
    \end{cor} 
\begin{proof}
    This follows immediately from \Cref{strong mot is strong a1} and \Cref{pi0 of monoid}.
\end{proof}
Let us denote by $\mathcal{G}rp_{mot}(\mathcal{H}^{mot}(k))$ the full subcategory of $\mathrm{Grp}(\mathcal{H}^{mot}(k))$ consisting of motivic monoid spaces whose motivic path component $\pi_0^{mot}$ is strongly $L_{mot}$-local.
\begin{cor}\label[cor]{mot bar omega}
Over a perfect field, there is an equivalence of categories $$ \mathrm{B}^{nis}:  \mathcal{G}rp_{mot}(\mathcal{H}^{mot}(k))\simeq\mathcal{H}_\bullet^{mot}(k)_{\geq 1}:\Omega$$  where $ Grp_{mot}(\mathcal{H}^{mot}(S))$ denotes the full subcategory of $\mathcal{G}rp(\mathcal{H}^{mot}(k))$ consisting of motivic monoids with strongly $L_{mot}$-local $\pi_0^{mot}.$ This provides further evidence that $\mathcal{H}^{mot}$ is not an $\infty$-topos.
\end{cor}  
\begin{proof}
By \Cref{strong motivic groups and motivic connected spaces}, the claim follows from \Cref{strongly Lmot iff pi0 strongly Lmot}.
    
To justify the last line, it suffices to know that there exist grouplike motivic spaces whose motivic path components $\pi_0^{mot}$ ( $\simeq \pi_0^{\mathbb{A}^1}$) are not strongly $\mathbb{A}^1$-invariant. This is the content of [\cite{choudhury2014connectivity}, the paragraph after Lemma 5.4], which we have used earlier in \Cref{mot not topos}.\end{proof}
\subsection{Motivic theory of covering spaces}

\textbf{Motivic theory of geometric coverings}
\begin{defn}
A collection $\{\esc{U}_i\to \esc{X}\}$ of morphisms of Nisnevich local spaces is said to be a motivic Geometric covering if $L_{mot}\check{C}_\bullet p\to L_{mot}\esc{X}$ is a colimit diagram in $\mathcal{H}^{mot}(k)$, where $p: \bigsqcup\esc{U}_i\to \esc{X}$ is the disjoint union $\sqcup p_i$.
\end{defn}
\begin{ex}
    Every Nisnevich cover of a (even, not necessarily smooth) scheme is a Motivic geometric cover.
\end{ex}
\begin{cor}\label[cor]{a1 geom conn}
A $k$-space $\esc{X}$ is $\mathbb{A}^1$-connected if and only if there exists an $\mathbb{A}^1$-geometric covering (which is if and only if a Nisnevich-geometric covering) $\esc{U}_i\to \esc{X}$ such that each $\esc{U}_{i}$ is $\mathbb{A}^1$-connected and each $\esc{U}_{ij}$ admits a global point. Then $\esc{X}$ is $\mathbb{A}^1$-connected.
\end{cor}
\begin{proof}
    If $\esc{X}$ is $\mathbb{A}^1$-connected then the identity $\esc{X}\to \esc{X}$ is a Nisnevich cover with this property, which in turn is an $\mathbb{A}^1$-grometric covering. Note that the existence of the rational point from connectivity follows from the effectivity of $L_{mot}$. The remaining implication follows from \Cref{geom L connectivity} in the presence of \Cref{motivic 0 truncation is injective} and the fact that $L_{mot}$ is effective, so that for any space $\esc{Y}$, existence of a global point of $\pi_0^{\mathbb{A}^1}\esc{Y}$ is equivalent to $\esc{Y}$ being pointed. 
\end{proof}
\begin{cor}\label[cor]{local-global-connectivity}
  A scheme $X/k$ is $\mathbb{A}^1$-connected iff there is a nisnevich covering $U_i\to X$ such that each $U_i$ is $\mathbb{A}^1$-connected and for every pair $i,j$ the fiber product $U_i\times_XU_j$ has a ($k$-)point.
\end{cor}

\begin{cor}[Motivic Van Kampen]
Let $\esc{U}_i\to \esc{X}$ be an $\mathbb{A}^1$-geometric covering (for example, a Nisnevich local covering) such that $\esc{U}_{ijk}$ is $\mathbb{A}^1$-connected for every $i,j,k$ (1-, 2-, and 3-fold). Then $\esc{X}$ is $\mathbb{A}^1$-connected, and the $\mathbb{A}^1$ homotopy group of $\esc{X}$ is given by the following coequalizer in $Grp^{\mathbb{A}^1}_k$:
$$\underset{i,j}{*}\pi_1^{\mathbb{A}^1}({\esc{U}_i\times_\esc{X}\esc{U}_j})\rightrightarrows \underset{i}{*}\pi_1^{\mathbb{A}^1}({\esc{U}_i})\to \pi_1^{\mathbb{A}^1}(\esc{X})\to *$$
\end{cor}
\begin{proof}
    This follows at once from \Cref{L van campen} and \Cref{motivic 0 truncation is injective}, using the identification $\pi_1^{\mathbb{A}^1}=\pi_1^{mot}$ in \Cref{pi_i mot=nis}. 
\end{proof}
\begin{cor}[Generalized Morel's $\mathbb{A}^1$ Van Kampen, [\cite{morel2012a1}, Theorem 7.12\text{]}]
Let $\{{U}_i\to {X}\}_{i\in I}$ be a Nisnevich covering of a smooth scheme $X$ such that each ${U}_{ijk}$ is pointed $\mathbb{A}^1$-connected for every $i,j,k\in \{I\}\cup I$ (where $U_I$ is 1-, 2-, or 3-fold). Then ${X}$ is $\mathbb{A}^1$-connected, and the $\mathbb{A}^1$ homotopy group of ${X}$ is given by the following coequalizer in $Grp^{\mathbb{A}^1}_k$:
$$\underset{i,j}{*}\pi_1^{\mathbb{A}^1}({{U}_i\times_{X}{U}_j})\rightrightarrows \underset{i}{*}\pi_1^{\mathbb{A}^1}({{U}_i})\to \pi_1^{\mathbb{A}^1}({X})\to *$$
\end{cor}
\begin{rem}
While Morel establishes an $\mathbb{A}^1$-local Van Kampen theorem specifically for Zariski covers over perfect fields [\cite{morel2012a1}, Theorem 7.12], the preceding result applies to more general covers. Furthermore, our statement includes a connectivity assumption on the triple intersections to resolve a technical gap in \textit{loc. cit.} The original argument relies on the exactness of the following diagram:$$Hom_{\mathcal{G}r}(\pi_1^{\mathbb{A}^1}(\check{C}(\mathcal{U})), G) \to \prod_i Hom_{\mathcal{G}r}(\pi_1^{\mathbb{A}^1}(U_i), G) \rightrightarrows \prod_{i,j} Hom_{\mathcal{G}r}(\pi_1^{\mathbb{A}^1}(U_i \cap U_j), G).$$ Technically, this diagram represents the $\pi_0$ of the corresponding limit diagram of spaces:$$\prod_i Map(U_i,\mathrm{B} G) \rightrightarrows \prod_{i,j} Map(U_i \cap U_j,\mathrm{B} G)\cdots$$Since $\pi_0$ may fail to commute with limits, the required exactness is not automatic. Resolving the resulting cocycle condition on the double intersections inherently requires the connectivity of the triple intersections.
\end{rem} 
\textbf{Motivic theory of spatical coverings}

\begin{defn}[\Cref{defn for covering}, \Cref{definition for L covering}]
\begin{enumerate}
\item A morphism $f:\esc{Y}\to \esc{X}$ of motivic spaces is a motivic covering if $f\in \big({\mathcal{H}^{mot}_{/\esc{X}}}\big)_{\leq 0}$. Denote the category of such coverings by $Cov_{\mathcal{H}^{mot}}(\esc{X})$. 
\end{enumerate}
When $f:\esc{Y}\to \esc{X}$ is a morphism of Nisnevich local spaces, then:
\begin{enumerate}[resume] 
\item We say that it is an $\mathbb{A}^1$-fibration if it is a local object of ${\mathcal{P}_{nis}}_{/\esc{X}}$ with respect to motivic equivalences contained in ${\mathcal{P}_{nis}}_{/\esc{X}}$. We denote the category of $\mathbb{A}^1$-fibrations of $\esc{X}$ by $\mathcal{F}ib_{\esc{X}}^{\mathbb{A}^1}.$
\item We call it an $\mathbb{A}^1$-covering if it is an $\mathbb{A}^1$-fibration and is a $0$-truncated object in the $\infty$-category of $\mathbb{A}^1$-fibrations over $\esc{X}$. Denote the category of $\mathbb{A}^1$-coverings of $\esc{X}$ as $Cov_{\mathbb{A}^1}(\esc{X})$. 
\end{enumerate}
\end{defn}
\begin{prop}
A morphism $\esc{X}\to \esc{Y}$ of motivic spaces is a motivic covering if and only if it is an $\mathbb{A}^1$ covering if and only if it is a Nisnevich covering.
\end{prop}
\begin{proof}
Follows immediately from \Cref{Covering of Localization} and \Cref{L covering of L local object}.
\end{proof}
\begin{rem}
Because the motivic localization functor $L_{mot}$ is effective (see \Cref{Effective localization remark}), we obtain, by \cref{covering pullback description}, a pullback square of $1$-categories:
\[
\xymatrix{
\mathcal{C}ov_{\mathbb{A}^1}(\esc{X})\ar@{^(->}[r]\ar[d]&\mathcal{H}^{mot}_{/L_{mot}\esc{X}}\ar[d]\\
\mathcal{C}ov_{nis}(\esc{X})\ar@{^(->}[r] &{\mathcal{P}_{nis}(S)}_{/X}
}
\]
\end{rem}
\begin{cor}
Let $\esc{Z},\esc{X}$ be Nisnevich-local spaces, with $\esc{X}$ pointed and $\mathbb{A}^1$-connected. Then a morphism $f:\esc{Z}\to L_{mot}\esc{X}$ is an $\mathbb{A}^1$-covering iff $fib(f)\in Sh^{\mathbb{A}^1}$.
\end{cor}
\begin{proof}
Follows immediately from \Cref{L covering iff fiber is L discrete} and from the fact that $L_{mot}$ is closed under extensions over $\mathbb{A}^1$-connected objects [\cite{morel2012a1}, Lemma 6.51].
\end{proof}
\begin{cor}
Suppose $\esc{X}$ is a pointed, connected motivic space. Then $Cov_{\mathcal{H}^{mot}}(\esc{X})_*$ has an initial object, namely the ordinary simply connected cover. 
\end{cor}
\begin{proof}
    By [\cite{morel2012a1}, Corollary 6.2], $\Pi^{\mathbb{A}^1}\esc{X}$ is a motivic space. Whence, the claim follows from \Cref{Universal cover for fundamental localization}. 
\end{proof}

\begin{cor}\label[cor]{Lmot covering is covering of Lmot}
  Suppose $\esc{X}$ is a Nisnevich sheaf of spaces. We then have an equivalence of categories $$\eta^*_{mot}:Cov_{{\mathbb{A}^1}}({L_{mot}}\esc{X)}=Cov_{\mathcal{H}^{mot}}({L_{mot}}\esc{X)}\leftrightarrows  Cov_{{\mathbb{A}^1}}(\esc{X)}: L_{mot}$$ where $\eta_{mot}$ is the localization unit $1\to L_{mot}$. The category on the left-hand side is the category of covering spaces of $L_{mot}\esc{X}$ internal to the motivic homotopy category, while the right-hand side is the $\infty$-categorical analogue of Morel's category of $\mathbb{A}^1$-coverings of $\esc{X}$, as developed in \textup{[\cite{morel2012a1}, \S7]}.
\end{cor}
\begin{proof}
    Since $L_{mot}$ is locally cartesian [\cite{hoyois2017six}, Proposition 3.15] and effective (\Cref{Effective localization remark}), the result follows from \Cref{L covering is covering of L}.
\end{proof}
\begin{thm}\label{motivically simply connected covering space}
Let $\esc{X}$ be a pointed and motivically connected $k$-motivic space over a perfect field $k$.
\begin{enumerate}
    \item Then the category $\mathcal{C}ov_{\mathcal{H}^{mot}}(\esc{X})_*$ has an initial object ${\widetilde{{\esc{X}}}}^{^{mot}}\to \esc{X}$.
    \item This initial object is the unique one with a motivically simply connected source, i.e., it satisfies $\tau_{\leq 1}^{mot}{\widetilde{{\esc{X}}}}^{^{mot}}\simeq *$. The inclusion  $$\mathcal{C}ov_{\mathcal{H}^{mot}(k)}(\esc{X})_*\subset \mathcal{C}ov_{\mathcal{P}_{nis}(k)}(\esc{X})_*$$ preserves this initial object, i.e., this initial object is the unique Nisnevich covering with a Nisnevich simply connected source.
    \item The canonical map $ \pi_1^{mot}(\esc{X})\to \underline{Aut}_{\esc{X}}(\widetilde{{\esc{X}}}^{^{mot}})$ is an equivalence.
\end{enumerate}
\end{thm}
\begin{proof}
    Since motivically connected spaces are $\mathbb{A}^1$-connected by \Cref{motivic 0 truncation is injective}, we know that $\esc{X}$ is $\mathbb{A}^1$-connected. By [\cite{morel2012a1}, Corollary 6.2], it follows that $\Pi^{\mathbb{A}^1}\esc{X}$ is a motivic space. Whence \Cref{Universal cover for fundamental localization} applies, and statement (1) follows. 
    
Statement (2) is a formal consequence of the fact that, from the referred proposition, one has $(\widetilde{{\esc{X}}}^{^{mot}})\simeq \tau_{>1}^{nis}\esc{X}$, which is the unique simply connected (simplicial) cover of $\esc{X}$. Since this is motivic local and $\mathbb{A}^1$ connected, we have $$\tau_{\leq 1}^{mot}{\widetilde{{\esc{X}}}}^{^{mot}}\simeq \tau_{\leq 1}^{nis}{\widetilde{{\esc{X}}}}^{^{mot}},$$ which is trivial, as obtained in the previous line. 
    
The third statement follows from \Cref{aut=pi_1}.
\end{proof}

\begin{thm}[compare the methods with [\cite{morel2012a1}, Theorem 7.8\text{]}]\label{a1 simply connected covering space}
Let $\esc{X}$ be a pointed, $\mathbb{A}^1$-connected, Nisnevich-local space. 
\begin{enumerate}
    \item Then the category $\mathcal{C}ov_{\mathbb{A}^1}(\esc{X})_*$ has an initial object $\widetilde{{\esc{X}}}^{^{\mathbb{A}^1}}$. 
    \item The initial object is the unique $\mathbb{A}^1$-covering with $\mathbb{A}^1$ simply connected source.
    \item   The canonical map $ \pi_1^{\mathbb{A}^1}(\esc{X})\to \underline{Aut}_{\esc{X}}(\widetilde{{\esc{X}}}^{^{\mathbb{A}^1}})$ is an equivalence.
\end{enumerate}
\end{thm}
\begin{proof}
When $\esc{X}$ is pointed $\mathbb{A}^1$-connected, $L_{mot}\esc{X}$ is pointed and Nisnevich locally connected. Therefore, by the above corollary, \Cref{Lmot covering is covering of Lmot}, we may replace $\esc{X}$ with $L_{mot}\esc{X}$ and assume that $\esc{X}$ is $\mathbb{A}^1$-local. In that case, \Cref{motivically simply connected covering space} applies. For the third statement to follow from the third statement of the previous theorem, it suffices to know that $\pi_1^{\mathbb{A}^1}\esc{X}\simeq \pi_1^{mot}\esc{X}$ (see \Cref{pi_i mot=nis}).
\end{proof}
However, it might not be true that the ordinary simplicial covering of a space maps to the $\mathbb{A}^1$-simply connected covering under $L_{mot}$. This works for nisnevich-locally connected spaces when $\tau^{nis}_{\leq 1}\esc{X}$ is $\mathbb{A}^1$-local, which happens precisely when $\pi_1^{nis}\esc{X}$ is strongly $\mathbb{A}^1$-invariant.
\begin{prop}
Suppose $\esc{X}$ is a Nisnevich-local, pointed, connected space with a strongly $\mathbb{A}^1$-invariant fundamental group. Then $L_{mot}$ takes the universal Nisnevich-local cover of $\esc{X}$ to the universal $\mathbb{A}^1$-simply connected cover.
\end{prop}
\begin{proof}
    This follows immediately from [\cite{morel2012a1}, Theorem 6.59] for $n=1$ and from the identification of the $\mathbb{A}^1$-simply connected covering with $(\widetilde{{\esc{X}}}^{^{\mathbb{A}^1}})\simeq \tau_{>1}^{nis}L_{mot}\esc{X}$.
\end{proof}
\begin{cor}[\cite{morel2012a1}, Lemma 7.5(1)]
Let $G$ be a strongly $\mathbb{A}^1$-invariant Nisnevich sheaf of groups. Then any $G$-torsor is an $\mathbb{A}^1$-covering. 
\end{cor}
\begin{proof}
    This follows at once from \Cref{torsors are covering} and the definition of strong $\mathbb{A}^1$ invariance.
\end{proof}
 
\begin{prop}[This was stated in \text{[}\cite{morel2012a1}, Remark 7.10\text{]}]
Suppose, $\esc{Y}\in \mathcal{P}_{nis}(k)$ and $y:*\to \esc{Y}$ a base point. Then pull back along $y$ induces a functor $$y^*:Cov_{\mathbb{A}^1}(\esc{{Y}})\to Cov_{\mathbb{A}^1}(*)=Sh^{\mathbb{A}^1}_k.$$ This functor factors through the category of $\mathbb{A}^1$-invariant nisnevich sheaves (of sets) having an action by $\pi_1^{\mathbb{A}^1}(\esc{Y},y)$ (by the canonnical functor forgetting the action): $$\Gamma_y^{\mathbb{A}^1}:Cov_{\mathbb{A}^1}(\esc{{Y}})\longrightarrow \pi_1^{\mathbb{A}^1}(\esc{Y},y)\mbox{-}Sh^{\mathbb{A}^1}_k.$$
When $\esc{Y}$ is $\mathbb{A}^1$ connected, $\Gamma_y^{\mathbb{A}^1}$ is an equivalence of categories.
\end{prop}
\begin{proof}
The existence of the pullback functor follows from the pullback stability of ($\mathbb{A}^1$) coverings, \Cref{covering lemma} (3). That it factors through $\pi_1^{\mathbb{A}^1}(Y,y)\mbox{-}Sh^{\mathbb{A}^1}_k$ is a consequence of \Cref{action of the fundamerntal groupoid}, in view of the fact that when $\esc{Y}$ is $\mathbb{A}^1$ connected, there is an equivalence $$\Pi^{\mathbb{A}^1}\esc{Y}\simeq \Pi^{mot}L_{mot}\esc{Y},$$ or, in other words, that $\Pi^{\mathbb{A}^1}\esc{Y}$ is $\mathbb{A}^1$ local (see \Cref{layers of connected motivic space}).

Since $\mathcal{H}^{mot}$ is closed under extensions among connected objects (see [\cite{morel2012a1}, Lemma 6.51]), the claim follows immediately from \Cref{L galois theory}.
\end{proof}

\section{The internal homotopy theory of \texorpdfstring{$\mathcal{H}^{bir}$}{h}} 
The main observation of the first three subsections of the previous section is that the motivic homotopy category, constructed out of smooth schemes, even over a field, is not an $\infty$-topos, in contrast to the ordinary homotopy theory of smooth manifolds. To elaborate, let us have a quick recap of how these categories are born from universality:
\begin{prop}\label[prop]{robalo}
    Let $S$ be a Qcqs schem. Then the functor $$h_S^{mot}:=L_{mot}h_S:Sm_S\to  \mathcal{H}^{mot}(S)$$ 
    is universal among functors to a presentable $\infty$-category that contracts the affine line and sends the Nisnevich squares to pushout squares. In other words, the large co-presheaf $$Cat_{\infty}^{\mathbb{A}^1,nis}(Sm_S, -): Pr^L\to \mathcal{G}_1\mbox{-}Spc$$ sending a presentable $\infty$-category $\mathcal{D}$ to the full subcategory of $Cat_\infty(Sm_S,\mathcal{D})$ satisfying the two properties mentioned above, is representable by $h_S^{mot}: Sm_S\to \mathcal{H}^{mot}(S)$. 
\end{prop}
\begin{proof}
    See [\cite{robalo2012noncommutative}, Theorem 5.2].
\end{proof}
In the classical case of smooth manifolds, the answer is well known and justifies the construction of the motivic homotopy category:
\begin{prop}
    Let $\mathrm{SmMan}$ be the category of smooth manifolds. Then the co-presheaf $$Cat_\infty^{\mathbb{R}^1,open}(\mathrm{SmMan}, -): Pr^L\to \mathcal{G}_1\mbox{-}Spc$$ sending a presentable $\infty$-category $\mathcal{D}$ to the full subcategory of $Cat_\infty(SmMan, \mathcal{D})$ consisting of functors that contract the real line $\mathbb{R}^1$ and map twofold open cover squares to pushout squares is representable by the embedding $\mathrm{SmMan}\hookrightarrow Spc$.
\end{prop} 
\begin{proof}
    This is [\cite{dugger1999sheaves}, Remark 3.4.10].
\end{proof}  
Evidently, the resulting category in the classical case is not just a presentable $\infty$-category but, in fact, an $\infty$-topos. However, the schematic case is not as beautiful. Specifically, in \Cref{motivic em} we have seen that the motivic homotopy category need not be an $\infty$-topos. The problem is that, in the case of manifolds, one can always construct a `good' cover, in the sense that, 'open cover locally,' every manifold is homeomorphic to some $\mathbb{R}^n$; i.e., given a manifold $M$, one can always find an open cover $\{U_i\subset M\}$ such that each simplex of its \v{C}ech nerve is homeomorphic to a union of copies of $\mathbb{R}^{dimM}$. Even though one knows that (say over a field) etale locally every smooth scheme $X$ is isomorphic to $\mathbb{A}^{dimX}$, this does not pass to the corresponding \v{C}ech nerve, so that even etale locally schemes might have $\mathbb{A}^1$-locally non-zero homotopy dimension. Thus, even the etale local motivic homotopy category is not an $\infty$-topos. 

Instead of seeking a suitable topology in which such a dimension reduction (from pure homotopy dimension to $\mathbb{A}^n$'s) is possible, algebraic geometry offers an interesting way to brute-force this reduction. Namely, because fields have homotopy dimension zero, the desired $\infty$-topos might arise by identifying every variety with its function field, or equivalently, by identifying every open subscheme of a given scheme in $Sm_k$. This is precisely the birational motivic homotopy category of [\cite{0bat}] (where it is denoted $\mathcal{H}^{b\mathbb{A}^1}(S)$), [\cite{bachmann2019voevodsky}] (where it is denoted $L^0_{bir}\mathcal{H}^{\mathbb{A}^1}(S)$), and [\cite{pelaez2014unstable}], which we now recall below.

Given a Qcqs scheme $S$, we define the birational motivic homotopy category, denoted $\mathcal{H}^{bir}(S)$, as the localization of $\mathcal{H}^{mot}(S)$ at the set of all dense open immersions in $Sm_S$. The key observation of [\cite{0bat}] is what makes the contents of this section work and thus suggests that the birational motivic homotopy category is close to being an $\infty$-topos. Let us recall it here for the convenience of the reader:
\begin{thm}\label{Hbir is Hb}
    There is an equivalence of $\infty$-categories:
    $$\mathcal{H}^{bir}(S)=\mathcal{H}^b(S):=L_{dense}\mathcal{P}_\Sigma(S).$$
\end{thm} 
Here, $\mathcal{P}_\Sigma(S)$ is the full subcategory of $\mathcal{P}(S)$ consisting of presheaves that turn coproducts of schemes into products of spaces. Being (identically) the category of sheaves $\mathcal{P}_{\sqcup}(S)$ for the disjoint cover topology $\sqcup$, this is also an $\infty$-topos [\cite{bachmann2017norms}, Lemma 2.4], just like the Nisnevich one $\mathcal{P}_{nis}(S)$. But this new topos is quite well-behaved compared to the Nisnevich one. To see this, let $\tau^\sigma_{\leq n}$ denote the truncation functors for the various topologies we deal with. The ordinary truncation $\tau_{\leq n}$ of $\mathcal{P}(S)$ preserves products and hence takes $\mathcal{P}_\Sigma(S)$ to itself. Since $\mathcal{P}_\Sigma(S)=\mathcal{P}_\sqcup(S)$ is a left exact localization of $\mathcal{P}(S)$, the previous sentence, together with [\cite{lurie2009higher}, 5.5.6.28], implies that $\tau_{\leq n}^{\sqcup}\simeq {\tau_{\leq n}}_{|\mathcal{P}_{\Sigma}(S)}$. This implies, in particular, that $\pi_i^\sqcup\simeq {\pi_i}_{|\mathcal{P}_\Sigma(S)}$. (Such a statement is not true nisnevich locally, and the correct relation there is only $\tau^{nis}_{\leq n}\simeq L_{nis}\tau_{\leq n}$.)
\subsection{Internal Birational motivic truncations and Birational motivic homotopy group}

We now discuss the birational local analogue of \S2.1. We won't follow the lines of \S3.1 to restate the definitions for the birational category; we recommend that the reader look at \S2.1 or \S3.1. Let us denote the internal truncations and the internal homotopy groups of the presentable $\infty$-category $\mathcal{H}^{bir}(S)$ by $\tau_{\leq n}^{bir}$ and $\pi_i^{bir}$, respectively.

\begin{lem}\label[lem]{checking bir locality on pi}
 An $S$-space $\esc{X}$ is a birational motivic space if and only if, for every $i\geq 0$, $\pi_i\esc{X}$ is birational motivic. 
\end{lem}

\begin{proof}
By \Cref{Hbir is Hb}, this follows from the observation that ${\pi_i}_{|\mathcal{P}_\Sigma(S)}$ is automatically a $\Sigma$-presheaf (i.e., a $\sqcup$-sheaf).
\end{proof}

Therefore, we see that $\tau_{\leq n }$ maps $\mathcal{H}^{bir}(S)$ to itself:
\begin{thm}\label{birational truncation}
  There are identifications (of truncations restricted to the the birational homotopy category) $$\tau^{bir}_{\leq n}={\tau_{\leq n}}_{|\mathcal{H}^{bir}(S)}={\tau^{\sqcup}_{\leq n}}_{|\mathcal{H}^{bir}(S)}={\tau^{nis}_{\leq n}}_{|\mathcal{H}^{bir}(S)}$$ and of homotopy groups $$\pi_i^{bir}={\pi_i}_{|\mathcal{H}^{bir}(S)}.$$
\end{thm}
\begin{proof}
    This follows directly from the fact that $\tau_{\leq n}$ preserves birational local objects, say by [\Cref{checking bir locality on pi}]. Indeed, for $$\esc{X}\in\mathcal{H}^{bir}(S)\text{ and }\esc{Y}\in\mathcal{H}^{bir}(S)_{\leq n}=\mathcal{H}^{bir}(S)\bigcap \mathcal{P}_{}$$ one computes \begin{flalign*} Map_{{\mathcal{H}^{bir}(S) }_{\leq n}}(\tau^{b}_{\leq n}\esc{X}, \esc{Y}) \simeq Map_{{\mathcal{H}^{bir}(S) }}(\esc{X}, \esc{Y}) &=Map_{\mathcal{P}(S)
    }(\esc{X}, \esc{Y}) \\&\simeq  Map_{\mathcal{P}(S)_{\leq n} }({\tau_{\leq n}}\esc{X}, \esc{Y})\\&=Map_{\mathcal{H}^{bir}(S)_{\leq n} }({\tau_{\leq n}}\esc{X}, \esc{Y})\end{flalign*}
    (here the last equality follows from the fact that $\tau_{\leq n}\esc{X}\in \mathcal{P}(S)_{\leq n}$ is birational and hence belongs to $\mathcal{H}^{bir}_{\leq n}$ as well, while it was given that $\esc{Y}\in\mathcal{H}^{bir}(S)_{\leq n}$. The result then follows by the Yoneda lemma.
Moreover, because $L_{nis}$ of a birational local object is the object itself [\cite{0bat}, Theorem 2.1.1], we have $\tau^{bir}_{\leq n}={\tau^{nis}_{\leq n}}_{|\mathcal{H}^{bir}{(S)}}$.

It then follows from definition that $$\pi_i^{bir}:=\tau_{\leq 0}^{bir}\Omega^i=\tau_{\leq 0}\Omega^i={\pi_i}_{|\mathcal{H}^{bir}(S)}.$$
\end{proof}
\begin{thm}\label{path injectivity of Lbir}
    $L_{\leq 0}=\pi_0^{b}$, and the restriction $$\tau_{\leq 0}^{bir}: \tau_{\leq 0}\mathcal{H}^{bir}(S)\to \text{Disc}\mathcal{P}_{nis}(S)=Sh^{nis}_S$$ is simply the inclusion of birational sheaves.
\end{thm}
\begin{proof}
    This follows immediately from the observation that $\pi_0$ of birational motivic spaces is birational (\Cref{checking bir locality on pi}).
\end{proof}

\begin{thm}\label{post comp Hbir}
    The birational motivic homotopy category is postnikov complete.
\end{thm}
\begin{proof}
   By construction, $\mathcal{H}^{bir}(S)\subset \mathcal{P}(S)$ is closed under limits and, by \Cref{birational truncation}, under truncations. Since $\mathcal{P}(S)$ is Postnikov complete, so is $\mathcal{H}^{bir}(S)$. 
\end{proof}

\subsection{Birational Eilenberg Maclane spaces}
In Section \S3.2, we saw that the ordinary motivic homotopy category fails to be an $\infty$-topos. The argument stemmed from the observation that, in the ordinary motivic case, not every discrete motivic group admits an internal delooping. In this section, we shall show that this issue disappears after birational localization. The key algebraic reason is that birational sheaves are flasque. 


We denote by $Sh^b_S$ the category of sheaves of sets that are birational local. This is precisely the full subcategory of discrete/$0$-truncated objects of $\mathcal{H}^{bir}(S)$. Similarly, $Grp^b_S$ and $Ab^b_S$ denote the full subcategories of birational local sheaves of groups and abelian groups, respectively. Regarding \Cref{strictly n local}, we say that,
\begin{defn}
    A presheaf of (abelian) groups $G$ is ($n$-strictly) strongly $L_{bir}$-local if ($\mathrm{B}^n_{nis}G$) $\mathrm{B}_{nis}G$ is $L_{bir}$-local. We denote the full subcategory of $n$-strictly $L_{bir}$-local groups by $n\mbox{-}Grp(Disc(\mathcal{H}^{bir}(S)))$.
\end{defn}
\begin{lem}
A presheaf of groups $G$ is $L_{bir}$-local if and only if it is strongly $L_{bir}$-local. In other words, there is an equivalence $Grp^b_S=Grp_{L_{bir}}(Disc(\mathcal{H}^{bir}(S)))$.
\end{lem} 
\begin{proof}
 This is [\cite{0bat}, Corollary 3.2.9].
\end{proof} 
Similarly, 
\begin{lem}\label[lem]{n-strict bir groups}
    Let $n\geq 2$. A presheaf of groups $G$ is $n$-strictly $L_{bir}$-local if and only if $G$ is $L_{bir}$-local.
\end{lem} 
\begin{defn}
    A pointed birational local space $x:*\to \esc{X}$ is said to be:
    \begin{enumerate}
        \item birationally connected ($1$-connective) if $\tau_{\leq 0}^{bir}\esc{X}\simeq *$.
        \item birationally globally $(n-1)$-connected (or, $n$-connective) if it is connected and for all $0<i< n$ $\pi_i^{bir}(\esc{X},x)\simeq *$. 
    \end{enumerate}
\end{defn} 
\begin{defn}
    We say that a pointed birational local space is birationally $n$-Eilenberg-MacLane if it is birationally globally $n$-connective and birationally $n$-truncated.
\end{defn}
\begin{thm}\label{bir em}
Let $S$ be a QCQS scheme. Then, similar to that of an $\infty$-topos \textup{[\cite{lurie2009higher}, Lemma 7.2.2.11 (1)]}, then there are equivalences of categories: 
\begin{flalign*}
   Sh_{*,S}^{b}\simeq  \mathrm{Disc}(\mathcal{H}^{bir})_\bullet&\simeq \mathcal{EM}_0(\mathcal{H}^{bir}):\pi_0^{bir}\\
   Grp_S^{b}\simeq  Grp(\mathrm{Disc}(\mathcal{H}^{bir}))&\simeq \mathcal{EM}_1(\mathcal{H}^{bir}): \pi_1^{bir}\\
   Ab_S^{b}\simeq Ab(\mathrm{Disc}(\mathcal{H}^{bir}))&\simeq \mathcal{EM}_n(\mathcal{H}^{bir}):\pi_n^{bir}  \text{ for all }n\geq 2
\end{flalign*}
 
\end{thm}
\begin{proof}
Since $L_{bir}$ is path-injective (\Cref{path injectivity of Lbir}), this follows from \Cref{grp em} and \Cref{n-strict bir groups}.
\end{proof}
\begin{rem}\label[rem]{EM}
 Here is a rather straightforward argument. The induced presheaf Eilenberg-MacLane functors take a birational group object to birational Eilenberg-MacLane spaces, with inverse given by presheaf homotopy groups. Explicitly, there are equivalences 
    \begin{enumerate}
        \item $\mathrm{K}(-,0):Shv^b_S\simeq \mathcal{EM}_0(\mathcal{H}^{bir}(S)):\pi_0$.
        \item $\mathrm{K}(-,1):Grp^b_S\simeq \mathcal{EM}_1(\mathcal{H}^{bir}(S)): \pi_1$.
        \item $\mathrm{K}(-,i):Ab^b_S\simeq \mathcal{EM}_i(\mathcal{H}^{bir}(S)):\pi_i$ for all $i\geq 2$.
    \end{enumerate}
In short, we have equivalences of categories 
$\mathrm{K}(-,n):n\mbox{-}Grp^b_S\simeq \mathcal{EM}_n(\mathcal{H}^{bir}_S)$ for each $n\geq 0$. The identification of the theorem above with this statement follows from two facts (compare with the motivic case). First, $\pi_i$ of a birational local space is birational. Second, $\text{K}(A,n)$ is birational local iff $A$ is a birational presheaf. Indeed, these can be checked on the presheaf $\pi_i$'s by \Cref{checking bir locality on pi}.

     
Now the equivalence $$\text{K}(-,n):n\mbox{-}Grp_S\leftrightarrows \mathcal{EM}_n(\mathcal{P}(S)): \pi_n$$ restricts to the desired equivalence by restricting both sides to $\mathcal{H}^{bir}$ and noting that $$\mathcal{EM}_n(\mathcal{H}^{bir}(S))=\mathcal{EM}_n(\mathcal{P}(S))\bigcap \mathcal{H}^{bir}(S)$$ [\Cref{em of localization}] and that $\pi_i^{bir}={\pi_i}_{|\mathcal{H}^{bir}(S)}$ \Cref{birational truncation}.
\end{rem}

\begin{rem}
    An algebraic reason this is not surprising is that birational sheaves are flasque over arbitrary $qcqs$ base schemes. (To see this, let $U\xhookrightarrow{}X$ be an open immersion (not necessarily dense). Obviously, $(X\setminus\bar{U})\sqcup U\xhookrightarrow{}X$ is a dense open immersion. Thus, if $F$ is birational, $$F(X)\to F(X\setminus \bar{U}\sqcup U)\cong F(X\setminus \bar{U})\times F(U)$$ is an isomorphism. By composing with the (surjective) projection $ F(X\setminus \bar{U})\times F(U)\twoheadrightarrow  F(U)$, we see that the canonical map $F(X)\to F(U) $ is indeed surjective.) But flasque sheaves of abelian groups cannot have higher nisnevich cohomology groups. (Indeed, this follows from [Lemma 1.39, 1.40 \cite{riou2002theorie}] over noetherian bases and thus over all qcqs bases by noetherian approximation.) To follow the above theorem for $n\geq 2$, it suffices to verify the same on generators of $\mathcal{H}^{bir}_S$, namely $h^b(X)$ for smooth $X/S$. But,
\begin{flalign}
    \pi_iMap(h^b(X),\mathbf{K}(A-,n))&=\pi_i\mathcal{P}_{nis}(Sm_S)\big(X,\mathbf{K}(A-,n)\big )\\
&=\begin{cases}
A \text{ for }0=i\\
    H^{n-i}(X;A)=0 \text{ for }0<i\leq n\\
    \pi_{i-n}Map_{\mathcal{P}_{nis}(S)}\big(X,A\big )=0 \text{ for }i>n \text{ [as } A \text{ is discrete} ]
\end{cases}
\end{flalign}
The case $n=1$ is a little delicate, though it is also straightforward, since we are expecting $H^1_{nis}(-;G)$ to be birational and hence a nisnevich sheaf in particular. 
    \end{rem}
 \begin{rem}
This equivalence is an $\infty$-topos-theoretic property (\Cref{em equiv of a topos}) that fails in the standard motivic context (\Cref{mot not topos}). The theorem thus shows that a further birational localization forces the motivic homotopy category to exhibit properties of an $\infty$-topos.
 \end{rem}
 \begin{thm}\label{coh dim of bir}
Over a qcqs scheme, every object $\esc{E}\in\mathcal{H}^{bir}(S)$ has $0$-cohomological dimension. Hence, in particular, $\mathcal{H}^{bir}(S)$ has $0$-cohomological dimension.
\end{thm}
\begin{proof}
    By definition, for any $n>0$ and $X$ smooth over $S$,
$$\mathbf{H}^{n}_{\mathcal{H}^{bir}(S)}(A)(X/S):=\pi_0\Big(\mathrm{Map}_{\mathcal{H}^{bir}(S)}\big(L_{bir}X,\mathbf{K}(A-, n) \big)\Big)=\pi_0(\mathbf{K}(AX, n) \big)$$
vanishes.
    
    Since $\mathcal{H}^{bir}(S)$ is generated under colimits by $L_{bir}h_S(X)$ (for smooth $X/S$), we conclude that $\mathbf{H}^{n}(\mathcal{H}^{bir}(S)(\esc{E)}$ is $0$ for all $n>0$.  
\end{proof}
\begin{rem}
 These properties do not already hold in any other higher birational motivic categories (see [\cite{sfbat}] for the definition). In fact, the same example as in [\cite{choudhury2014connectivity}] shows that the $1$-birational motivic homotopy category is not an $\infty$-topos either. Indeed, his $\mathbb{A}^1$-local example $\mathbb{Z}^{nis}(\mathbb{G}_m)$ is not strongly $\mathbb{A}^1$-invariant. From his proof, or otherwise, it is clear that $\mathbb{Z}^{nis}(\mathbb{G}_m)$ is a summand of $\mathbb{Z}^\sqcup (\mathbb{G}_m)$, and is thus $1$-birational local since $\mathbb{G}_m$ is so. But $\mathcal{EM}_1(\mathcal{H}^{1}(k))$ consists of strongly $\mathbb{A}^1$-invariant, strongly $1$-birational-invariant sheaves of groups. In fact, this failure is not just to the $\mathbb{A}^1$-local aspect; it is of birational nature as well. There are $n$-birational, strongly $\mathbb{A}^1$-invariant abelian groups that are not strictly $n$-birational. One quick example is $K_{n+1}^M$.
 \end{rem}

\subsection{Birational groups/monoids}
We can generalize the results of the previous section about birational sheaves of groups to arbitrary birational monoid spaces. Recall that the birational homotopy category is presentable, and hence there is an internal bar construction. We shall denote this by $\mathrm{B}_{bir}: \mathrm{Mon}(\mathcal{H}^{bir})\to \mathcal{H}^{bir}$. As in \S2.3, one may define a notion of strongly $L_{bir}$-local group objects, but the following theorem shows, as in \Cref{n-strict bir groups}, that such a concept is unnecessary.
\begin{thm}\label{B_bir is B}
    There is an equivalence $\mathbf{B}_{bir}\simeq \mathbf{B}_{|Mon(\mathcal{H}^{bir}(S))}$.
\end{thm}
\begin{proof}
It follows from \Cref{bar of cart localization} that $\mathrm{B}_{bir}\simeq L_{bir} \mathrm{B}$, where $\mathrm{B}$ is the presheaf-level bar construction. Thus, it suffices to know that $\mathrm{B}$ preserves birational local objects. This is [\cite{0bat}, Lemma 3.2.4 (3)].
\end{proof}
We caution that $\mathbf{B}_{mot}\not\simeq \mathbf{B}_{|Mon(\mathcal{H}^{mot}(k))}$.
\begin{thm}\label{bir commute group competition}
        $L_{bir}$ commutes with the (Nisnevich) group completion of commutative monoids in $\mathcal{P}(S)$ (as well as in $\mathcal{P}_{nis}(S)$).
\end{thm}
\begin{proof}
 Since $\esc{M}^+\simeq \Omega\mathbf{B} \esc{M}$ and $L_{bir} $ commutes with $\Omega $ on connected objects and with $\mathbf{B}$ on all objects, the result follows immediately.
\end{proof}

Recall that the following is a property of an $\infty$-topos:
\begin{prop}\label[prop]{bir omega bar}
    The functor $$\Omega: \mathcal{H}^{bir}(S)_{\geq 1} \to  Grp(\mathcal{H}^{bir}_\bullet(S))\simeq Grp(\mathcal{H}^{bir}(S))$$ is an equivalence of infinity categories, with inverse given by the bar construction $\mathbf{B}_{|bir}$.
\end{prop}
\begin{proof}
Since $\mathbf{B}_{bir}\simeq \mathbf{B}_{|Mon(\mathcal{H}^{bir}(S))}$, the co-unit $$\mathbf{B}_{bir}\Omega \to 1_{\mathcal{H}^{bir}(S)_{\geq 1} }$$ is an equivalence, as it is equivalent to the restriction of the equivalence $ B\Omega\to 1_{\mathcal{P}(S)_{\geq 1}}$. 

Because grouplike objects are group-complete, this is just a restatement of [\cite{bas}, Theorem 3.3.7].     
\end{proof}
Note that this is a property of an $\infty$-topos that fails to hold in the ordinary motivic homotopy category \Cref{strong motivic groups and motivic connected spaces}. 
We conclude this section with one more topos-theoretic property of the birational motivic homotopy category: 
\begin{prop}\label[prop]{sub obj in bir}
    Suppose S is a geometrically unibranch scheme with finitely many generic points. Then the birational homotopy category of $S$ has a subobject classifier.
\end{prop}
\begin{proof}
    It suffices to show that the subobject classifier of $\mathcal{P}_{\Sigma}(Sm_S)$ is birational local. Indeed, since a localization is closed under limits, the subobject functor is simply the restriction of the subobject functor from the parent category to the localized one. 
    
    In general, the subobject classifier of $\mathcal{P}_{\Sigma}(Sm_S)=\mathcal{P}_{\sqcup}(Sm_S)$ is given by the sifted sheaf:
    $$X/S \mapsto \Omega^{\sqcup}_S(X):= \{ 
    W/S\big | W \text{  is a clopen subset of } X\}$$
    Since $X$ has a finite decomposition by its irreducible components (which are thus both open and closed) (see the lemma below), $$\Omega^{\sqcup}_S(X):= \{ 
    W/S\big | W \text{  is a component of } X\}.$$
   It is obvious that $\Omega^\sqcup _S$ is birational local.
\end{proof}
\begin{lem}\label[lem]{unibranch lemma}
   Suppose S is a Qcqs scheme. And $f:X\to S$ a smooth (qc and separated) morphism.
   \begin{enumerate}
       \item If $S$ is geometrically unibranch then the irreducible components of $X$ are disjoint.
       \item  If $S$ has finitely many generic points, then so does $X$.
   \end{enumerate}
\end{lem}
\begin{proof}
(1). We first show that $X$ has the same set of properties. Since being geometrically unibranch is a local condition, we reduce to the following algebra problem: if $A\to B$ is an essentially smooth local homomorphism of local rings such that $A$ is geometrically unibranch, then so is $B$. This follows from [\cite[\href{https://stacks.math.columbia.edu/tag/0DQ1}{Lemma 0DQ1 (1)}]{stacks-project}] and [\cite[\href{https://stacks.math.columbia.edu/tag/0C37}{Lemma 0C37 (3)}]{stacks-project}]. 

Therefore, it suffices to show that the irreducible components of a geometrically unibranch scheme are disjoint. This is, in fact, true for all unibranch schemes: if $x$ is a point that belongs to two irreducible components, then $\mathcal{O}_{X,x}$ will have at least two minimal primes.

(2). This follows from the fact that $f$ maps generic points to generic points and smooth schemes over fields have this property.
\end{proof}
\begin{rem}
    The motivic homotopy category most likely does not possess a subobject classifier. In fact, the subobject classifier $\Omega^{nis} _S$ of $Shv_S^{nis}$ cannot be $\mathbb{A}^1$-local, since there are nisnevich covers of $X$ does not extend (along the $0$ section) uniquely to a nisnevich cover of $\mathbb{A}^1\times X$.
\end{rem}

\begin{thm}
    The adjunction $\mathcal{H}^{bir}(S)\leftrightarrows \mathcal{P}(S)$ is monadic. In other words, the idempotent monads induced by the adjunctions identify $\mathcal{H}^{bir}$ as their Eilenberg-Moore category of algebras. In detail, this means that the forgetful map $\mathcal{H}^{bir}\hookrightarrow \text{LMod}_{L_{bir}}(\mathcal{P}
    {(S)})$ is an equivalence of categories over $\mathcal{P}(S)$.
\end{thm}
\begin{proof}
    In light of the theorem above, this is a consequence of Beck's Monadicity Theorem [\cite{lurie2017higher}, Theorem 4.7.0.3]. In fact, this holds for any localization closed under geometric realizations. Indeed, the conservative condition is redundant for fully faithful functors. 
\end{proof}

\begin{rem}
  It turns out that any Bousfield localization is monadic. Consequently, $$\mathcal{H}^{bir}(S)\subset\mathcal{P}_{\Sigma}(S),\mathcal{H}^{bir}(S)\leftrightarrows\mathcal{P}_{nis}(S), \mathcal{H}^{bir}(S)\leftrightarrows\mathcal{H}^{mot}(S),\text{ and }\mathcal{H}^{mot}(S)\leftrightarrows\mathcal{P}_{nis}(S) $$ are all monadic. However, localization via descent is not monadic in general. The only monadic localization here that is obtained by imposing descent is $\mathcal{P}_{\Sigma}(S)\subset \mathcal{P}(S)$. Even the localization $\mathcal{P}_{nis}(S)\subset \mathcal{P}(S)$ is not monadic. To see this, observe that in a monadic localization, an object $X$ is local if and only if the localization map $X\to LX$ is a split monomorphism. An example of such an object for the localization $\mathcal{P}_{nis}(k)\subset \mathcal{P}(k)$ is the separated presheaf generated by $\mathbb{Z}$ (or any other constant set, for that matter):
    $$U\mapsto \underline{\mathbb{Z}}(U)=\begin{cases}\mathbb{Z} \text{  if  } U\neq \emptyset \\
    0 \text{ otherwise}        
    \end{cases}$$
    Evidently, the sheafification $\underline{\mathbb{Z}}\to \mathbb{Z}^{nis}$ is a nontrivial split-injective map with retract given by $(-)^{0}\to *$. The same example shows that this issue persists for the motivic adjunction $\mathcal{H}^{mot}(S)\subset \mathcal{P}(S)$.  
\end{rem}


\subsection{Birational coverings}

\textbf{Birational Geometric coverings}
\begin{defn}
A collection $\{\esc{U}_i\to \esc{X}\}$ of morphisms of presheaves of spaces is said to be a birational geometric covering if $L_{bir}\check{C} p\to L_{bir}\esc{X}$ is a colimit diagram in $\mathcal{H}^{bir}(k)$, where $p: \bigsqcup\esc{U}_i\to \esc{X}$.
\end{defn}
\begin{lem}
    An $\mathbb{A}^1$-geometric covering is a birational geometric covering.
\end{lem}
\begin{proof}
    Since $L_{bir}\simeq L_{bir}L_{mot}$, this is obvious.
\end{proof}
So,
$$\text{Nis-geom-cover}\implies \mathbb{A}^1\text{-geom-cover}\implies \text{bir-geom-cover}$$
\begin{cor}
An $S$-space $\esc{X}$ is birationally connected if and only if there exists a birational (which is if and only if a Nisnevich if and only if an $\mathbb{A}^1$-) geometric covering $\esc{U}_i\to \esc{X}$ such that each $\esc{U}_{i}$ is birationally connected and, for every pair $i,j$, the sheaf $\pi_0^b(U_i\times_XU_j)$ has an $S$-rational point.
\end{cor}
\begin{proof}
    If $\esc{X}$ is birationally connected, then the identity map $\esc{X}\to\esc{X}$ is a Nisnevich cover with the given property. Using the above chain of implications, we reduce to the case of a birational geometric cover. But this follows from \Cref{geom L connectivity} in the presence of \Cref{path injectivity of Lbir}.
\end{proof}
\begin{cor}
   A scheme $X$ is birationally-connected iff there is a $U_i\to X$ be a Nisnevich covering of a scheme $X/S$ and such that for every pair $i,j$ the $\pi_0^b(U_i\times_XU_j)$ has an $S$-rational point. If $U_i$ is birationally-connected for all $i$, then 
\end{cor}
\begin{rem}
    Compare with the ordinary motivic counterparts \Cref{a1 geom conn}, \Cref{local-global-connectivity}. The condition there is slightly stronger due to the effectiveness of $L_{mot}$.
\end{rem}
\begin{cor}[Birational Motivic Van Kampen]
Let $\esc{U}_i\to \esc{X}$ be a birational-geometric covering (for example, a Nisnevich local covering or a motivic geometric covering) such that $\esc{U}_{ijk}$ is birationally connected for every $i,j,k$ (1-, 2-, or 3-fold). Then $\esc{X}$ is birationally connected, and the birational motivic homotopy group of $\esc{X}$ is given by the following coequalizer in $Grp^b_S$:
$$\underset{i,j}{*}\pi_1^{b}({\esc{U}_i\times_\esc{X}\esc{U}_j})\rightrightarrows \underset{i}{*}\pi_1^{b}({\esc{U}_i})\to \pi_1^{b}(\esc{X})\to *$$
\end{cor}
\begin{proof}
    This follows immediately from \Cref{L van campen} and \Cref{path injectivity of Lbir}.
\end{proof}


\begin{thm}[Birational version of Morel's Van Kampen, [\cite{morel2012a1}, Theorem 7.12\text{]}]\label{bvk}
Assume $S$ is qcqs. Let $\mathcal{U}:=\{U_i\to X\}$ be a basic Nisnevich cover in $Sm_S$ such that the 1-, 2-, and 3-fold fiber products are birationally connected and compatibly pointed across all fiber products of all finite pairs. Then there is a coequalizer diagram in $Grp_S^b$:
$$\underset{i,j}{*}\pi_1^b({U_i\times_XU_j})\rightrightarrows \underset{i}{*}\pi_1^b({U_i})\to \pi_1^b(X)\to *$$
\end{thm}
\begin{rem}
   When all the connectivity conditions are replaced by $\mathbb{A}^1$-connectivity, we can derive our birational Van Kampen theorem \Cref{bvk} directly from the $\mathbb{A}^1$ motivic Van Kampen theorem. Indeed, if all the relevant schemes are $\mathbb{A}^1$-connected, then [\cite{0bat}, Proposition 3.6.7] tells us that the sequence involving $\pi_1^b$ is simply the application of $\pi_0^b$ to Morel's Van Kampen sequence. The result then follows from the fact that, being a localization, $\pi_0^b: Grp^{\mathbb{A}^1}\to Grp^b$ preserves colimits.
\end{rem}

\textbf{Birational Spatial coverings}

\begin{defn}[\Cref{defn for covering}, \Cref{definition for L covering}]
\begin{enumerate}
\item  A morphism $f:\esc{Y}\to \esc{X}$ of birational motivic spaces is a birational covering if $f\in \big({\mathcal{H}^{bir}_{/\esc{X}}}\big)_{\leq 0}$. Denote the category of such coverings by $Cov_{\mathcal{H}^{bir}}(\esc{X})$. 
\end{enumerate}
When $f:\esc{Y}\to \esc{X}$ is a morphism of $S$-spaces, then:
\begin{enumerate}[resume] 
    \item We say that it is a birational fibration if it is a local object of ${\mathcal{P}}_{/\esc{X}}$ with respect to birational motivic equivalences contained in ${\mathcal{P}}_{/\esc{X}}$. We denote the category of birational fibrations of $\esc{X}$ by $\mathcal{F}ib_{\esc{X}}^{b}.$
\item We call it an $\mathbb{A}^1$ covering if it is a birational fibration and is a $0$-truncated object in the $\infty$-category of birational fibrations of $\esc{X}$. Denote the category of birational coverings of $\esc{X}$ by $Cov_{b}(\esc{X})$.
\end{enumerate}
\end{defn}
\begin{prop}
A morphism $\esc{X}\to \esc{Y}$ of birational motivic spaces is a birational covering if and only if it is a (Nisnevich) covering if and only if it is an $\mathbb{A}^1$-covering.
\end{prop}
\begin{proof}
Follows immediately from \Cref{Covering of Localization}.
\end{proof}

\begin{prop}\label[prop]{birational torsor}
    Let $G$ be a birational sheaf of groups. Then any Nisnevich $G$-torsor $\esc{Y}\to \esc{X}$ is a birational cover. 
\end{prop}
\begin{proof}
Since $ Grp_{L{bir}}(\mathcal{H}^{bir})\simeq  Grp(\mathcal{H}^{bir})$ (\Cref{B_bir is B}), this follows from \Cref{torsors are covering}.
\end{proof}
\begin{rem}
Using the model categorical definition of Morel, one may use the original idea from [\cite{morel2012a1}, 7.5 (1)] to prove the above theorem. Let us repeat it for the convenience of the reader. Suppose we are given a commutative square on the left, where $\esc{A}\to \esc{B}$ is a birational weak equivalence \[\xymatrix{
 \esc{A}\ar[r]\ar[d]& \esc{Y}\ar[d]&& \esc{A}\ar[rd]\ar[r]&\esc{Y}\times_\esc{X}\esc{B}\ar[r]\ar@<-.2em>[d]_{a}& \esc{Y}\ar[d]\\
\esc{B}\ar[r]& \esc{X}&& &\esc{B}\ar@<-.2em>@{-->}[u]\ar[r]& \esc{X}
 }\]
 Consider the diagram on the right, obtained by pulling the torsor to $\esc{B}$. It follows that the projection map $\esc{Y}\times_\esc{X}\esc{B}\xrightarrow{\phi}\esc{B}$ is a $G$ torsor. The commutativity of the starting diagram yields that this torsor is trivial on $\esc{A}$. Since $\mathrm{B}\esc{G}$ is birational local by \Cref{B_bir is B}, and $\esc{A}\to \esc{B}$ is given to be a birational weak equivalence, it follows that the map $[\esc{B}, \mathrm{B}G]\to [\esc{A}, \mathrm{B}G]$ is an isomorphism. So the triviality of the torsor on $\esc{A}$ implies that it trivializes on $\esc{B}$ as well. Let $s: \esc{B}\to \esc{Y}\times_\esc{X}\esc{B}$ be the corresponding trivialization. Then the composite  $$\esc{B}\xrightarrow{s} \esc{Y}\times_\esc{X}\esc{B}\xrightarrow[]{a} \esc{Y}$$ makes the bottom triangle commute. Obviously, the upper triangle may not commute. However, since the torsor $\esc{Y}\to \esc{X}$ is a $\esc{G} $-torsor trivializing on $\esc{A}$, any two such trivializations are $\esc{G}$ automorphic. Therefore, there is a section $g: \esc{A}\to\esc{G}$ such that $g(a\circ s)$ works. 
\end{rem}
\begin{thm}\label{initial object of bir cov}
Let $\esc{X}$ be a pointed birational motivic space (not necessarily connected). 
\begin{enumerate}
    \item Then the category $\mathcal{C}ov_{\mathcal{H}^{bir}}(\esc{X})_*$ always has an initial object, say denoted by ${\widetilde{{\esc{X}}}}^{^{bir}}$.
    \item This initial object is the unique one with a birationally simply connected source (i.e., $\tau_{\leq 1}^{bir}{\widetilde{{\esc{X}}}}^{^{bir}}\simeq *$). Each of the inclusions  $$\mathcal{C}ov_{\mathcal{H}^{bir}}(\esc{X})_*\subset \mathcal{C}ov_{\mathcal{P}_{nis}}(\esc{X})_*\subset \mathcal{C}ov_{\mathcal{P}_{}}(\esc{X})_* $$ preserves this initial object. In other words, this initial object is the unique (Nisnevich) covering with a (Nisnevich) simply connected source.
    \item The canonical map $$ \pi_1^{bir}(\esc{X})\to \underline{Aut}_{\esc{X}}(\widetilde{{\esc{X}}}^{^{bir}})$$ is an equivalence.
\end{enumerate}
\end{thm}
\begin{proof}
  (The first one is mentioned at [\cite{0bat}, Corollary 3.6.2], whose proof is similar to \Cref{Covering of Localization}. We shall, however, prove this using the categorical techniques developed in \S2 just to keep up with the style of the paper.)
     
 From \cref{birational truncation}, it follows that $\Pi^{b}\esc{X}$ is a birational motivic space. Thus, \Cref{Universal cover for fundamental localization} applies, and statement (1) follows. 
    
Statement (2) is a formal consequence of the fact that, from the referred proposition, one has $(\widetilde{{\esc{X}}}^{^{bir}})\simeq \tau_{>1}^{}\esc{X}$, where this last term is the unique (nisnevich) simply connected (simplicial) cover of $\esc{X}$. Since this is already birational local, we have $$\tau_{\leq 1}^{bir}{\widetilde{{\esc{X}}}}^{^{bir}}\simeq \tau_{\leq 1}^{}{\widetilde{{\esc{X}}}}^{^{bir}}\simeq \tau_{\leq 1}\tau_{>1}^{}\esc{X}\simeq *.$$
    
The third statement follows from \Cref{aut=pi_1} because $L_{bir}$ is fundamental on every object.
\end{proof}

\begin{rem}
    Compare the above theorem with its ordinary motivic counterpart, namely \Cref{motivically simply connected covering space}. There, a similar statement was proven for internally connected motivic spaces. Here, in the birational case, however, we did not require the base to be connected. Thus, we again see that the birational motivic homotopy category, $\mathcal{H}^{bir}$, behaves like an $\infty$-topos (see \Cref{initial object of cov in a topos}). Though in the motivic case, we did not really prove that an initial internal motivic covering may not exist when the base is not internally (motivically) connected. It is interesting in its own right to know whether this has a positive answer.
    \end{rem}

Since $L_{bir}$ is neither locally cartesian (see [\cite{0bat}, Counterexample 3.2.13]) nor effective (see [\cite{0bat}, Counterexample 3.5.7]), proving the following by pulling back the (birational) universal cover of $L_{bir}\esc{E}$ to $\esc{E}$ (similar to Morel [\cite{morel2012a1}, $\mathbb{A}^1$-coverings]) seems difficult. 
 \begin{cor}[compare with the $\mathbb{A}^1$ analogue, \Cref{Lmot covering is covering of Lmot}, no connectivity condition is required there]\label{Lbir covering is covering of Lbir}
Let $\esc{X}$ be an $\mathbb{A}^1$-connected (equivalently, motivically connected when the base is a field, \Cref{motivic 0 truncation is injective}) space. Then we have an equivalence of categories $$\eta^*_{bir}:Cov_{b}({L_{bir}}\esc{X)}=Cov_{\mathcal{H}^{bir}}({L_{bir}}\esc{X)}\leftrightarrows  Cov_{{b}}(\esc{X)}: L_{bir}$$ where $\eta_{bir}$ is the localization unit $1\to L_{bir}$.
\end{cor}
\begin{proof}
Recall that $L_{bir}$ is locally cartesian on motivically connected objects [\cite{0bat}, Corollary 3.2.16]. Moreover, since $\esc{X}$ is an $\mathbb{A}^1$-connected space, [\cite{0bat}, Corollary 3.5.2] implies that $L_{bir}\esc{X}$ is connected. Therefore, the localization map $\esc{X}\to L_{bir}\esc{X}$ is an effective epimorphism in $\mathcal{P}_{nis}(S)$ [\cite{lurie2009higher}, Proposition 7.2.1.14]. With these observations, the claim follows from \Cref{L covering is covering of L}.
\end{proof}


 \begin{thm}\label{universal birational covering}
 Let $\esc{X}\in \mathcal{P}(S)$ be an $\mathbb{A}^1$ connected space. Then 
\begin{enumerate}
    \item   $\esc{X}$ has a birationally simply connected birational cover $\widetilde{\esc{X}}^b$.
    \item Choice of a base point makes this birationally simply connected cover into the initial object of $Cov_b\esc{X}_*$.
    \item The canonical map $\pi_1^b\esc{X}\to Aut_{\esc{X}_*}(\widetilde{\esc{X}}^b)$ is an isomorphism.

\end{enumerate}    
 \end{thm}
\begin{proof}
Since $\esc{X}$ is $\mathbb{A}^1$-connected, so is $L_{bir}\esc{X}$ by [\cite{0bat}, Corollary 3.5.2]. Therefore, by the above corollary, replacing $\esc{X}$ with $L_{bir}\esc{X}$ we may assume that $\esc{X}$ is birational local. By [\cite{morel2012a1}, Corollary 6.2], it follows that $\Pi^{b}\esc{X}$ is a birational motivic space. So \Cref{Universal cover for fundamental localization} applies, and statement (1) follows. 

Statement (2) is a formal consequence of the fact that, from the referred proposition, one has $\tilde{\esc{X}}\simeq \tau_{>1}\esc{X}$, which is the unique simply connected (simplicial) cover of $\esc{X}$. 

The third statement follows from \Cref{aut=pi_1}.
\end{proof}

\begin{rem}
     The main hurdle when $\esc{X}$ is not motivically connected, but only birationally connected, is that $$\esc{Z}\simeq \esc{X}\times _{L_{bir}\esc{X}}\esc{Y}\to \esc{Y}$$ need not be a birational weak equivalence, since our category is not locally cartesian. Thus, the above theorem is weaker than its $\mathbb{A}^1$-counterpart [\cite{morel2012a1}, Theorem 7.8].
\end{rem}

\begin{rem}
When the base space is actually Nisnevich connected, the result is a little stronger. In that case, the simply connected cover maps to the birational simply connected cover under $L_{bir}$ (see [\cite{0bat}, Corollary 3.6.4]).\end{rem}
\begin{prop}
Suppose $\esc{Y}\in \mathcal{P}(S)$ and $y:*\to \esc{Y}$ is a base point. Then pullback along $y$ induces a functor $$y^*:Cov_{b}(\esc{{Y}})\to Cov_{b}(*)=Sh^{b}_k.$$ This functor factors through the category of birational sheaves (of sets) with an action by $\pi_1^{b}(\esc{Y},y)$ (via the canonical functor forgetting the action):  $$\Gamma_y^{b}:Cov_{b}(\esc{{Y}})\longrightarrow \pi_1^{b}(Y,y)\mbox{-}Sh^{b}_k$$ which is an equivalence of categories when $\esc{Y}$ is birationally connected.
\end{prop}
\begin{proof}
The existence of the pullback functor follows from the pullback stability of (birational) coverings, \Cref{covering lemma} (3). That it factors through $\pi_1^{b}(\esc{Y},y)\mbox{-}Sh^{b}_k$ is a consequence of \Cref{action of the fundamerntal groupoid}, in view of the fact that $$\Pi^{b}\esc{Y}\simeq \Pi^{bir}L_{bir}\esc{Y}$$ (in other words, $\Pi^{b}\esc{Y}$ is birational local) (see \Cref{birational truncation}).

Since $\mathcal{H}^{b}$ is closed under extensions over connected objects [\cite{0bat}, Lemma 3.3.1], the claim follows immediately from \Cref{L galois theory}.
\end{proof}

\subsection{On the birational motivic homotopy category being an \texorpdfstring{$\infty$}{}-topos}

At the beginning of \S4, we explained how the transition from ordinary motivic homotopy to birational motivic homotopy suggests a near `classical topological' behavior, allowing us to brute-force various topos-theoretic properties. In fact, through \S4.1, \S4.2, and \S4.3, we have proven many of those topos-theoretic properties. However, the birational localization functor $L_{bir}: \mathcal{P}_{nis}(S)\to \mathcal{P}_{nis}(S)$ is not even locally cartesian [\cite{0bat}, Counterexample 3.2.13], let alone left exact. At this juncture, we are thus unable to answer the following question over an arbitrary scheme: 
\begin{Q}\label[Q]{bir topos question}
    For a Qcqs scheme $S$, is the presentable $\infty$-category $$\mathcal{H}^{bir}(S):=L_{bir}\mathcal{H}^{mot}(S)\simeq L_{dense} \mathcal{P}_{\Sigma}(S)$$ (where the last equivalence is due to \textup{[\cite{0bat}, Corollary 2.2.12(2)]}) an $\infty$-topos?
\end{Q}

In fact, we have only been able to show the existence of a subobject classifier for geometrically unibranch Qcqs schemes with finitely many irreducible components; see \Cref{sub obj in bir}. In the result below, we will show that, under this condition on the base scheme, \Cref{bir topos question} actually has a positive answer:
\begin{thm}\label{Hbir a topos}
    Suppose $S$ is a Qcqs, geometrically unibranch scheme with finitely many generic points (for example, the spectrum of a field). Then $\mathcal{H}^{bir}(S)$ is an $\infty$-topos.
\end{thm}
\begin{proof}
    The proof follows from the abstract lemma below (\Cref{coreflexive sub topos}) and the fact that, under the conditions of the theorem, $\mathcal{H}^{bir}$ is closed under colimits in $\mathcal{P}_{nis}(S)$ (as well as in $\mathcal{P}_\Sigma(S)$; see \Cref{closed under coproducts} below).
\end{proof}
\begin{lem}\label[lem]{coreflexive sub topos}
    Suppose $\mathcal{C}$ is an $\infty$-topos and $L\mathcal{C} \subset \mathcal{C}$ is an accessible localization which is closed under all colimits. Then $L\mathcal{C}$ is an $\infty$-topos.
\end{lem}
\begin{proof}
    Since $L\mathcal{C}$ is a presentable $\infty$-category, it suffices to check conditions (3) (ii), (iii), and (iv) from [\cite{lurie2009higher}, Theorem 6.1.0.6 (3)]. Each of these conditions concerns interactions between certain limits and colimits. Since these hold in $\mathcal{C}$ and $L\mathcal{C}$ is closed under both limits and colimits in $\mathcal{C}$, the result follows.
\end{proof}
\begin{prop}\label[prop]{closed under coproducts}
    Let $S$ be a Qcqs, geometrically unibranch scheme with finitely many generic points. Then $L_{bir}: \mathcal{P}_\Sigma(S)\to \mathcal{P}_{\Sigma}(S)$ and $L_{bir}: \mathcal{P}_{nis}(S)\to \mathcal{P}_{nis}(S)$ preserve all colimits. In other words, both $\mathcal{H}^b(S)\subset \mathcal{P}_\Sigma(S)$ $\mathcal{H}^b(S)\subset \mathcal{P}_{nis}(S)$ are closed under all colimits.
\end{prop}
\begin{proof}
    Using  [\cite{0bat}, Corollary 3.1.6], it suffices to show that $L_{bir}$ preserves coproducts. Suppose $j:U\hookrightarrow X$ is a dense open immersion in $Sm_S$. Since $X$ is a finite disjoint union decomposition of its irreducible connected components (\Cref{unibranch lemma}), we can assume that $j$ is a finite disjoint union of dense open immersions of connected schemes. Since connected schemes are points in the $\sqcup$-topology, it follows that coproducts (computed in $\mathcal{P}_\Sigma(S)$) of birational local objects are birational local.

    Since $\mathcal{P}_{nis}(S)$ is closed under coproducts in $\mathcal{P}_\Sigma(S)$, the claims about the Nisnevich topology follow from the one proven above.
\end{proof}

\begin{cor}
    Under the above condition on $S$, there are essential geometric morphisms of $\infty$-toposes $$ \mathcal{P}_{nis}(S), \mathcal{P}_{\Sigma}(S)\to \mathcal{H}^{bir}(S).$$
\end{cor}
\begin{proof}
Owing to \Cref{closed under coproducts}, there are right adjoints to the right adjoints $\mathcal{H}^{bir}(S)\subset \mathcal{P}_\Sigma(S)$ and $\mathcal{H}^{bir}(S)\subset \mathcal{P}_{nis}(S)$ by the adjoint functor theorem.
\end{proof}
Actually, using the main result of [\cite{sfbat}], we can easily get rid of the geometrically unibranch condition:

\begin{thm}\label{bir topos on fin gen pt}
    Let $S$ be a Qcqs scheme with finitely many generic points. Then $\mathcal{H}^{bir}(S)$ is an $\infty$-topos.
\end{thm}
\begin{proof}
    By [\cite{sfbat}, Theorem 4.3.4] we know that there is an equivalence of $\infty$-categories:
    $$\mathcal{H}^{bir}(S)\to\underset{\eta\in S^{(0)}}{\prod}\mathcal{H}^{bir}(k(\eta))$$
    where the product on the right-hand side is computed in $Cat_\infty$. By \Cref{Hbir a topos}, we know that each category $\mathcal{H}^{bir}(k(\eta))$ in the product above is an $\infty$-topos. So it suffices to know that the Cartesian product of $\infty$-toposes is an $\infty$-topos. This is [\cite{lurie2009higher}, proposition 6.3.2.1].
\end{proof}
To compare with \Cref{robalo}, we summarize the above result in the following manner:
\begin{cor}\label[cor]{bir yoneda}
  Let $S$ be a Qcqs scheme with finitely many generic points. Then $$h^b_S(-): Sm_S\to \mathcal{H}^{bir}(S)$$ contracts the affine line, sends the Nisnevich and cdh squares to pushout squares, and has an $\infty$-topos as its target that is Postnikov complete and of cohomological dimension $0$.
\end{cor}
\begin{proof}
    The $\mathbb{A}^1$ and Nisnevich parts follow from the definition: $\mathcal{H}^{bir}(S):=L_{bir}L_{\mathbb{A}^1}L_{nis}\mathcal{P}(S)$. The cdh descent case follows from [\cite{0bat}, Remark 2.1.5]. That the target is an $\infty$-topos is the theorem above. Its Postnikov completeness is proven in \Cref{post comp Hbir}, and $0$-cohomological dimension in \Cref{coh dim of bir}.
\end{proof}

\textbf{Conclusion and Open question}

We conclude this paper with a set of open questions that the author plans to return to later:

(1). Although $\mathcal{H}^{bir}(S)$ exhibits several topos-theoretic properties over an arbitrary base scheme $S$ (as shown in \S4.1–4.3), its status as an $\infty$-topos in full generality (for schemes with infinitely many generic points) remains open. 


(2). It also remains an open question whether the Birational Yoneda functor of \Cref{bir yoneda} has a universal property in $\mathcal{T}op_\infty$ (similar to the one in \Cref{robalo}).


(3). Finally, since in $pro\mbox{-}\mathcal{H}^{bir}$ every scheme is identified with its discrete set of generic points (which has Nisnevich homotopy dimension $0$), it is reasonable to hope that the birational $\infty$-topos is locally of homotopy dimension $0$.